
\documentclass[11pt]{article}

\usepackage[margin=1in]{geometry}
\usepackage{amsmath,amssymb,amsthm,mathtools}
\usepackage{enumitem}
\usepackage{booktabs}
\usepackage{aliascnt}
\usepackage[hidelinks]{hyperref}
\usepackage[nameinlink,noabbrev]{cleveref}
\usepackage{microtype}
\usepackage{comment}

\numberwithin{equation}{section}

\setlength{\parindent}{0pt}
\setlength{\parskip}{6pt}

\newtheorem{theorem}{Theorem}[section]
\newaliascnt{proposition}{theorem}
\newtheorem{proposition}[proposition]{Proposition}
\aliascntresetthe{proposition}
\newaliascnt{lemma}{theorem}
\newtheorem{lemma}[lemma]{Lemma}
\aliascntresetthe{lemma}
\newaliascnt{corollary}{theorem}
\newtheorem{corollary}[corollary]{Corollary}
\aliascntresetthe{corollary}
\newaliascnt{definition}{theorem}
\newtheorem{definition}[definition]{Definition}
\aliascntresetthe{definition}
\newaliascnt{remark}{theorem}
\newtheorem{remark}[remark]{Remark}
\aliascntresetthe{remark}

\crefname{theorem}{Theorem}{Theorems}
\Crefname{theorem}{Theorem}{Theorems}
\crefname{proposition}{Proposition}{Propositions}
\Crefname{proposition}{Proposition}{Propositions}
\crefname{lemma}{Lemma}{Lemmas}
\Crefname{lemma}{Lemma}{Lemmas}
\crefname{corollary}{Corollary}{Corollaries}
\Crefname{corollary}{Corollary}{Corollaries}
\crefname{definition}{Definition}{Definitions}
\Crefname{definition}{Definition}{Definitions}
\crefname{remark}{Remark}{Remarks}
\Crefname{remark}{Remark}{Remarks}

\newcommand{\TV}{\mathrm{TV}}
\newcommand{\Unif}{\mathrm{Unif}}
\newcommand{\PreCL}{\mathrm{PreCL}}
\newcommand{\CL}{\mathrm{CL}}
\newcommand{\cA}{\mathcal A}
\newcommand{\cI}{\mathcal I}
\newcommand{\cT}{\mathcal T}
\newcommand{\cU}{\mathcal U}
\newcommand{\E}{\mathbb E}
\newcommand{\Pp}{\mathbb P}
\newcommand{\Var}{\operatorname{Var}}
\newcommand{\Z}{\mathbb Z}
\newcommand{\1}{\mathbf 1}
\newcommand{\cloc}{c_{\mathrm{loc}}}
\newcommand{\Hh}{\mathsf H}
\newcommand{\Mloc}{M_{\mathrm{loc}}^*}
\newcommand{\dTV}{d_{\TV}}
\newcommand{\nuUnif}{\nu_{\Unif}}
\newcommand{\UB}{\mathrm{UB}}
\newcommand{\ileft}{\mathrm I_{\rm left}}
\newcommand{\iright}{\mathrm I_{\rm right}}
\newcommand{\stable}{\mathrm{stable}}

\title{Time-Inhomogeneous Riffle Shuffles: Upper and Matching Lower Bounds}
\author{Time-inhomogeneous riffle shuffle note}
\date{}

\begin{document}
\maketitle

\section{Purpose and status}

We prove buffered upper and lower total-variation bounds for time-inhomogeneous riffle shuffles with deterministic column counts. The two effective margins are the maximal-word margin and the cold-spot margin.

The upper bound is first proved directly from empirical column costs. An exposed initial pile leaves a reserve of labels for endpoint smoothing; exact conditional sampling then controls natural exploration histories and their exponential moments. The macroscopic upper bound follows by transferring only bounded scalar costs, under the ordered approximation hypotheses of Theorem~\ref{thm:deterministic-upper-bound}. The lower bound in Theorem~\ref{thm:buffered-reduced-lower-bound} uses the largest-word and cold-spot obstructions, with concentration, finite type construction, collision estimates and the ascent test proved directly in the fixed-column model.

\section{Common setup}
\label{sec:common-setup}

For every positive integer $m$, write $[m]=\{1,\ldots,m\}$. Fix $k\ge2$,
write $[k]_0=\{0,\ldots,k-1\}$, and set
\begin{equation}
\label{eq:probability-simplex}
  D_k=\left\{p=(p_0,\ldots,p_{k-1})\in[0,1]^k:
  \sum_{\ell=0}^{k-1}p_\ell=1\right\}.
\end{equation}
Fix $T<\infty$ and $K=\lfloor T\log N\rfloor$. Let
\[
  n^{(j)}=(n^{(j)}_0,\ldots,n^{(j)}_{k-1}),
  \qquad
  \sum_{\ell=0}^{k-1}n^{(j)}_\ell=N,
\]
be deterministic pile sizes, and write $p^{(j)}=n^{(j)}/N\in D_k$. The macroscopic profile is a family $(\mu_t)_{0\le t\le T}$ of probability measures on $D_k$; we abbreviate this family by $\mu_\bullet$.
For this deterministic pile sequence, let $P(n,K)$ denote the law on the
symmetric group $\mathfrak S_N$ of the associated inverse $K$-shuffle. Here the inverse $K$-shuffle is the inverse of the chronological composition of the forward shuffles; its one-step inverse operations are consequently applied in reverse chronological order. When
the pile sequence is fixed by context, we also write $P_N^K=P(n,K)$.
Let $\nuUnif$ be the uniform law on $\mathfrak S_N$, and write $d_{\TV}$ for
total-variation distance. For a fixed pile sequence, we use
\begin{equation}
\label{eq:mixing-time-definition}
  t_{\mathrm{mix}}^{(N)}(\varepsilon)
  =\inf\{m\ge0:d_{\TV}(P(n,m),\nuUnif)\le\varepsilon\}.
\end{equation}
Unless a subscript is displayed, \(\Pp\) and \(\E\) refer to the
shuffle-matrix randomness associated with the pile sequence under discussion.
For $H\subseteq\Z$, write
\[
  E(H)=\{(i,i+1):i,i+1\in H\},
  \qquad
  \partial H=\{(i,i+1):\mathbf 1_H(i)\ne\mathbf 1_H(i+1)\}.
\]
Thus $E(H)$ and $\partial H$ are respectively the internal edges and the
edge boundary of $H$ in the integer path. For an integer $x$ and $r\ge0$,
write $(x)_r=x(x-1)\cdots(x-r+1)$, with $(x)_0=1$. Finally,
$\log_N x=(\log x)/(\log N)$, and $X_N=O_\Pp(a_N)$ means that $X_N/a_N$ is bounded in probability under the
probability law currently under discussion.
For $p\in D_k$, define
\begin{equation}
\label{eq:one-column-costs}
\begin{aligned}
  p_{\max}(p)=\max_\ell p_\ell,
  &\qquad
  a_p=\log\frac1{p_{\max}(p)},
  &\qquad
  \psi_p(\theta)=-\log\sum_{\ell=0}^{k-1}p_\ell^\theta .
\end{aligned}
\end{equation}
For $\theta\ge1$, also define the tilted law
\[
  p^{[\theta]}_\ell=\frac{p_\ell^\theta}{\sum_a p_a^\theta},
\]
and the one-column cost and entropy
\[
  I_p(\theta)=\sum_\ell p^{[\theta]}_\ell\log\frac1{p_\ell},
  \qquad
  H_p(\theta)=H(p^{[\theta]})=\theta I_p(\theta)-\psi_p(\theta).
\]
Here \(H(q)=-\sum_\ell q_\ell\log q_\ell\) is Shannon entropy, with the
usual convention \(0\log0=0\).
All time integrals in this article are ordinary Lebesgue integrals of the
displayed scalar integrands, whose measurability and integrability are part
of the notation. This applies in particular to the cell probabilities in
Definition~\ref{def:almost-time-profile-like-cuts}. No convention assigning
zero to a nonintegrable function is intended. In the tilted expressions,
$0^\theta\log(1/0)$ is interpreted as its limiting value zero.
At macroscopic time $t$, set
\begin{equation}
\label{eq:profile-costs}
\begin{aligned}
  a(t)=\E_{\mu_t}a_p,
  &\qquad
  \psi_t(\theta)=\E_{\mu_t}\psi_p(\theta),
  &\qquad
  I_t(\theta)=\E_{\mu_t}I_p(\theta),
  &\qquad
  H_t(\theta)=\E_{\mu_t}H_p(\theta).
\end{aligned}
\end{equation}
For $0\le r\le s\le T$, write
\begin{equation}
\label{eq:integrated-profile-costs}
\begin{aligned}
  A(r,s)=\int_r^s a(t)\,dt,
  &\qquad
  B_\theta(r,s)=\int_r^s\psi_t(\theta)\,dt,
  &\qquad
  J_\theta(r,s)=\int_r^s I_t(\theta)\,dt.
\end{aligned}
\end{equation}

\section{Main upper-bound result with reduced margins}
\label{sec:reduced-upper-margins}

For the whole deck, set
\begin{equation}
\label{eq:global-beta}
  \beta(s)=1-A(0,s),
\end{equation}
and define the four global upper-bound margins by
\begin{equation}
\label{eq:global-lambda-zero}
  \Lambda_0(T)=1-A(0,T),
\end{equation}
\begin{equation}
\label{eq:global-lambda-one}
  \Lambda_1(T)
  =\sup_{0\le s\le u\le T}\inf_{\theta\in[3,4]}
  \left\{\frac{\theta+3}{2}\beta(s)-B_\theta(s,u)-2B_2(u,T)\right\}.
\end{equation}
\begin{equation}
\label{eq:global-lambda-two}
  \Lambda_2(T)
  =\sup_{0\le s\le u\le T}\inf_{\theta\in[3,4]}
  \left\{\frac{\theta-1}{2}\beta(s)-B_\theta(s,u)-2B_2(u,T)\right\},
\end{equation}
and
\begin{equation}
\label{eq:global-lambda-three}
  \Lambda_3(T)=\sup_{0\le s\le T}\{\beta(s)-B_2(s,T)\}.
\end{equation}
The full upper-bound margin and its critical time are
\begin{equation}
\label{eq:global-upper-margin}
  \Lambda_{\UB}(T)=
  \max\{\Lambda_0(T),\Lambda_1(T),\Lambda_2(T),\Lambda_3(T)\}.
\end{equation}
\begin{equation}
\label{eq:full-upper-threshold}
  T_{\UB}=\inf\{T>0:\Lambda_{\UB}(T)<0\}.
\end{equation}
Let
\begin{equation}
\label{eq:reduced-upper-threshold}
  T_{\UB}^{\mathrm{red}}
  =\inf\{T>0:\max(\Lambda_0(T),\Lambda_1(T))<0\}.
\end{equation}
We use $\inf\varnothing=+\infty$ for both critical-time thresholds.
Theorem~\ref{thm:dominance-reduction} proves that
\[
  \Lambda_2(T)\le\Lambda_0(T),
  \qquad
  \Lambda_3(T)\le\Lambda_0(T),
\]
and consequently that the effective margin is
\[
  \Lambda_{\UB}(T)=\max\{\Lambda_0(T),\Lambda_1(T)\}.
\]
It also gives
\[
  T_{\UB}^{\mathrm{red}}=T_{\UB},
\]
so the reduced and four-margin thresholds agree.



\begin{definition}[Almost time-profile-like cuts]
\label{def:almost-time-profile-like-cuts}
Fix a finite measurable partition $D_k=D_0\cup\cdots\cup D_m$. The deterministic cuts are $(\rho,\varphi)$-almost-$\mu_\bullet$-like on $[0,T]$ if, for every positive-length grid interval
\[
  I=[a,\min(a+\rho,T)]\subset[0,T],
  \qquad a\in\rho\Z,
\]
and every cell $D_i$,
\[
  \left|
  \frac1{\log N}\#\{j:j/\log N\in I,\ p^{(j)}\in D_i\}
  -\int_I\mu_t(D_i)\,dt
  \right|\le \varphi |I|+\frac1{\log N}.
\]
\end{definition}

\begin{definition}[Randomized time-profile-like cuts]
\label{def:randomized-time-profile-like-cuts}
Fix $T>0$, a finite partition $\mathcal P$, and positive
$\rho,\varphi$, all independent of $N$.  Let
$X=(X^{(j)})_{1\le j\le K_N}$ be a random array with
$X^{(j)}\in\mathbb Z_{\ge0}^k$ and
$\sum_\ell X^{(j)}_\ell=N$, where $K_N=\lfloor T\log N\rfloor$.
The array is randomized-$(\mathcal P,\rho,\varphi)$-almost-
$\mu_\bullet$-like if its realized cuts satisfy
Definition~\ref{def:almost-time-profile-like-cuts} with these fixed
parameters with probability $1-o(1)$.  The parameters in an upper-bound
application are chosen in the order required by
Theorem~\ref{thm:deterministic-upper-bound}.
Conditional on $X$, the shuffle matrix has independent uniform
columns with the prescribed counts.  The law $P(X,K_N)$ is the
resulting mixture of the deterministic inverse-shuffle laws.
No independence between the random count vectors $X^{(j)}$ is assumed.
\end{definition}

\begin{definition}[$\chi$-regular partitions]
\label{def:chi-regular-partitions}
A finite Borel partition $D_k=D_0\cup\cdots\cup D_m$, with representatives $p^{(i)}\in D_i$, is $\chi$-regular, with a constant $C=C_k$ independent of $\chi$, if
\[
  \sup_{p,p'\in D_i}\sup_{\theta\in[2,4]}
  |\psi_p(\theta)-\psi_{p'}(\theta)|\le C\chi,
  \qquad
  \sup_{p,p'\in D_i}
  \left|\log\frac1{p_{\max}(p)}-\log\frac1{p_{\max}(p')}\right|
  \le C\chi,
\]
and if each coordinate satisfies
\[
  p_\ell^{(i)}>2\chi,\ p\in D_i
  \quad\Longrightarrow\quad
  p_\ell>\chi,
  \qquad
  \left|\log\frac1{p_\ell}-\log\frac1{p^{(i)}_\ell}\right|\le C\chi,
\]
while
\[
  p_\ell^{(i)}\le2\chi,
  \ p\in D_i
  \quad\Longrightarrow\quad
  p_\ell\le3\chi .
\]
Such partitions with a common constant $C_k$ are obtained by intersecting
the simplex with a finite rectangular grid of coordinate diameter at most
$c_k\chi^2$ and choosing a representative in every nonempty cell.
For a retained coordinate, the logarithm has derivative at most $1/\chi$;
for an omitted coordinate, the diameter bound gives $p_\ell\le3\chi$.
The map $p\mapsto p_{\max}$ is Lipschitz and $p_{\max}\ge1/k$, so $a_p$
has a uniform Lipschitz bound. The functions $\psi_p(\theta)$ have uniformly
bounded first derivatives for $2\le\theta\le4$, since their power-sum
denominators are at least $k^{-3}$. These facts give the remaining bounds.
\end{definition}

\begin{definition}[Good columns]
\label{def:good-columns}
A column $j$ is $\chi$-good if $p^{(j)}$ has at least two coordinates at least $\chi$. The sequence satisfies the local $(\chi,\gamma,\rho)$-good-shuffle condition on $[r,T]$ if every macroscopic interval $I\subset[r,T]$ of length $\rho$ contains at least $\gamma\rho\log N$ indices $j$ with $j/\log N\in I$ that are $\chi$-good.
\end{definition}

\begin{definition}[Arbitrarily fine profile and local-good approximation]
\label{def:ub-arbitrarily-fine-approximation}
A family of deterministic cut arrays, one for each $N$, has this property if, for every
$T>0$ and every $\chi_*>0$, there are a fixed
$0<\chi\le\chi_*$, a finite $\chi$-regular partition $\mathcal P$ with
the common constant $C_k$, and fixed $g,\gamma,\bar\rho>0$, with
$g<1/k$, such that the following holds.  For every fixed
$0<\rho\le\bar\rho$ and every fixed $\varphi>0$, all sufficiently
large $N$ satisfy both the $(\mathcal P,\rho,\varphi)$-almost-profile
condition and the local $(g,\gamma,\rho)$-good-shuffle condition on
$[0,T]$.  The threshold in $N$ may depend on these fixed choices.
The scalar cell integrands of the selected partition satisfy the
standing integrability convention.  For each $N$ the cuts form one
fixed sequence; its restrictions at different horizons agree.
\end{definition}

Here $g$ and $\gamma$ are fixed before $\rho$. Thus the local-good
constants do not deteriorate as the chosen mesh decreases. This is
the precise approximation hypothesis for the critical-time result;
the buffered theorem below needs only one sufficiently fine fixed
partition and sufficiently small fixed mesh.

\begin{theorem}[Upper critical-time theorem, reduced form]
\label{thm:upper-critical-reduced}
Let $(\mu_t)_{t\ge0}$ satisfy the standing ordinary-integral convention,
and let the deterministic cut sequence satisfy
Definition~\ref{def:ub-arbitrarily-fine-approximation} on every compact
time interval.  Then, for every finite $T>T_{\UB}^{\mathrm{red}}$,
\[
 d_{\TV}(P(n,\lfloor T\log N\rfloor),\nuUnif)\longrightarrow0.
\]
Consequently, for every fixed $\varepsilon\in(0,1)$,
\[
 \limsup_{N\to\infty}
 \frac{t_{\mathrm{mix}}^{(N)}(\varepsilon)}{\log N}
 \le T_{\UB}^{\mathrm{red}}.
\]
When the negativity set is empty, its critical time is $+\infty$ and
these assertions impose no finite upper bound.
\end{theorem}

The proof of Theorem~\ref{thm:upper-critical-reduced} is given after
Theorem~\ref{thm:dominance-reduction} and uses the buffered estimate
in Theorem~\ref{thm:deterministic-upper-bound}.



\section{Upper bound}

\begin{theorem}[Buffered deterministic upper bound]
\label{thm:deterministic-upper-bound}
Fix $k\ge2$, $T>0$ and $\eta>0$.
There is $\chi_0>0$, depending only on these three quantities, such
that the following holds.  Choose a fixed
$0<\chi\le\chi_0$ and any finite $\chi$-regular partition $\mathcal P$
with the common regularity constant $C_k$, and then fix
$0<g<1/k$ and $\gamma>0$.  There is $\rho_0>0$, and
for every fixed $0<\rho\le\rho_0$ there is $\varphi_0>0$, such that
for every fixed $0<\varphi\le\varphi_0$ the following conclusion holds.
Let $(\mu_t)_{0\le t\le T}$ satisfy the standing scalar-integrability
convention for $a$, $\psi_\theta$, $2\le\theta\le4$, and the cells of
this same partition $\mathcal P$.  Suppose
\[
 \max\{\Lambda_0(T),\Lambda_1(T)\}\le-\eta.
\]
For $K_N=\lfloor T\log N\rfloor$, consider deterministic cuts satisfying
Definition~\ref{def:almost-time-profile-like-cuts} with
$(\mathcal P,\rho,\varphi)$ and the local $(g,\gamma,\rho)$-good-shuffle
condition on $[0,T]$.  Then
\[
 d_{\TV}(P(n,K_N),\nuUnif)\longrightarrow0
\]
uniformly over all such cuts.  Explicitly, for every $\varepsilon>0$
there is $N_0$ such that the distance is at most $\varepsilon$ for all
$N\ge N_0$ and every admissible cut array.  All parameters displayed
before $N_0$ are fixed as $N\to\infty$.
\end{theorem}

Theorem~\ref{thm:deterministic-upper-bound} is proved after
Theorem~\ref{thm:ub-empirical-upper}. We first prove a bound directly
from finite column costs, including the estimates needed after
conditioning on the natural sorted-word histories.


\subsection{Shuffle graph and truncation}

Sample an $N\times K$ shuffle matrix $S$ uniformly among matrices whose $j$th
column contains exactly $n_\ell^{(j)}$ entries equal to $\ell$, and let $S'$
be an independent copy. Sort the row strings of $S$ lexicographically as
$s_1\le\cdots\le s_N$. The shuffle graph $G=G(S)$ is the subgraph of the
path on $[N]$ with edge $(i,i+1)$ present when $s_i=s_{i+1}$; write
$G'=G(S')$ for the independent graph. For a graph $G$ and
$\sigma\in\mathfrak S_N$, let $f_{G,\sigma}$ denote the density at
$\sigma$, relative to $\nuUnif$, of the law obtained by sorting a uniform
permutation increasingly on every connected component of $G$. For an
independent $\sigma\sim\nuUnif$, set
\[
  f_{G,G'}=\E_\sigma[f_{G,\sigma}f_{G',\sigma}].
\]

\begin{lemma}[Graph representation]
\label{lem:ub-graph-representation}
Let $G$ be the shuffle graph generated by the deterministic cut sequence, and independently let $\pi\sim\nuUnif$. Then the inverse $K$-shuffle has the same law as $\pi_G$, obtained by sorting $\pi$ increasingly inside each connected component of $G$.
\end{lemma}

\begin{proof}
One inverse shuffle can be sampled by assigning digits uniformly to the
current card positions with the prescribed pile counts, and then stably
sorting by digit. This follows by inverting the forward interleaving:
within each pile the original order is retained, and all words with the
prescribed digit counts are equally likely. Inverting the product of the
forward shuffles applies these one-step inverse operations in reverse
chronological order. At every such operation, pull the new digit assignment
back to the original card labels. Conditional on the preceding operations,
this is just a permutation of a uniform fixed-count column. Its law does
not depend on that conditioning. Thus these columns, indexed by the original
shuffle times, are independent uniform fixed-count columns.

Successive stable sorts in reverse chronological order sort the resulting
row words lexicographically, with the earlier column taking precedence;
ties retain increasing original card labels. Conditional on the multiset
of row words, row exchangeability makes their assignment to the \(N\)
original labels uniform among all assignments with that multiset. Equivalently,
assign a uniform permutation of the labels to the ordered word occurrences,
and sort the labels increasingly within each block of identical words.
Indeed every assignment of label sets to these blocks has the same
\(\prod_C|C|!\) preimages under this procedure. The blocks are exactly
the components of \(G\), so the conditional law is that of \(\pi_G\).
It depends on the word multiset only through \(G\), which proves the stated
representation with \(\pi\) independent of \(G\).
\end{proof}

\begin{lemma}[Truncated chi-square bound]
\label{lem:ub-truncated-chi-square}
Let $S_1$ be a measurable set of single shuffle matrices, and use the same
set for the two independent matrices $S,S'$. Then
\begin{equation}
\label{eq:ub-truncated-chi-square}
  \dTV(P(n,K),\nuUnif)
  \le \frac12\left(\E[(f_{G,G'}-1)
             \1_{\{S\in S_1,S'\in S_1\}}]\right)^{1/2}
       +\Pp(S\notin S_1).
\end{equation}
The expression under the square root is nonnegative.
\end{lemma}

\begin{proof}
Put $a=\Pp(S\in S_1)$ and
$h(\sigma)=\E[f_{G,\sigma}\1_{\{S\in S_1\}}]$.
The density of $P(n,K)$ is $h+r$, where $r\ge0$ and
$\E_{\nuUnif}r=1-a$. Since $\E_{\nuUnif}h=a$, the triangle inequality
and Cauchy--Schwarz give
\[
 \dTV(P(n,K),\nuUnif)
 \le\tfrac12\E_{\nuUnif}|h-a|+1-a
 \le\tfrac12\bigl(\E_{\nuUnif}(h-a)^2\bigr)^{1/2}+1-a.
\]
Independence of $S,S'$ and finite summation give
\[
 \E_{\nuUnif}(h-a)^2
 =\E[f_{G,G'}\1_{\{S\in S_1,S'\in S_1\}}]-a^2
 =\E[(f_{G,G'}-1)\1_{\{S\in S_1,S'\in S_1\}}].
\]
This proves both assertions. No pointwise nonnegativity of $f_{G,G'}-1$
is required.
\end{proof}

For a path subgraph $H$, let $\mathcal C(H)$ be its connected components,
including isolated vertices, and put
\[
 c(H)=\prod_{C\in\mathcal C(H)}|C|!.
\]
For an integer $B\ge2$, set
\begin{equation}
\label{eq:ub-component-constant}
 d_B=\binom{B}{\lfloor B/2\rfloor}.
\end{equation}

\begin{lemma}[Component factorial bound]
\label{lem:ub-component-factorial}
For any two path subgraphs $G,H$ on $[N]$,
\begin{equation}
\label{eq:ub-factorial-overlap-identity}
 f_{G,H}=\frac{c(G)c(H)}{c(G\cup H)}.
\end{equation}
If every component of $H$ has at most $B$ vertices, then
\begin{equation}
\label{eq:ub-component-factorial-bound}
 f_{G,H}\le d_B^{|E(G)\cap E(H)|}.
\end{equation}
In particular, edge-disjoint $G,H$ satisfy $f_{G,H}\le1$.
\end{lemma}

\begin{proof}
Sorting a uniform permutation on the components of $G$ gives density
$c(G)$ on permutations increasing across every edge of $G$, and zero
elsewhere. Each permutation in this support has exactly $c(G)$
preimages. A uniform permutation is increasing across $G\cup H$ with
probability $1/c(G\cup H)$, proving
\eqref{eq:ub-factorial-overlap-identity}.

First suppose the edge sets are disjoint. In a union component of $v$
vertices, let $\ell_1,\ldots,\ell_s$ be the lengths of the consecutive
monochromatic edge runs. Their sum is $v-1$, and the numerator contributed
by that component is $\prod_i(\ell_i+1)!$. For positive integers $a,b$,
\[
 a!b!\le(a+b-1)!;
\]
indeed $\binom{a+b-1}{a}\ge b$. Applying this inequality successively
shows that the numerator is at most $v!$. Isolated vertices contribute
one, so \eqref{eq:ub-factorial-overlap-identity} is at most one.

For general $G,H$, delete the common edges one at a time from $H$.
The union with $G$ is unchanged. A deletion splits one component into
parts of sizes $a,b\ge1$, with $a+b\le B$, and changes $c(H)$ by the
factor $\binom{a+b}{a}\le d_B$. The last inequality follows because the
largest binomial coefficient in row $m$ is
$\binom m{\lfloor m/2\rfloor}$ and these maxima are nondecreasing in
$m$, by Pascal's identity. The graph remaining after all deletions is
edge-disjoint from $G$. Multiplying the deletion factors proves
\eqref{eq:ub-component-factorial-bound}.
\end{proof}

\begin{definition}[Local sparsity]
\label{def:local-sparsity}
For an integer $L\ge10$, a shuffle graph is $L$-sparse if every interval
of $L$ consecutive vertices contains at most $L/3$ graph edges.
\end{definition}

\begin{lemma}[Sparse factorial bound]
\label{lem:ub-sparse-factorial}
If $G,G'$ are $L$-sparse, then
\[
 f_{G,G'}\le d_L^{|E(G)\cap E(G')|}
            \le(L!)^{|E(G)\cap E(G')|}.
\]
\end{lemma}

\begin{proof}
Every component of an $L$-sparse graph has fewer than $L$ vertices:
a component containing $L$ vertices would give $L-1>L/3$ edges in an
$L$-vertex interval. If $N<L$, the component bound is automatic.
Apply Lemma~\ref{lem:ub-component-factorial} with $B=L$.
\end{proof}

For a full word $w\in[k]_0^K$, define
\[
 C_w=\#\{i:S_i=w\},\qquad
 \lambda_w=\prod_{j=1}^Kp^{(j)}_{w[j]},\qquad
 \mu_*=N\max_w\lambda_w=N\prod_{j=1}^K\max_\ell p^{(j)}_\ell.
\]
Empty products equal one when $K=0$.

\begin{lemma}[Largest-component tail]
\label{lem:ub-largest-component-tail}
For integers $2\le q\le N$,
\begin{equation}
\label{eq:ub-largest-component-tail}
 \Pp\bigl(\max_w C_w\ge q\bigr)
 \le\binom Nq\left(\frac{\mu_*}{N}\right)^{q-1}
 \le\frac{N\mu_*^{q-1}}{q!}.
\end{equation}
For $q>N$ the probability is zero. Consequently, if
$\mu_*\le N^{-\eta}$ for fixed $\eta>0$, then for every fixed
$A>0$ one can choose an integer $B\ge2$ such that
\begin{equation}
\label{eq:ub-bounded-component-probability}
 \Pp\bigl(G\text{ has a component with more than }B\text{ vertices}\bigr)
 \le \frac{N^{-A}}{(B+1)!}.
\end{equation}
\end{lemma}

\begin{proof}
For $q$ distinct labels, the probability that all their words equal $w$
is
\[
 \prod_{j=1}^K\frac{(n^{(j)}_{w[j]})_q}{(N)_q}
 \le\lambda_w^q.
\]
For $0\le n\le N$, the inequality $(n)_q/(N)_q\le(n/N)^q$
follows factor by factor when $n\ge q$; when $n<q$ the numerator
vanishes. Hence
\[
 \E\binom{C_w}{q}\le\binom Nq\lambda_w^q.
\]
Sum over $w$, use $\sum_w\lambda_w=1$, and apply Markov's inequality.
The components of $G$ are precisely the nonempty equal-word blocks.
Taking $q=B+1$ and choosing $\eta B\ge A+1$ proves
\eqref{eq:ub-bounded-component-probability}; the case $B+1>N$ is
already zero.
\end{proof}

\begin{proposition}[Collision upper bound]
\label{prop:ub-collision-upper}
For $N\ge2$, every deterministic pile sequence satisfies
\begin{equation}
\label{eq:ub-collision-upper}
 \dTV(P(n,K),\nuUnif)
 \le\binom N2\prod_{j=1}^K
       \sum_{\ell=0}^{k-1}\frac{(n^{(j)}_\ell)_2}{(N)_2}
 \le\binom N2\exp\left\{-\sum_{j=1}^K\psi_{p^{(j)}}(2)\right\}.
\end{equation}
In particular, if $\sum_j\psi_{p^{(j)}}(2)\ge(2+\zeta)\log N$
for fixed $\zeta>0$, then the distance is at most $N^{-\zeta}/2$.
\end{proposition}

\begin{proof}
On the event that no two full words agree, $G$ has no edges and the
independent permutation in Lemma~\ref{lem:ub-graph-representation} is
unchanged. Thus total variation is at most the probability of a full-word
collision. A union bound over label pairs gives the first inequality in
\eqref{eq:ub-collision-upper}. For each digit,
$(n)_2/(N)_2\le(n/N)^2$; summing over digits and using
\eqref{eq:one-column-costs} gives the second inequality.
\end{proof}

\subsection{An empirical first moment with a reserved population}

We retain the original scale $N$ when applying a bound to a smaller deck.
Fix integers $2\le n\le N$, $0\le a<K\le T\log N$, and independent
uniform fixed-count columns on $n$ labels, indexed by $a<j\le K$.
Write $q^{(j)}\in D_k$ for their proportions and $L=\log N$.
A column is $g$-good in the sense of Definition~\ref{def:good-columns},
with proportions relative to its current $n$ labels. For integers
$a\le i\le j\le K$, define
\begin{equation}
\label{eq:ub-empirical-costs}
 \widehat A(i,j)=\frac1L\sum_{h=i+1}^j a_{q^{(h)}},\qquad
 \widehat B_\theta(i,j)=\frac1L\sum_{h=i+1}^j\psi_{q^{(h)}}(\theta),
 \qquad \widehat\beta(s)=\log_N n-\widehat A(a,s).
\end{equation}
The empirical negative-margin condition, with margin $\eta>0$, is
\begin{equation}
\label{eq:ub-empirical-negative-margin}
 \widehat\beta(K)\le-\eta,\qquad
 \inf_{3\le\theta\le4}
 \left\{\frac{\theta+3}{2}\widehat\beta(s)
       -\widehat B_\theta(s,u)-2\widehat B_2(u,K)\right\}
 \le-\eta\quad(a\le s\le u\le K).
\end{equation}
All endpoints in \eqref{eq:ub-empirical-negative-margin} are integers.
No time profile or simplex partition is used in this condition.

\begin{proposition}[Empirical first moment]
\label{prop:ub-empirical-first-moment}
\label{prop:ub-first-moment}
Fix $k,T,\eta,g>0$, with $k\ge2$ and $g<1/k$.
There exist $c>0$ and $N_0$, depending only on these fixed parameters,
such that the following holds for every $N\ge N_0$. If the first column
$q^{(a+1)}$ is $g$-good and
\eqref{eq:ub-empirical-negative-margin} holds, then independent shuffle
graphs $G,G'$ on the $n$ labels satisfy
\begin{equation}
\label{eq:ub-empirical-first-moment}
 \E|E(G)\cap E(G')|\le N^{-c}.
\end{equation}
The bound is uniform in $n,a,K$ and the deterministic cut arrays.
\end{proposition}

The first column fixes the endpoints of its piles. Its deterministic
counts are therefore excluded from all fluctuation scales. We condition
on its assignment to labels. For each nonempty pile $R$, its size $M$
satisfies $M\le(1-g)n$. All remaining columns still use the full
population of $n$ labels; the labels outside $R$ are a population reserve.
Row-permutation invariance shows that the conditional sorted graph law
in a pile depends only on its deterministic size and rank offset, not
on the particular labeled set $R$. Thus, after exposing the two
independent first-column assignments, their corresponding pile graphs
still have the same law and are independent.

Here and below, a word is legal when all of its selected probabilities
are positive. A legal prefix $x$ beginning in $R$ has absolute final
column $M_x$, mean count
\[
 \mu_x=n\prod_{a<j\le M_x}q^{(j)}_{x[j]},\qquad
 \ell_x=\log_N\mu_x,\qquad e_x=-\widehat B_2(M_x,K).
\]
The associated block $B_x$ consists of its full-word extensions. Its
sorted rank interval is denoted by $\mathcal I(B_x)=[\iota_x,\tau_x)$,
with one-based ranks and $\tau_x=\iota_x+C_x$, where $C_x$ is its count.
The graph $G_{B_x}$ contains the equal-word edges internal to that block.

Fix a small cutoff $0<\chi<1/k$. For each later column, put
\[
 l_j=\min\{d:q^{(j)}_d\ge\chi\},\qquad
 r_j=\max\{d:q^{(j)}_d\ge\chi\}.
\]
For a full word beginning in $R$, its left boundary $y$ stops at the first
column $j>a+1$ where its digit is greater than $l_j$; its right boundary
$z$ stops at the first column where its digit is less than $r_j$.
If the corresponding crossing never occurs, use the full word instead.
Let $v$ be the longer of $y,z$. A pair is terminal if at least one
boundary stops only at column $K$; it will be treated directly.
Every other boundary pair has two genuine crossings. Within its join
$B_v$, the pair $(y,z)$ is fixed.

For a nonterminal pair, at every column after the first and up to $M_v$,
let $P_j$ and $Q_j$ be the subsets of $R$ whose strings first differ
from $v$ at column $j$ by a smaller or larger digit, respectively.
Their counts and expectations are taken in the fixed-column model.
Define
\begin{equation}
\label{eq:ub-anchored-fluctuation}
 L_v=\max_{a+1<j\le M_v}\E|P_j|,\qquad
 R_v=\max_{a+1<j\le M_v}\E|Q_j|,\qquad
 F_v=\frac12\log_N\min(L_v,R_v),\qquad V_v=\ell_v.
\end{equation}
A genuine crossing gives a side mass at least $\chi\mu_v$ on each
side. Thus both maxima are positive and
\begin{equation}
\label{eq:ub-anchored-fluctuation-comparison}
 2F_v\ge V_v-q_N,\qquad q_N=\frac{\log(1/\chi)}{\log N}.
\end{equation}
The scale $F_v$ remains fixed for every descendant of $v$.

Choose a small stability parameter $\delta>0$. If $V_v-F_v<\delta$,
use just the block $B_v$. Otherwise, descend from $v$ and stop each
branch at its first prefix $x$ with $\ell_x-F_v<\delta$.
The resulting blocks partition $B_v$: the terminal margin gives
$\ell_x\le-\eta$ at depth $K$, while
\eqref{eq:ub-anchored-fluctuation-comparison} and $V_v\ge\ell_x$ imply
$\ell_x-F_v\le-\eta/2+q_N/2<\delta$ there for large $N$.
A block other than the immediate block $v$ is called proper. Its parent, denoted $x^-$,
has mean at least $N^{F_v+\delta}$. Moreover,
\begin{equation}
\label{eq:ub-proper-fluctuation-lower}
 F_v\ge\delta-q_N\ge\delta/2
\end{equation}
for large $N$, because the join was not itself stopped.

\begin{lemma}[Hypergeometric concentration and atoms]
\label{lem:ub-hypergeometric-estimates}
Write $\operatorname{Hyp}(r,n,m)$ for the number of marked positions
in a uniform sample of $m$ positions from a population of $n$ with
$r$ marked positions. A hypergeometric variable with population $n\ge2$, marked size $r$,
and sample size $m$ has mean $\mu=rm/n$ and variance
$\sigma^2=rm(n-r)(n-m)/(n^2(n-1))$. It satisfies
\[
 \Pp(|Y-\mu|\ge t)\le2\exp\{-t^2/[4(\mu+t)]\}\quad(t>0),
 \qquad \sup_z\Pp(Y=z)\le C/\sigma\quad(\sigma>0),
\]
where $C$ is absolute. The concentration bound also holds for the count
of a fixed prefix in a deterministic row set, with its actual mean.
\end{lemma}

\begin{proof}
The exponential transform and tail proof in
Lemma~\ref{lem:successive-thinning-tail} applies to successive independent
uniform subsets of the full label population. A fixed initial row set
can be replaced by an independent uniform set of the same size without
changing the count law, by permutation invariance. This proves the
concentration assertions.

For the atom bound, the explicit hypergeometric probabilities have
successive ratios decreasing on their support, so they form a
log-concave sequence. We give the needed general estimate. Let $p_*$
be the maximum of a log-concave integer probability mass function
with interval support,
attained at $z_0$, and set $h=\lceil4/p_*\rceil$. Unimodality and total
mass one imply $p(z_0+h),p(z_0-h)\le p_*/2$, with values outside the
support set to zero. Log-concavity then gives
$p(z_0\pm j)\le2p_*2^{-j/h}$ for every $j\ge0$. Summing the geometric
second moments yields
\[
 \Var(Y)\le\E(Y-z_0)^2\le C_0p_*h^3\le C_1/p_*^2.
\]
Rearranging proves $p_*\le\sqrt{C_1}/\sigma$.
\end{proof}


For the next three lemmas, use local coordinates after the exposed
first column: local time $j$ means absolute column $a+1+j$, and
$d=K-a-1$. A local prefix denotes the tail of the corresponding
anchored prefix. Its count, mean, endpoints and internal graph are
unchanged by this deletion of its fixed first digit. We take
$g_{\rm anchor}=g$ and later $\chi_{\rm cross}=\chi$.
If $d=0$, all graph components have sizes equal to the first-column
pile counts, so the terminal negative margin is impossible for a
nonempty pile and sufficiently large $N$. Thus only $d\ge1$ occurs.

The smoothing argument uses the population outside an exposed anchor.
Here is the precise conditional model.  Fix integers $N\ge n\ge2$ and
$d\ge1$.  Conditionally on an initial sigma-field $\mathcal H_0$, let
$R\subseteq[n]$, its size $M\ge1$, and an integer offset $c_0$ be fixed.
For $1\le j\le d$, the column $(\xi_{r,j})_{r\in[n]}$ is uniform among
all arrangements of its deterministic counts $(n_{j,a})_{a\in[k]_0}$,
and these columns are conditionally independent.  Write
$p_{j,a}=n_{j,a}/n$.  Assume
\begin{equation}
\label{eq:ub-anchor-reserve}
  M\le(1-g_{\rm anchor})n,\qquad 0<g_{\rm anchor}<1.
\end{equation}
For example, expose the first column of the current $n$-row model and
take $R$ to be one of its piles.  If its largest pile has size at most
$(1-g_{\rm anchor})n$, then \eqref{eq:ub-anchor-reserve} holds for every
nonempty choice of this pile.  The remaining columns still have the
conditional law just specified.  All statements below are conditional
on $\mathcal H_0$; in their proofs we fix an atom of this finite
sigma-field and suppress this conditioning from the notation.

For a word $x=(x_1,\ldots,x_m)$ of length $m\le d$, set
\[
 B_x=\{r\in R:\xi_{r,j}=x_j\text{ for }1\le j\le m\},\qquad
 C_x=|B_x|,\qquad \mu_x=M\prod_{j=1}^m p_{j,x_j}.
\]
Only legal words, for which $\mu_x>0$, are considered.  Put
\[
 \iota_x=c_0+\#\{r\in R:(\xi_{r,1},\ldots,\xi_{r,m})<_{\rm lex}x\},
 \qquad \tau_x=\iota_x+C_x,
 \qquad I_x=[\iota_x,\tau_x)\cap\mathbb Z.
\]
Thus the integer slots are numbered from $c_0$; changing to ranks
numbered from $c_0+1$ translates every interval by one and changes none
of the assertions.  Let $\mathcal U_x$ record the labels in $B_x$ and
all their digits after time $m$.  In particular $C_x$ is
$\mathcal U_x$-measurable.

For $1\le i\le m$, let $P_{x,i}$ and $Q_{x,i}$ be the subsets of $R$
whose digits agree with $x$ before $i$ and whose $i$-th digit is,
respectively, less than or greater than $x_i$.  Write
\[
 p_i(<x_i)=\sum_{a<x_i}p_{i,a},\qquad
 p_i(>x_i)=\sum_{a>x_i}p_{i,a},
\]
\[
 \ell_i=M\prod_{j<i}p_{j,x_j}\,p_i(<x_i),\qquad
 r_i=M\prod_{j<i}p_{j,x_j}\,p_i(>x_i).
\]
These are exactly the conditional means of $|P_{x,i}|$ and $|Q_{x,i}|$.
An exposed column is never included in these side pieces or in a
fluctuation maximum.

\begin{lemma}[Cylinder concentration in a fixed anchor]
\label{lem:ub-anchor-cylinder-tail}
In the conditional model just specified, prescribe an arbitrary set of
allowed digits at each of any collection of future columns.  If $X$ is
the number of rows of $R$ satisfying these prescriptions and $\mu=\E X$,
then, for $s>0$ and $t>0$,
\begin{equation}
\label{eq:ub-anchor-cylinder-transform}
 \E s^X\le(1-\mu/M+(\mu/M)s)^M\le e^{\mu(s-1)},
\end{equation}
\begin{equation}
\label{eq:ub-anchor-cylinder-tail}
 \Pp(|X-\mu|\ge t)\le
 2\exp\left(-\frac{t^2}{4(\mu+t)}\right),\qquad
 \Pp(X\le\mu/2)\le e^{-c_{\rm Ch}\mu},
 \quad c_{\rm Ch}=\frac{1-\log2}{2}>0.
\end{equation}
The same bounds hold for a hypergeometric intersection with any fixed
set in the $n$-point population, with its exact hypergeometric mean.
\end{lemma}

\begin{proof}
The allowed positions in each prescribed column form a uniform subset
of $[n]$ of some deterministic size $nq_j$, independently across
columns.  If the current intersection has $z$ rows, its size after
the next column has law $\operatorname{Hyp}(nq_j,n,z)$.  For a
positive real number $s$, its transform is the elementary symmetric
mean of order $z$ of a list with $nq_j$ copies of $s$ and $n(1-q_j)$
copies of one.  This mean is at most $(1-q_j+q_js)^z$.
Indeed, averaging any two entries preserves their sum and increases
their product, so it cannot decrease the elementary symmetric
polynomial.  Repeatedly averaging a largest and a smallest entry
converges to the constant list of their common arithmetic mean:
if $V$ is the sum of squared deviations, an averaging reduces it by
half the squared range, which is at least $V/(2n)$.  Continuity proves
the asserted inequality, also for $z=0,1$.

Start with the fixed set $R$ of size $M$.  Backward iteration of the
affine transformations $s\mapsto1-q_j+q_js$ gives
\[
 \E s^X\le(1-Q+Qs)^M,\qquad Q=\prod_jq_j,\qquad \mu=MQ.
\]
This proves \eqref{eq:ub-anchor-cylinder-transform}.  For $0\le h\le1$,
$e^{h}-1-h\le h^2$ and $e^{-h}-1+h\le h^2$; consequently
$\E e^{\pm h(X-\mu)}\le e^{\mu h^2}$.
Markov's inequality with $h=t/(2(\mu+t))$ gives the first bound in
\eqref{eq:ub-anchor-cylinder-tail}.  The choice $s=1/2$ in
\eqref{eq:ub-anchor-cylinder-transform}, followed by Markov's inequality
at $2^{-\mu/2}$, gives its second bound.  One prescribed column gives
the hypergeometric assertion.  The argument starts from the fixed
set $R$, but every sampling population is $[n]$.
\end{proof}

Fix compatible boundary words $y,z$ that are prefixes of a word $x$.
Let $h_- = |y|$, $h_+=|z|$, $h=\max(h_-,h_+)$ and let $v$ be their
longer word.  Suppose both boundary words are genuine crossings, in
the following precise sense: for a fixed $0<\chi_{\rm cross}<1$,
\begin{equation}
\label{eq:ub-anchor-genuine-crossings}
 p_{h_-}(<y_{h_-})\ge\chi_{\rm cross},\qquad
 p_{h_+}(>z_{h_+})\ge\chi_{\rm cross}.
\end{equation}
Set
\begin{equation}
\label{eq:ub-anchor-side-scales}
 \ell=\max_{1\le i\le h}\ell_i,\qquad
 r=\max_{1\le i\le h}r_i,\qquad
 \omega=\min(\ell,r),\qquad
 D=2h+\chi_{\rm cross}^{-1}.
\end{equation}
All these quantities depend only on $(y,z)$, not on the continuation
of $x$ after $h$.  Since $v$ is legal,
\begin{equation}
\label{eq:ub-anchor-side-comparison}
 0<\chi_{\rm cross}\mu_v\le\omega,
 \qquad M\le D\max(\ell,r).
\end{equation}
Indeed, the crossing side at $h_-$ has mean at least
$\chi_{\rm cross}\mu_v$, and the crossing at $h_+$ gives the other
inequality.  Partitioning the rows according to their first difference
from $v$ gives the exact mean identity
$M=\sum_{i\le h}(\ell_i+r_i)+\mu_v$, which proves the second
inequality in \eqref{eq:ub-anchor-side-comparison}.

\begin{lemma}[Anchored endpoint anti-concentration]
\label{lem:ub-anchor-endpoint-atoms}
Assume \eqref{eq:ub-anchor-reserve} and
\eqref{eq:ub-anchor-genuine-crossings}.  For every legal continuation
$x$ of $v$, there is a sigma-field
$\mathcal F_x\supseteq\sigma(\mathcal H_0,\mathcal U_x)$ and an
$\mathcal F_x$-measurable event $\mathcal S_x$ such that
\begin{equation}
\label{eq:ub-anchor-smoothing-failure}
 \Pp(\mathcal S_x^c\mid\mathcal H_0)
 \le3e^{-c_{\rm Ch}g_{\rm anchor}\omega}.
\end{equation}
Conditionally on $\mathcal F_x$, the endpoints have the form
\begin{equation}
\label{eq:ub-anchor-endpoint-shift}
 \iota_x=A_x+Y_x,\qquad \tau_x=A_x+Y_x+C_x,
\end{equation}
where $A_x,C_x$ are $\mathcal F_x$-measurable and $Y_x$ is an exact
hypergeometric variable.  On $\mathcal S_x$,
\begin{equation}
\label{eq:ub-anchor-variance}
 \Var(Y_x\mid\mathcal F_x)
 \ge\frac{g_{\rm anchor}\omega}{16D}.
\end{equation}
If $C_{\rm hyp}$ is the absolute point-mass constant in
Lemma~\ref{lem:ub-hypergeometric-estimates}, then almost surely
\begin{equation}
\label{eq:ub-anchor-conditional-atoms}
 \max\left\{
 \sup_{a\in\mathbb Z}\Pp(\iota_x=a,\mathcal S_x\mid\mathcal U_x,\mathcal H_0),
 \sup_{a\in\mathbb Z}\Pp(\tau_x=a,\mathcal S_x\mid\mathcal U_x,\mathcal H_0)
 \right\}
 \le4C_{\rm hyp}\sqrt{\frac{D}{g_{\rm anchor}\omega}}.
\end{equation}
\end{lemma}

\begin{proof}
We expose all but one future column, and only part of that last column.
The unused population outside $R$ supplies the sample-complement
factor in the variance.  Relative lower tails for the remaining
marked and unmarked counts supply its other factors; no subtraction
of two approximate means is used.

Reverse the order of the alphabet if necessary so that
$\ell\le r$.  This interchanges left and right endpoints, up to a
fixed reflection inside $R$.  Reflection changes a shift plus a
hypergeometric variable $Y$ into a shift minus $Y$.  If its sample
size is $b$, replace $Y$ by $b-Y$, which is hypergeometric with
marked and unmarked counts interchanged and has the same variance;
the sample size will be measurable in every revealing sigma-field
below.  Thus it suffices to prove the assertion in this order.
Choose $i_-,i_+\le h$ attaining $\ell$ and $r$.
Thus $\omega=\ell$ and $M\le Dr$.

Suppose first that $i_->i_+$, and put $j=i_+$,
$d_j=p_{j,x_j}$, $p=p_j(\le x_j)$ and $q=1-p=p_j(>x_j)$.
Reveal all columns other than $j$.  The following row sets are then
known:
\[
 T=\{a\in R:\xi_{a,i}=x_i\text{ for }i<j\},
\]
\[
 B=\{a\in T:(\xi_{a,j+1},\ldots,\xi_{a,i_-})
                    <_{\rm lex}(x_{j+1},\ldots,x_{i_-})\}.
\]
Put $A=T\setminus B$.  Also reveal column $j$ on $A$, and let $a_3$ be the number of its
revealed digits at most $x_j$.  Conditional on this reveal,
\begin{equation}
\label{eq:ub-anchor-hyp-case-right}
 Y_x=\#\{a\in B:\xi_{a,j}\le x_j\}
 \sim\operatorname{Hyp}(np-a_3,n-|A|,|B|).
\end{equation}
Every row in $R\setminus T$ has its order relative to $x$ decided
before column $j$.  A row in $B$ precedes $x$ exactly when its digit
at $j$ is at most $x_j$.  The order of every row in $A$ is known.
These facts give \eqref{eq:ub-anchor-endpoint-shift}.  All rows with
prefix $x$ belong to $A$, and their later digits have been exposed;
hence this revealing sigma-field contains $\mathcal U_x$.

Let $B_0\subseteq B$ consist of rows matching $x$ before $i_-$ except
at column $j$, with their digit at $i_-$ less than $x_{i_-}$.
This is a cylinder with column $j$ omitted and mean $\ell/d_j$.
Require the three inequalities
\[
 |B_0|\ge\frac{\ell}{2d_j},\qquad
 np-a_3\ge\frac{(n-|A|)p}{2},\qquad
 n(1-p)-(|A|-a_3)\ge\frac{(n-|A|)q}{2}.
\]
They define an event in the revealing sigma-field.  The first failure
has probability at most $e^{-c_{\rm Ch}\ell/d_j}$ by
Lemma~\ref{lem:ub-anchor-cylinder-tail}.  Conditional on the other
columns, the two remaining counts are hypergeometric with means
$(n-|A|)p$ and $(n-|A|)q$.  These means are at least
$g_{\rm anchor}np$ and $g_{\rm anchor}nq$.
Here $np\ge nd_j\ge\ell$ and $nq\ge r\ge\ell$.
Lemma~\ref{lem:ub-anchor-cylinder-tail} and a union bound therefore
give \eqref{eq:ub-anchor-smoothing-failure}.

On this event, the marked and unmarked fractions in
\eqref{eq:ub-anchor-hyp-case-right} are at least $p/2$ and $q/2$.
Since $A\cup B\subseteq R$, its sample-complement fraction is at least
\[
 \frac{n-|A|-|B|}{n-|A|}\ge g_{\rm anchor}.
\]
The variance formula consequently gives
\[
 \Var(Y_x\mid\mathcal F_x)
 \ge\frac{g_{\rm anchor}|B|pq}{4}
 \ge\frac{g_{\rm anchor}\ell q}{8}
 \ge\frac{g_{\rm anchor}\ell}{8D}.
\]
The second inequality uses $p\ge d_j$, and the last uses
$q\ge r/M\ge1/D$.

Suppose next that $i_-<i_+$, and put $j=i_-$,
$d_j=p_{j,x_j}$, $p=p_j(<x_j)$ and $q=1-p$.
Reveal all other columns and set
\[
 T=\{a\in R:\xi_{a,i}=x_i\text{ for }i<j\},
\]
\[
 B=\{a\in T:(\xi_{a,j+1},\ldots,\xi_{a,i_+})
                    >_{\rm lex}(x_{j+1},\ldots,x_{i_+})\}.
\]
Put $A=T\setminus B$.  Reveal column $j$ on $A$, writing $b_3$ for the number of revealed
digits less than $x_j$.  Then
\begin{equation}
\label{eq:ub-anchor-hyp-case-left}
 Y_x=\#\{a\in B:\xi_{a,j}<x_j\}
 \sim\operatorname{Hyp}(np-b_3,n-|A|,|B|).
\end{equation}
For a row in $B$, preceding $x$ is equivalent to having its digit at
$j$ less than $x_j$.  Every prefix-$x$ row lies in $A$.  Thus the
same explicit classification of rows proves
\eqref{eq:ub-anchor-endpoint-shift} and
$\mathcal U_x\subseteq\mathcal F_x$ in this case also.

Now $B$ contains the cylinder obtained from the right piece at $i_+$
by omitting column $j$, whose mean is $r/d_j$.  Require its count to
be at least $r/(2d_j)$, and require the remaining marked and unmarked
counts to be at least $(n-|A|)p/2$ and $(n-|A|)q/2$.
Their three means are at least $r$, $g_{\rm anchor}np$ and
$g_{\rm anchor}nq$.  Since $np\ge\ell$ and
$nq\ge nd_j\ge r\ge\ell$, the failure probability is again at most
$3e^{-c_{\rm Ch}g_{\rm anchor}\ell}$.
The sample-complement fraction is still at least $g_{\rm anchor}$,
so on this event
\[
 \Var(Y_x\mid\mathcal F_x)
 \ge\frac{g_{\rm anchor}|B|pq}{4}
 \ge\frac{g_{\rm anchor}rp}{8}
 \ge\frac{g_{\rm anchor}r\ell}{8M}
 \ge\frac{g_{\rm anchor}\ell}{8D}.
\]
Here $q\ge d_j$ and $p\ge\ell/M$.

Finally suppose $i_-=i_+=j$.  Set
$p=p_j(<x_j)$, $d_j=p_{j,x_j}$ and $q=p_j(>x_j)$.
Let
\[
 T=\{a\in R:\xi_{a,i}=x_i\text{ for }i<j\},\qquad
 \bar a=M\prod_{i<j}p_{i,x_i},
\]
so that $\ell=\bar a p$, $r=\bar a q$ and $p\le q$.
Reveal all columns except $j$, and in column $j$ reveal the entire
global set $E=\{a\in[n]:\xi_{a,j}=x_j\}$.  Put $B=T\setminus E$.
Conditional on this information, the remaining column is uniform
on its $n(p+q)$ unexposed positions, with exactly $np$ digits less
than $x_j$.  Therefore
\begin{equation}
\label{eq:ub-anchor-hyp-case-equal}
 Y_x=\#\{a\in B:\xi_{a,j}<x_j\}
 \sim\operatorname{Hyp}(np,n(p+q),|B|).
\end{equation}
The rows in $T\cap E$ have their order decided by the already exposed
suffixes, those outside $T$ by earlier digits, and those in $B$ by
the indicated strict inequality at $j$.  All prefix-$x$ rows and their
continuations are known.  This proves the shift and measurability
assertions in this final case.

Require
\[
 |T|\ge\bar a/2,\qquad
 |B|\ge |T|(p+q)/2,\qquad
 n(p+q)-|B|\ge(n-|T|)(p+q)/2.
\]
The first failure has probability at most $e^{-c_{\rm Ch}\bar a}$.
Conditional on all columns other than $j$, the two counts in the
last inequalities are intersections of the uniform $(n(p+q))$-set
$[n]\setminus E$ with $T$ and $[n]\setminus T$.  Their means are
$|T|(p+q)$ and $(n-|T|)(p+q)$.  On the first event the former is at
least $\bar a(p+q)/2\ge\ell$; the latter is always at least
$g_{\rm anchor}n(p+q)\ge g_{\rm anchor}\ell$.
Thus their total failure probability, including the first event, is
at most $3e^{-c_{\rm Ch}g_{\rm anchor}\ell}$.
On this event the mean of \eqref{eq:ub-anchor-hyp-case-equal} is at
least $\ell/4$, its unmarked fraction $q/(p+q)$ is at least $1/2$,
and its sample-complement fraction is at least $g_{\rm anchor}/2$.
It follows that
\[
 \Var(Y_x\mid\mathcal F_x)\ge
 \frac{g_{\rm anchor}\ell}{16},
\]
which implies \eqref{eq:ub-anchor-variance} because $D\ge1$.

In every case the good event is measurable in the stated revealing
sigma-field.  On that event the positive lower bounds for marked
and unmarked counts ensure that the hypergeometric population has
at least two positions; hence its variance formula is applicable.
No lower bound such as $\omega\ge1$ is required.  Apply the point-mass
bound of Lemma~\ref{lem:ub-hypergeometric-estimates} on each good atom
of $\mathcal F_x$.  Both shifts in
\eqref{eq:ub-anchor-endpoint-shift} are known there, so each endpoint
has conditional atom at most
$4C_{\rm hyp}\sqrt{D/(g_{\rm anchor}\omega)}$.
Taking conditional expectation back to
$\sigma(\mathcal H_0,\mathcal U_x)$ proves
\eqref{eq:ub-anchor-conditional-atoms}.  This argument does not
condition the entire law on the good event.
\end{proof}

\begin{lemma}[Endpoint windows from the smaller side]
\label{lem:ub-anchor-endpoint-window}
In the same model, fix genuine boundary words $(y,z)$ and the
quantities in \eqref{eq:ub-anchor-side-scales}.  For any legal
continuation $x$ of their join $v$, let
$\bar\iota_x=\E(\iota_x\mid\mathcal H_0)$ and
$\bar\tau_x=\E(\tau_x\mid\mathcal H_0)$.
For every $w>0$,
\begin{equation}
\label{eq:ub-anchor-finite-window}
 \Pp\left(
   \max(|\iota_x-\bar\iota_x|,|\tau_x-\bar\tau_x|)>w
       \,\middle|\,\mathcal H_0\right)
 \le2(d+1)\exp\left\{
 -\frac{w^2}{4(d+1)^2(\chi_{\rm cross}^{-1}\omega+w)}
 \right\}.
\end{equation}
This estimate needs neither the reserve
\eqref{eq:ub-anchor-reserve} nor a lower bound on $M/n$.
\end{lemma}

\begin{proof}
Suppose $\ell\le r$.  The exact endpoint decomposition is
\[
 \iota_x=c_0+\sum_{i\le|x|}|P_{x,i}|,
 \qquad \tau_x=\iota_x+C_x.
\]
For $i\le h$ the mean of a displayed side piece is at most
$\ell=\omega$.  For $i>h$ it is a subset of $B_v$, and its mean
is at most $\mu_v\le\chi_{\rm cross}^{-1}\omega$ by
\eqref{eq:ub-anchor-side-comparison}.  The same bound holds for
$\mu_x$.  Each count is a cylinder count of the type in
Lemma~\ref{lem:ub-anchor-cylinder-tail}; a strict inequality on one
digit is simply an allowed set of digits in that column.
Apply that lemma to these at most $d+1$ counts with threshold
$w/(d+1)$.  If every deviation is at most this threshold, both
endpoint deviations are at most $w$.  A union bound gives
\eqref{eq:ub-anchor-finite-window}, after enlarging its denominator
from $\chi_{\rm cross}^{-1}\omega+w/(d+1)$ to
$\chi_{\rm cross}^{-1}\omega+w$.

If $r<\ell$, use instead
\[
 \tau_x=c_0+M-\sum_{i\le|x|}|Q_{x,i}|,
 \qquad \iota_x=\tau_x-C_x.
\]
The identical argument uses the smaller right-hand maximum
$r=\omega$.  The constant $c_0+M$ is known after the anchor exposure.
\end{proof}

For the counting consequence, use $\log N$ as the normalization even
when the sampling population is $n<N$.  For a legal word $x$ put
$c_L(x)=\log_N\mu_x$.  Fix $\delta>0$ and genuine boundary words
$(y,z)$, write $v$ for their join, and set
\begin{equation}
\label{eq:ub-anchor-normalized-scales}
  V=c_L(v),\qquad F=\tfrac12\log_N\omega,\qquad
  q_N=\frac{\log(1/\chi_{\rm cross})}{\log N}.
\end{equation}
Equation~\eqref{eq:ub-anchor-side-comparison} gives
$2F\ge V-q_N$.  If $V-F<\delta$, stop at the join and use the
single block $v$.  Otherwise a proper stable leaf is a legal strict
extension $x$ of $v$ that first satisfies $c_L(x)-F<\delta$;
its parent, denoted $x^-$, satisfies $c_L(x^-)-F\ge\delta$.
Full-length words that never satisfy this inequality are terminal
fallback words, and are not called proper stable leaves here.

\begin{lemma}[Relevant parents and per-edge sparsity]
\label{lem:ub-anchor-relevant-parents}
Fix $k,T,\delta>0$ and $\chi_{\rm cross}\in(0,1)$ before taking
$N\to\infty$, and assume $d\le K\le T\log N$.
For all sufficiently large $N$, uniformly over the conditional
column counts, anchors and genuine boundary pairs, the following
holds.  For an integer slot $m$, call a proper stable leaf relevant
when
\begin{equation}
\label{eq:ub-anchor-relevance-threshold}
 \Pp(m\in I_x\mid\mathcal H_0)
       \ge\exp(-N^{\delta/16}).
\end{equation}
There are at most $2kK$ relevant proper stable leaves for the fixed
boundary pair.  If the join itself is the stable block, it contributes
at most one block instead.  Moreover, the probability that $m$
belongs to some proper stable leaf but no relevant proper stable
leaf is at most
\begin{equation}
\label{eq:ub-anchor-irrelevant-loss}
 (d+1)k^d e^{-N^{\delta/16}}
 \le e^{-N^{\delta/32}}.
\end{equation}
The assertion applies to all proper stable leaves, regardless of
whether their expected sizes are designated small or large.
\end{lemma}

\begin{proof}
The existence of a proper stable leaf implies that the join was not
stopped.  Therefore $V-F\ge\delta$, and
\eqref{eq:ub-anchor-normalized-scales} gives
\begin{equation}
\label{eq:ub-anchor-proper-positive-fluctuation}
 F\ge\delta-q_N\ge\delta/2
\end{equation}
for all sufficiently large $N$.  This argument uses the first test
at the join, before any extension is explored.

Use the same halfwidth $w=N^{F+\delta/4}$ for every parent in the
fixed boundary pair.  Since $\omega=N^{2F}$ and
\eqref{eq:ub-anchor-proper-positive-fluctuation} holds, the exponent
in \eqref{eq:ub-anchor-finite-window} is at least
\[
 \frac{\chi_{\rm cross}}{8(d+1)^2}N^{\delta/2}
\]
for all sufficiently large $N$.  The factor $2(d+1)$ is harmless
because $d\le T\log N$.  Thus, for every parent $a$ of a proper
stable leaf,
\begin{equation}
\label{eq:ub-anchor-parent-window}
 \Pp\left(
 \max(|\iota_a-\bar\iota_a|,|\tau_a-\bar\tau_a|)>w
       \,\middle|\,\mathcal H_0\right)
 \le e^{-N^{\delta/8}}.
\end{equation}
Also $\bar\tau_a-\bar\iota_a=\mu_a\ge N^{F+\delta}>2w$.

Fix a depth.  Distinct parents at this depth are distinct cylinders
of the same length.  Their deterministic mean intervals
$[\bar\iota_a,\bar\tau_a]$ are therefore ordered and have disjoint
interiors: in lexicographic order the right endpoint of one is at
most the left endpoint of the next.  If $x$ is relevant, then
$I_x\subseteq I_{x^-}$ and the probability in
\eqref{eq:ub-anchor-relevance-threshold} exceeds the bad-window
probability in \eqref{eq:ub-anchor-parent-window}.  Some realization
therefore has both $m\in I_{x^-}$ and the good parent window.  It
follows that
\[
 m\in[\bar\iota_{x^-}-w,\bar\tau_{x^-}+w].
\]
At most two parent intervals at the fixed depth can have their
$w$-enlargements contain $m$.  Indeed, if three did, order them
from left to right.  The right endpoint of the first is at most
the left endpoint of the middle, and the left endpoint of the last
is at least the right endpoint of the middle.  Membership in both
outer enlargements would then imply that the length of the middle
interval is at most $2w$, contradicting $\mu_a>2w$.
Each parent has at most $k$ children, and there are at most $d\le K$
possible parent depths.  This proves the bound $2kK$.

There are at most $(d+1)k^d$ words of length at most $d$.  A union
bound over proper leaves failing
\eqref{eq:ub-anchor-relevance-threshold} proves the first inequality
in \eqref{eq:ub-anchor-irrelevant-loss}; the second follows from
$d\le T\log N$.  If the join is already stable, its family is a
singleton by definition.  Terminal fallback words and boundary
pairs with a nongenuine terminal side are outside this lemma and
must be excluded by a terminal margin or estimated separately.
\end{proof}

In particular, on a proper stable class the anti-concentration estimate
\eqref{eq:ub-anchor-conditional-atoms} has order
\[
 O\left(\sqrt{K+\chi_{\rm cross}^{-1}}\,
                  g_{\rm anchor}^{-1/2}N^{-F}\right),
\]
and \eqref{eq:ub-anchor-smoothing-failure} is a stretched exponential
error.  Events controlling block sizes or endpoint windows may be
intersected with $\mathcal S_x$: doing so only decreases the
conditional atoms, while their failure probabilities are added.
The sampling population throughout this proof is the actual global
population $n$, including the labels outside $R$.  Replacing it by
$M$ or conditioning on all future counts inside $R$ would change the
conditional law and is not part of the argument.

We now return to the absolute indices and notation of
\eqref{eq:ub-empirical-costs} and
\eqref{eq:ub-anchored-fluctuation}. In the preceding lemmas,
$\omega=N^{2F_v}$, $\chi_{\rm cross}=\chi$ and
$g_{\rm anchor}=g$ for the relevant boundary pair. All probabilities
below remain conditional on its first-column assignment; row
relabeling makes the bounds uniform in that assignment.

\begin{lemma}[Finite collision and density bounds]
\label{lem:ub-edge-density-bounds}
For every legal prefix $x$ in the anchored construction, put
$p_x(m)=\Pp((m,m+1)\in E(G_{B_x}))$ and
$D_x=\sum_{m=1}^{n-1}p_x(m)^2$. Then
\begin{equation}
\label{eq:ub-unconditional-prefix-density}
 \E|E(G_{B_x})|\le\tfrac12\mu_x^2N^{e_x},\qquad
 D_x\le\tfrac14N^{4\ell_x+2e_x}.
\end{equation}
If $x$ is a full word, then $D_x\le\mu_x$.
For a nonterminal stable block with $\ell_x\ge2\delta$, Lemma~\ref{lem:ub-anchor-endpoint-atoms} gives, uniformly over the construction,
\begin{equation}
\label{eq:ub-large-prefix-density}
 D_x\le N^{4\ell_x-F_v+2e_x+2\delta+o(1)}.
\end{equation}
\end{lemma}

\begin{proof}
An internal edge can be assigned injectively to a pair of equal rows.
The expected number of prefix-matching row pairs is at most
$\mu_x^2/2$, by the falling-factorial sampling bound in
Lemma~\ref{lem:ub-largest-component-tail}. Their tail matching
probability is at most
$\prod_{j>M_x}\sum_d(q^{(j)}_d)^2=N^{e_x}$, independently of the
prefix columns. This proves the first inequality in
\eqref{eq:ub-unconditional-prefix-density}. The second follows from
$\sum_m p_x(m)^2\le(\sum_m p_x(m))^2$.
For a full word, the internal edge count is $(C_x-1)_+$, so
$D_x\le\sum_m p_x(m)=\E(C_x-1)_+\le\mu_x$.

For a large stable block, $F_v>\delta$ by stability. Let $J_x$
be its set of internal edge offsets relative to $\iota_x$. The
sigma-field $\mathcal U_x$ records the labeled rows and their full
continuations, so it determines $J_x$. Put
\[
 A_x^{\rm atom}=4C_{\rm hyp}
       \sqrt{(2K+\chi^{-1})/g}\,N^{-F_v},\qquad
 a_x(m)=\Pp((m,m+1)\in E(G_{B_x}),\mathcal S_x).
\]
For each offset in $J_x$, the edge event fixes one possible value of
$\iota_x$. Lemma~\ref{lem:ub-anchor-endpoint-atoms}, conditioned on
$\mathcal U_x$, therefore gives
\[
 \sup_m a_x(m)\le A_x^{\rm atom}\E|J_x|,
 \qquad \sum_m a_x(m)\le\E|J_x|.
\]
Writing $p_x(m)=a_x(m)+b_x(m)$, we have
$p_x(m)^2-a_x(m)^2\le2b_x(m)$ and
$\sum_m b_x(m)\le n\Pp(\mathcal S_x^c)$. Consequently
\[
 D_x\le\tfrac14 A_x^{\rm atom}\mu_x^4N^{2e_x}
                       +2n\Pp(\mathcal S_x^c).
\]
By \eqref{eq:ub-anchor-smoothing-failure}, the last term is a
stretched-exponential error, uniformly because $F_v>\delta$.
Also $e_x\ge-T\log k$, $F_v\le1/2$, and $\ell_x\ge2\delta$,
so this error is negligible relative to the displayed power of $N$.
Since $K=O(\log N)$, the atom prefactor is $N^{o(1)}$.
This proves \eqref{eq:ub-large-prefix-density}, with its stated slack.
\end{proof}

We next group stable words by their actual costs, rather than by
frequencies in cells of the simplex. Fix a width $\upsilon>0$.
For an extension from $v$ to $x$, let
\[
 Q_x=\{j:M_v<j\le M_x,\ q^{(j)}_{x[j]}<\chi\},\quad
 S_x=\frac1L\sum_{\substack{M_v<j\le M_x\\j\notin Q_x}}
               \log\frac1{q^{(j)}_{x[j]}},\quad
 T_x=\frac1L\sum_{j\in Q_x}\log\frac1{q^{(j)}_{x[j]}}.
\]
In particular,
\begin{equation}
\label{eq:ub-exact-word-cost-split}
 \ell_x=V_v-S_x-T_x,\qquad S_x,T_x\ge0.
\end{equation}
Fix a large constant $C_0$. Retain boundary words with at most
$C_0L/\log(1/\chi)$ digits strictly below $l_j$ before their left
crossing, or strictly above $r_j$ before their right crossing.
Retain extensions with $|Q_x|\le\alpha_\chi L$, where
$\alpha_\chi=(\log(1/\chi))^{-1/2}$.
A class $\mathcal C$ records the first pile, both boundary words, $M_x$,
the set $Q_x$ and its selected digits, the integer
$\lfloor S_x/\upsilon\rfloor$, and whether the block is proper or
immediate, large or small. Terminal boundary words form singleton
classes. Empty classes are omitted. Put $H_{\mathcal C}=\log_N|\mathcal C|$.

\begin{lemma}[Empirical cost classes]
\label{lem:ub-empirical-cost-classes}
For every fixed $\varepsilon>0$ and $D>0$, first choosing $C_0$ large,
then $\upsilon$ and $\chi$ sufficiently small, gives at most
$N^\varepsilon$ retained classes for all sufficiently large $N$,
uniformly in the arrays. The expected number of discarded rows is at
most $N^{-D}$. Every nonterminal retained class satisfies, for every
$x\in\mathcal C$ and every $3\le\theta\le4$,
\begin{equation}
\label{eq:ub-empirical-class-entropy}
 H_{\mathcal C}\le\theta S_x-\widehat B_\theta(M_v,M_x)+e_H,
 \qquad e_H=4\upsilon+4\log(k)\alpha_\chi.
\end{equation}
\end{lemma}

\begin{proof}
For $q\ge0$, set
$U_q=\sum_{j=0}^{\min(q,K)}\binom Kj k^j$.
A boundary word is determined by its crossing time and digit, its
exceptional positions and digits, and the fixed default digits $l_j$
or $r_j$ elsewhere. A terminal boundary is encoded in the same way.
The number of pairs is therefore at most a constant times
$(K+1)^2U_{q_b}^2$, with
$q_b=\lfloor C_0L/\log(1/\chi)\rfloor$.
The additional extension data cost at most a constant times
$(K+1)U_{q_e}J$, where
\[
 q_e=\lfloor\alpha_\chi L\rfloor,\qquad
 J=1+\left\lfloor T\log(1/\chi)/\upsilon\right\rfloor.
\]
The elementary bound $\binom Kj\le(eK/j)^j$ shows
$L^{-1}\log(U_{q_b}^2U_{q_e})\to0$ as $\chi\downarrow0$, uniformly
for $K\le TL$, after taking the large-$N$ limit. All other factors
are fixed or polynomial in $\log N$. This proves the class bound
with the stated order of choices.

For a fixed labeled row, its digits in different columns are independent.
The probability of using a digit of mass below $\chi$ in each of a
prescribed set of $j$ columns is at most $(k\chi)^j$.
Consequently the expected number of rows failing either truncation is
bounded by a fixed multiple of
\[
 N\sum_{j>q_b}\binom Kj(k\chi)^j
   +N\sum_{j>q_e}\binom Kj(k\chi)^j.
\]
Their logarithms divided by $L$ are at most
$1-C_0+o_\chi(1)$ and
$1-\alpha_\chi\log(1/\chi)+o_\chi(1)$, respectively.
This gives $N^{-D}$ after the ordered choices. No independence between
different rows is needed.

Within a class, restriction of a word to its nonexceptional extension
positions is injective. Under the product of the tilted laws at those
positions, its probability is
\[
 N^{-\theta S_x+B^G_\theta},\qquad
 B^G_\theta=L^{-1}\sum_{\substack{M_v<j\le M_x\\j\notin Q_x}}
                           \psi_{q^{(j)}}(\theta).
\]
If $b$ is the fixed cost-bin index, every word in the class has
$S_x<(b+1)\upsilon$. The sum of their tilted probabilities is at most
one, whence
$|\mathcal C|\le N^{\theta(b+1)\upsilon-B^G_\theta}$.
Now $b\upsilon\le S_x$ and
$0\le\psi_q(\theta)\le4\log k$ for $2\le\theta\le4$.
At most $\alpha_\chi L$ terms were removed. These inequalities give
\eqref{eq:ub-empirical-class-entropy}.
\end{proof}

\begin{lemma}[Empirical exponent bounds]
\label{lem:ub-numerical-lemma}
Under \eqref{eq:ub-empirical-negative-margin}, every large stable
class satisfies
\begin{equation}
\label{eq:ub-empirical-large-exponent}
 H_{\mathcal C}+4\ell_x-F_v+2e_x
 \le-\eta+\delta+q_N/2+e_H.
\end{equation}
Every small proper class satisfies
\begin{equation}
\label{eq:ub-empirical-small-exponent}
 H_{\mathcal C}+4\ell_x+2e_x
 \le-\eta+8\delta+21q_N/2+e_H.
\end{equation}
For an immediate singleton, $\ell_v+e_v\le-\eta$;
for a terminal word, $\ell_x\le-\eta$.
\end{lemma}

\begin{proof}
Set $s=M_v$, $u=M_x$. We have $V_v\le\widehat\beta(s)$.
For every $\theta\in[3,4]$, the exact identity
\[
 4\ell_x-F_v+\theta S_x
 =\theta V_v+(4-\theta)(\ell_x-F_v)
                  +(3-\theta)F_v-\theta T_x
\]
combined with stability, \eqref{eq:ub-anchored-fluctuation-comparison},
and $T_x\ge0$, is at most
$\frac{\theta+3}{2}\widehat\beta(s)+\delta+q_N/2$.
Use \eqref{eq:ub-empirical-class-entropy}, the exact value of $e_x$,
and take the infimum over $\theta$ to obtain
\eqref{eq:ub-empirical-large-exponent}.

For a proper first crossing, the parent has
$\ell_{x^-}\ge F_v+\delta$. The last digit adds at most
$q_N$ to the good cost $S_x$, whether it is good or exceptional.
Thus $S_x\le V_v-F_v+q_N\le\widehat\beta(s)/2+3q_N/2$.
Since $S_x\ge0$, also $\widehat\beta(s)\ge-3q_N$.
It follows that
\[
 \theta S_x
 \le\tfrac{\theta+3}{2}\widehat\beta(s)+21q_N/2
 \qquad(3\le\theta\le4).
\]
Use \eqref{eq:ub-empirical-class-entropy} and $4\ell_x<8\delta$
to prove \eqref{eq:ub-empirical-small-exponent}.

Finally, $\sum_dq_d^2\le\max_dq_d$ gives $\psi_q(2)\ge a_q$.
Therefore
\[
 \ell_v+e_v\le\widehat\beta(s)-\widehat B_2(s,K)
                   \le\widehat\beta(K)\le-\eta.
\]
The terminal assertion is the maximal-word inequality.
\end{proof}

\begin{proof}[Proof of Proposition~\ref{prop:ub-empirical-first-moment}]
We may replace $\eta$ by $\min(\eta,1)$. Choose
$\varepsilon<\eta/100$ and $\upsilon<\eta/1000$, fix $C_0$
sufficiently large for $D=10$, and then choose $\chi$ so small
that Lemma~\ref{lem:ub-empirical-cost-classes} holds with $D=10$ and
$e_H<\eta/50$. Choose $\delta<\eta/1000$ and then take $N$ large,
with all choices fixed before the cut arrays.

At each edge position, the retained relevant proper blocks number at
most $(2kK)$ times the number of retained boundary pairs, by
Lemma~\ref{lem:ub-anchor-relevant-parents}. Immediate and terminal singleton pairs
add at most one block per pair. The boundary-pair count is at most
$N^\varepsilon$. The squared sum of the corresponding edge probabilities
is consequently at most
$N^{2\varepsilon}\sum_x p_x(m)^2$ for all large $N$.
The probability of an omitted nonrelevant proper block is stretched-exponentially
small after a union bound over at most $(K+1)k^{K+1}$ prefixes.
Discarded rows contribute at most $N^{-10}$ to the expected number of
edges, and hence at most this amount to the sum of squared discarded
edge probabilities. Using $(u+v)^2\le2u^2+2v^2$ handles this error.
Thus
\begin{equation}
\label{eq:ub-empirical-first-moment-assembly}
 \E|E(G)\cap E(G')|
 \le N^{-9}+2N^{2\varepsilon}
        \sum_{\mathcal C}\sum_{x\in\mathcal C}D_x.
\end{equation}
There are at most $N^\varepsilon$ classes, and
$|\mathcal C|=N^{H_{\mathcal C}}$ exactly. Equations
\eqref{eq:ub-large-prefix-density} and
\eqref{eq:ub-empirical-large-exponent} bound a large class by
$N^{-\eta+3\delta+q_N/2+e_H+o(1)}$.
Equations \eqref{eq:ub-unconditional-prefix-density} and
\eqref{eq:ub-empirical-small-exponent} bound a small proper class by
$N^{-\eta+8\delta+21q_N/2+e_H}$.
For a small immediate singleton,
$4\ell_v+2e_v=2(\ell_v+e_v)+2\ell_v\le-2\eta+4\delta$;
its class consists of that one word. A terminal singleton contributes
at most $N^{-\eta}$. All powers are at most $-\eta/2$ for large $N$.
The three class/localization losses in
\eqref{eq:ub-empirical-first-moment-assembly} are absorbed by the
remaining margin. For example $c=\eta/4$ works after increasing $N_0$.
All concentration, counting and numerical bounds were uniform in the
arrays and in $n,a,K$, as required.
\end{proof}

\begin{lemma}[Stability of empirical margins]
\label{lem:ub-conditional-preservation}
Let $e,\zeta\ge0$. Suppose $p^{(j)},q^{(j)}\in D_k$ and
$\|p^{(j)}-q^{(j)}\|_\infty\le\zeta$ on the retained columns.
Then, uniformly in interval endpoints and $2\le\theta\le4$,
\begin{equation}
\label{eq:ub-empirical-cost-stability}
 |\widehat A_p-\widehat A_q|\le kT\zeta,\qquad
 |\widehat B_{\theta,p}-\widehat B_{\theta,q}|\le4k^4T\zeta.
\end{equation}
If the original deck satisfies
\eqref{eq:ub-empirical-negative-margin} with margin $\eta$ and a
prefix ending at column $r$ produces a deck of size $m$ with
\[
 \log_Nm\le1-\widehat A_p(0,r)+e,
\]
then its tail $q$ satisfies the same negative-margin inequalities with
margin at least $\eta-\frac72(e+kT\zeta)-12k^4T\zeta$.
Dropping at most $\rho\log N+O(1)$ initial columns without decreasing
population size costs at most an additional
$\frac72\rho\log k+O_k(1/\log N)$ in the margin.
\end{lemma}

\begin{proof}
For probability vectors $p,q$,
\[
 |a_p-a_q|\le k\|p-q\|_\infty,\qquad
 |\psi_p(\theta)-\psi_q(\theta)|\le4k^3\|p-q\|_1.
\]
The first inequality uses $\max p,\max q\ge1/k$.
For the second, the power sums are at least $k^{-3}$ and
$|x^\theta-y^\theta|\le4|x-y|$ on $[0,1]$.
Sum over at most $T\log N$ columns to obtain
\eqref{eq:ub-empirical-cost-stability}. The tail size function is at
most the original size function plus $e+kT\zeta$.
The coefficient of that function is at most $7/2$, and the sum of the
absolute coefficients of the two $B$ terms is three. This gives the
stated loss. Finally $a_p\le\log k$ bounds the discarded initial cost.
No preservation of membership in any partition cell is asserted.
\end{proof}


\subsection{Conditional exploration in the fixed-column model}

Throughout this subsection the population size is the original \(N\),
including after deletion of an initial sequence of columns. We use
a working good threshold $0<g<1/k$ and density $0<\gamma\le1$;
any positive thresholds in the original hypothesis can be decreased
to this range. Write
\(K\le T\log N\), \(p_{j\ell}=n^{(j)}_\ell/N\), and, for a prefix
\(x\), set
\[
 \lambda_x=\prod_{j\le |x|}p_{j,x[j]},\qquad
 \mu_x=N\lambda_x,\qquad
 \mu_*=N\prod_{j\le K}\max_\ell p_{j\ell}.
\]
Zero-probability prefixes are omitted.  Reindex the retained columns by $1,\ldots,K$. Every empirical
endpoint is translated by the same discarded length, and every
normalization remains $\log N$. A tail after a retained prefix $x$
therefore uses $a=|x|$ in \eqref{eq:ub-empirical-costs}.

\begin{lemma}[A short cluster of good columns]
\label{lem:ub-good-cluster}
Fix $g,\gamma,\rho>0$, and put $m_N=\lceil\rho\log N\rceil+1$.
Suppose every interval of $m_N$ consecutive columns contains at least
$\gamma\rho\log N-2$ columns having two piles of size at least $gN$.
For all sufficiently large $N$, every such interval contains three
good columns whose first and last indices differ by at most
$D=\lceil6/\gamma\rceil$.
\end{lemma}

\begin{proof}
Partition the interval into consecutive intervals of length at most
\(D\). If each contained at most two good columns, their total would
be at most \(2\lceil m_N/D\rceil\), which is strictly
smaller than \(\gamma\rho\log N-2\) for large \(N\).
\end{proof}

Definition~\ref{def:good-columns} implies the hypothesis of
Lemma~\ref{lem:ub-good-cluster} for all intervals fitting inside
the finite horizon: discard at most the two endpoint columns when
passing from a closed macroscopic interval to consecutive indices.
We restart at the first good column of such a cluster in the initial
\(\rho\log N+O(1)\) columns. Deleting these earlier shuffles is
permitted by total-variation contraction under convolution with a
probability measure; uniform measure is invariant under that
convolution. This restart deletes at most \(\rho\log N+O(1)\)
columns. Their maximal-pile costs are at most \(\log k\) per column.
Thus the terminal maximal-word exponent increases by at most
\(\rho\log k+o(1)\), and the numerical margins in \eqref{eq:ub-empirical-negative-margin} increase
by at most \((7/2)\rho\log k+o(1)\). The factor \(7/2\) is the
largest coefficient of the ambient exponent in \eqref{eq:ub-empirical-negative-margin}, as quantified by
Lemma~\ref{lem:ub-conditional-preservation}. We choose \(\rho\) small enough to retain a positive
margin. In particular, after decreasing that margin if necessary, fix
\(0<\eta\le1\) such that
\begin{equation}
\label{eq:ub-exploration-maximal-mass}
 \mu_*\le N^{-\eta}.
\end{equation}
The retained first column is \(g\)-good, and two further good columns
occur within its next \(D\) columns.

We specify the two spatial directions and their rank indices.
For a first-column digit \(p\), put
\[
 n_p=n^{(1)}_p,\qquad o_p=1+\sum_{\ell<p}n_\ell.
\]
Its sorted strings occupy the deterministic positions
\(o_p,\ldots,o_p+n_p-1\). Within this pile write them
\(s^p_1\le\cdots\le s^p_{n_p}\). In the next \(D\) coordinates,
choose two digits of probability at least \(g\) at each of the two
good columns and choose a maximal-probability digit elsewhere. The
four resulting tail prefixes include three in increasing order,
denoted \(v^-<v^0<v^+\). Every one has conditional product mass at
least
\begin{equation}
\label{eq:ub-reserve-mass}
 d_0=g^2 k^{-D}>0.
\end{equation}
Here and below these prefixes omit the initial digit \(p\).
A word beginning with \(p\) is called forward if its next \(D\)
digits are at most \(v^0\), and backward if they are at least
\(v^0\). Thus words in the middle cylinder have both directions.
The cylinders with tail prefixes \(v^+\) and \(v^-\) are respectively
right and left reserves, lying strictly beyond the corresponding
directional cutoff. All these choices are deterministic functions of
the cut array.

For a local edge \((u,u+1)\), \(1\le u<n_p\), the global edge is
\((o_p+u-1,o_p+u)\). It is a forward, respectively backward, edge
when its equal endpoint word has that direction. For independent
matrices, a common directional edge must have that direction in
\emph{both} matrices.

\begin{lemma}[Uniform concentration and directional coverage]
\label{lem:ub-exploration-concentration}
Fix \(0<\xi<\eta\). There is an event \(H\) for one matrix,
and constants \(c,A>0\), such that
\begin{equation}
\label{eq:ub-exploration-H-probability}
 \Pp(H^c)\le r_N^{\rm bad}:=A N^A\exp(-cN^{\xi/3}),
\end{equation}
and the following assertions hold on \(H\).
For every initial pile, its complement has, in each later column,
counts within \(N^{2/3}\) of its size times \(p_{j\ell}\).
For every prefix \(x\) with \(\mu_x\ge N^\xi\), writing \(C_x\)
for its row count and \(D_{j\ell}(x)\) for its \(\ell\)-count in
column \(j>|x|\),
\begin{equation}
\label{eq:ub-exploration-prefix-concentration}
 |C_x-\mu_x|\le\mu_x^{2/3},\qquad
 |D_{j\ell}(x)-C_xp_{j\ell}|\le\mu_x^{2/3}.
\end{equation}
For every initial pile with \(n_p\ge N^\xi\), its number of words
strictly before the middle cylinder and its number through that
cylinder are each within \(d_0 n_p/4\) of their expectations.
Consequently, on \(H\times H\), every shared edge in any such pile
has a common direction.
Moreover, the probability that any pile with \(n_p<N^\xi\) contains
an edge is at most \((k/2)N^{\xi-\eta}\).
\end{lemma}

\begin{proof}
We recall the elementary estimates used here. If a set of \(z\)
marked positions meets a uniform set of \(h\) positions in a
population of size \(M\), its intersection count \(Y\) satisfies,
for every \(s>0\),
\[
 \E s^Y\le(1-h/M+(h/M)s)^z.
\]
Indeed its generating function is the normalized \(z\)-th
elementary symmetric polynomial of \(h\) copies of \(s\) and
\(M-h\) copies of \(1\). At fixed sum, replacing any pair by their
average does not decrease this polynomial: the only changing term is
their product times a nonnegative elementary symmetric polynomial
of the other variables. Repeated averaging and continuity give the
displayed bound. This argument applies also when \(s<1\).
Applying this one-step estimate backwards through the independent
columns proves the same binomial generating-function bound for a
prefix count, with its actual mean \(\mu\). Exponential Markov,
with positive and negative exponents, then gives
\begin{equation}
\label{eq:ub-exploration-bernstein}
 \Pp(|Y-\mu|\ge u)
 \le2\exp\left(-\frac{u^2}{2(\mu+u)}\right),\qquad u>0.
\end{equation}
For example, optimize the upper bound
\(\exp(\mu(e^t-1)-t(\mu+u))\), and use
\((1+v)\log(1+v)-v\ge v^2/(2(1+v))\); the lower-tail
calculation uses \(-\log(1-v)\ge v\).

An initial-pile complement is independent of every later column,
so its count is hypergeometric. Taking \(u=N^{2/3}\) in
\eqref{eq:ub-exploration-bernstein} gives failure
\(\exp(-cN^{1/3})\). A prefix count has mean \(\mu_x\), so
\(u=\mu_x^{2/3}\) gives failure \(\exp(-c\mu_x^{1/3})\).
Given its labeled row set, a later-column count is hypergeometric
with mean \(C_xp_{j\ell}\); on the preceding size event,
\(C_x\le2\mu_x\), so the same bound applies. There are at most
\(\sum_{r\le K}k^r=N^{O(1)}\) prefixes and \(kK\) later
column/digit choices. This proves the simultaneous assertions in
\eqref{eq:ub-exploration-prefix-concentration} and their stated
failure bound.

Conditional on the first-column row set, each of the at most \(k^D\)
tail-prefix counts has mean \(n_p\) times its product mass and
satisfies \eqref{eq:ub-exploration-bernstein}. Summing the finitely
many prefixes before, or through, \(v^0\) gives the two guard
concentration assertions, with failure \(\exp(-c n_p)\), hence
within \eqref{eq:ub-exploration-H-probability}. Let their expected
counts be \(n_pL\) and \(n_pU\). Since \(U-L\ge d_0\), the
observed before-count in either matrix is strictly smaller than the
observed through-count in the other. If an edge were before the
middle cylinder in the first matrix and after it in the second, its
local left endpoint \(u\) would satisfy simultaneously
\[
 u+1\le\#\{\text{before in the first matrix}\},\qquad
 u\ge1+\#\{\text{through in the second matrix}\},
\]
a contradiction. Reversing the matrices proves coverage.

Finally, for a full word \(w\) the exact factorial identity and
\((h)_2/(N)_2\le(h/N)^2\) give
\(\E\binom{C_w}{2}\le N^2\lambda_w^2/2\). Therefore
\[
 \E\sum_{w[1]=p}\binom{C_w}{2}
 \le\frac{N^2}{2}\max_w\lambda_w\sum_{w[1]=p}\lambda_w
 =\frac{n_p\mu_*}{2}.
\]
Every graph edge implies such a collision. Sum over the at most
\(k\) small piles and use \eqref{eq:ub-exploration-maximal-mass}.
\end{proof}

We use the natural sorted-word histories within each initial pile.
For a nonempty pile the forward one-matrix history is
\(h_i=(s^p_1,\ldots,s^p_i)\), \(0\le i\le n_p\), with
\(h_0\) empty. Backward histories reveal the same strings in the
opposite spatial order. They do not reverse the columns. Each pile
and direction uses its own filtration: no later-column word outside
that initial pile is revealed. In particular the history does not
include previously explored initial piles. The outside-pile rows
remain available as the reservoir in every application below.
Put
\(\varepsilon_N=\sqrt{r_N^{\rm bad}}\), and, for either direction and
\(1\le i<n_p\), define the history-measurable event
\begin{equation}
\label{eq:ub-natural-history-good}
 F_i=\{\Pp(H^c\mid h_i)\le\varepsilon_N\}.
\end{equation}
Conditional Markov gives \(\Pp(F_i^c)\le\varepsilon_N\).
All conditional assertions below concern positive-probability
histories. Intersecting these events for both directions and all
piles loses at most \(2N\varepsilon_N\).

\begin{lemma}[Conditional counts from an actual reservoir]
\label{lem:ub-reservoir-counts}
Fix \(0<\tau<g\) and put
\(\alpha=\log(1/g)/\log(1/\tau)\). Assume \(\alpha<1/2\).
For all sufficiently large \(N\), uniformly in the cuts and positive
histories, the following holds. Consider a forward history \(h_i\)
on \(F_i\), whose last word \(s\) is forward, and let \(t\) be
the number of observed terminal copies of \(s\). For every prefix
cylinder \(B_x\) lying strictly above \(s\) and within the forward
cutoff,
\begin{align}
 \E[C_x\mid h_i]
 &\le C_1N\lambda_x^{1-\alpha}+N\varepsilon_N,
 \label{eq:ub-conditional-cylinder-count}\\
 \E[(C_s-t)_+\mid h_i]
 &\le C_1N\lambda_s^{1-\alpha}+N\varepsilon_N,
 \label{eq:ub-conditional-extra-copies}\\
 \E\left[\sum_{w\supseteq x}\binom{C_w}{2}\mathrel{\Big|}h_i\right]
 &\le C_2N^2\lambda_x^{2(1-\alpha)}
       \prod_{j>|x|}\sum_\ell p_{j\ell}^{2(1-\alpha)}
       +N^2\varepsilon_N.
 \label{eq:ub-conditional-tail-collisions}
\end{align}
The constants depend only on \(g,\tau,k,D,T\). The reversed spatial
statements hold for backward histories.
\end{lemma}

\begin{proof}
Fix the labeled first-column pile \(P\), of size \(n=n_p\).
This is without loss of generality: simultaneous relabeling of the
rows shows that the joint law of the sorted history and \(H\) is
independent of the particular labeled set \(P\).
This conditioning leaves all later columns independent and uniform
with their prescribed counts. Its complement has size
\(c=N-n\ge gN\): the first column has two piles of size at least
\(gN\), so no pile exceeds \((1-g)N\).

For a forward history, let \(h_u\) be the observed multiplicity of
each full word \(u<s\) beginning with \(p\). The history event is exactly
\begin{equation}
\label{eq:ub-exact-history-constraint}
 C_u=h_u\quad(u[1]=p,\ u<s),\qquad C_s\ge t.
\end{equation}
Choose \(m=1\) or \(2\) marked labeled rows in \(P\), and fix
all the other \(n-m\) rows, calling their configuration \(a\).
Before imposing \eqref{eq:ub-exact-history-constraint}, the marked
rows together with the outside reservoir are independent uniform
fixed-count columns on \(c+m\) rows. If \(R_{j\ell}\) denotes
the remaining column count, the exact probability that all marked
rows begin with \(x\) is
\begin{equation}
\label{eq:ub-reservoir-exact-prefix-weight}
 L_a(x^m)=\prod_{2\le j\le |x|}
                  \frac{(R_{j,x[j]})_m}{(c+m)_m}.
\end{equation}
The conditional law given \(a,h_i\) is this law restricted to the
marked tuples compatible with \eqref{eq:ub-exact-history-constraint},
and divided by the weight of that allowed set. This follows by
counting completions in each column: a prescribed ordered tuple
leaves \((c)!/\prod_\ell(R_{j\ell}-r_{j\ell})!\) completions,
where \(r_{j\ell}\) is its prescribed \(\ell\)-multiplicity.
The resulting normalized ratio is the ordinary without-replacement
probability. Thus the lexicographic conditioning has not been
discarded.

If all marked words lie strictly above \(s\), their removal leaves
\(a\) satisfying the entire history constraint. If \(a\) does
not satisfy it, that target event is impossible. If it does, every
marked tuple in the right reserve cylinder is allowed. On the event
\[
 G_m(a)=\{|R_{j\ell}-cp_{j\ell}|\le N^{2/3}+m
                           \text{ for all }j,\ell\},
\]
the exact reserve probability, using only its \(D\) prescribed
coordinates, is at least \(d_0^m/2\) for large \(N\).
All subsequent columns sum out. Notice also that \(H\) implies
\(G_m(a)\), for every marked tuple, since its counts differ from
the outside-reservoir counts by at most \(m\).

For \(p_{j\ell}\ge\tau\), on \(G_m(a)\),
\[
 \frac{R_{j\ell}}{c+m}
 \le p_{j\ell}\bigl(1+O_{g,\tau}(N^{-1/3})\bigr).
\]
The product of these errors over \(K\le T\log N\) columns is
bounded by \(2\) for large \(N\). For \(p_{j\ell}<\tau\),
the deterministic inequalities \(R_{j\ell}\le Np_{j\ell}\) and
\(c+m\ge gN\) give \(R_{j\ell}/(c+m)\le p_{j\ell}/g\).
Write
\(q_x=\prod_{2\le j\le |x|}p_{j,x[j]}=\lambda_x/(n/N)\).
If \(r\) of these selected digits have probability below \(\tau\),
then \(q_x\le\tau^r\), whence \(g^{-r}\le q_x^{-\alpha}\).
Using \((a)_m/(M)_m\le(a/M)^m\) in
\eqref{eq:ub-reservoir-exact-prefix-weight}, we obtain
\[
 L_a(x^m)\le2^m q_x^{m(1-\alpha)}.
\]
For a zero-mass prefix the actual probability is zero. Dividing by
the reserve weight proves, on \(G_m(a)\), the bound
\(2^{m+1}d_0^{-m}q_x^{m(1-\alpha)}\) for the conditional target
probability. Its failure set, conditionally on \(h_i\in F_i\),
has probability at most \(\varepsilon_N\).

Sum over ordered marked tuples. The initial-pile normalization
cancels in the correct direction:
\[
 (n)_m q_x^{m(1-\alpha)}
 \le N^m(n/N)^{m\alpha}\lambda_x^{m(1-\alpha)}
 \le N^m\lambda_x^{m(1-\alpha)}.
\]
This proves \eqref{eq:ub-conditional-cylinder-count}. For the
current word, use instead the removable-copy inequality
\[
 (C_s-t)_+
 \le\sum_{r\in P}\1_{\{\text{row }r=s,\ C_s-1\ge t\}}.
\]
On each summand, removing the marked row leaves the history
constraint satisfied, so the same reserve comparison applies.
This proves \eqref{eq:ub-conditional-extra-copies} without
treating any of the \(t\) forced observed copies as a fresh copy.

For \eqref{eq:ub-conditional-tail-collisions}, use two marked rows
and sum their equal-word probabilities over completions of \(x\).
The product structure of \(q_w\) gives the displayed tail factor.
The bad configurations are charged only once: for a given marked
pair the event that its words agree is one event, and its probability
on \(G_2(a)^c\) is at most \(\varepsilon_N\). Summing over
unordered marked pairs costs at most \(N^2\varepsilon_N/2\),
not that error multiplied by the number of full words. Reversing
the lexicographic inequalities and using the left reserve proves
the backward statements.
\end{proof}

\begin{lemma}[The conditional law inside future cylinders]
\label{lem:ub-exact-conditional-block-law}
Fix a positive forward history \(h_i\), and partition a strict
future lexicographic interval into disjoint prefix cylinders.
Condition additionally on their absolute rank intervals, their
labeled row sets, all their induced tail-column counts, and the
current word's total multiplicity. Inside any one cylinder \(B_x\),
its columns after \(|x|\) have the uniform independent fixed-count
law with those induced counts. Consequently its sorted internal
graph has precisely the fixed-count graph law on \(C_x\) vertices.
The same conclusion holds when the labeled row sets are averaged
out of this conditioning.
\end{lemma}

\begin{proof}
For each \(j>|x|\), permute the entries in column \(j\) within
the labeled row set of \(B_x\), independently for different
columns. These operations preserve the global column counts and
the prefix \(x\), and hence preserve membership in \(B_x\).
They preserve every other cylinder and its conditioned counts,
the absolute rank intervals, and the current word's multiplicity.
They also preserve the observed history: all words in \(B_x\)
are strictly larger than its last word. The original finite matrix
measure is uniform, so its conditional measure is invariant under
these operations. The operations act transitively on the product
of the sets of tail columns with the specified counts. Thus this
conditional law is uniform on that product, which is the asserted
independent-column law. All words in the block share \(x\), so
sorting them sorts precisely their tails. Finally the resulting
internal graph law depends only on the induced counts and the
interval, not on the labeled row set being averaged out.
\end{proof}

\begin{proposition}[Conditional first moment in every natural history]
\label{lem:ub-conditional-first-moment}
Fix $k\ge2$, $T>0$, $0<\eta\le1$, $0<g<1/k$ and
$0<\gamma\le1$, and put $\xi=\eta/10$ and $D=\lceil6/\gamma\rceil$.
Choose a fixed $\rho>0$ sufficiently small in terms of these parameters,
as in \eqref{eq:ub-exploration-exponent-budget}. Assume the retained
array satisfies \eqref{eq:ub-empirical-negative-margin} with $n=N$,
$a=0$ and margin $2\eta$, and the consecutive-window hypothesis of
Lemma~\ref{lem:ub-good-cluster}. Its first column must be $g$-good,
with two further $g$-good columns among the next $D$ columns.
Use the events $H,F_i$ constructed with this $\xi$.
There exist $\zeta,C>0$ and $N_0$, depending only on the fixed
parameters, such that for every $N\ge N_0$, every such cut array,
and all positive-probability forward or backward histories in any
pile with $n_p\ge N^\xi$,
\begin{equation}
\label{eq:ub-actual-conditional-remaining}
 \E[\text{unrevealed common directional edges}\mid h_i,h'_i]
 \le C N^{-\zeta}
 \quad\text{on }F_i\cap F'_i,\qquad 1\le i<n_p.
\end{equation}
For \(i=0\), the same bound holds with the two empty histories;
at \(i=n_p\) the remaining count is zero.
\end{proposition}

\begin{proof}
Set \(\xi=\eta/10\). Choose \(\tau\) so small that
\begin{equation}
\label{eq:ub-exploration-alpha-budget}
 \alpha\le\min\left\{\frac{\xi}{4(1+\xi)},
                         \frac{\eta}{40(1+\eta)}\right\}.
\end{equation}
Apply Proposition~\ref{prop:ub-empirical-first-moment} with a
margin $\eta$ and a first-column good threshold \(g/2\).
Denote its uniform decay exponent by \(\delta_F>0\); take the
minimum with its whole-deck decay exponent if necessary. Put
\begin{equation}
\label{eq:ub-exploration-exponent-budget}
 a=\min\{3\xi/4,\,3\eta/5,\,\delta_F\},\qquad
 \rho\log k\le a/4,\qquad \zeta=a/2.
\end{equation}
The parameters are fixed in this order: a smaller retained margin,
\(g,\gamma,k,T\), then \(\xi\), the uniform empirical
first-moment exponent \(\delta_F\), then \(\tau\) and a
sufficiently small \(\rho\), and finally \(N_0\). The initial
restart and all deterministic profile errors are also chosen within
the smaller-margin budget. No choice depends on the random history
or on the varying cut array.

Consider one forward history \(h_i\in F_i\). If its current word
is beyond the forward cutoff, no forward edge remains. Otherwise
the strict interval from its current word through that cutoff is
a disjoint union of at most \(2kK\) prefix cylinders. This follows
by following the two endpoint branches of the prefix tree: at
each depth at most \(2(k-1)\) complete sibling cylinders are
included. When at least $m_N=\lceil\rho\log N\rceil+1$ columns remain
after a cylinder prefix, refine it to the index immediately before
the first good column of a cluster supplied by
Lemma~\ref{lem:ub-good-cluster}. This adds at most $m_N$ digits. Every refined tail now has a good
first column and two more good columns within \(D\) positions.
The final disjoint family \(\mathcal B(h_i)\) has size
\begin{equation}
\label{eq:ub-exploration-block-count}
 |\mathcal B(h_i)|\le C_0(\log N)N^{\rho\log k}.
\end{equation}
Prefixes with fewer remaining columns need no such refinement.
Indeed their maximal tail mass \(q_*\) is at least
$k^{-m_N}$, while
\(\mu_xq_*\le\mu_*\). Thus
\(\mu_x\le O(N^{-\eta+\rho\log k})<N^{-\xi}\), by
\eqref{eq:ub-exploration-exponent-budget}, for large \(N\).

Let \(q_e(h_i)\) be the conditional probability that position
\(e\) is an unrevealed forward edge. Refine the history by the
information in Lemma~\ref{lem:ub-exact-conditional-block-law} for
this disjoint family, and call that information \(Z\). This is auxiliary information
used only in the following conditional expectation; it is never
added to the natural exploration filtration. Conditional
Jensen gives
\begin{equation}
\label{eq:ub-one-history-jensen}
 \sum_e q_e(h_i)^2
 \le\E\left[\sum_e q_e(h_i,Z)^2\mathrel{\Big|}h_i\right].
\end{equation}
Given \(Z\), every unrevealed edge lies either in the remaining
current-word run or inside one cylinder. An equal-word edge cannot
cross two distinct cylinders. The current run contributes at most
\(\E[(C_s-t)_+\mid h_i]\); this includes the possible edge
from the last revealed copy to the first unrevealed copy. The
right side of \eqref{eq:ub-one-history-jensen} is therefore at
most that term plus the sum of the conditional self-intersection
expectations inside the cylinders. We classify the \emph{final}
cylinders, after refinement, in three regimes.

If \(\mu_x\le N^{-\xi}\), self-intersection is at most the
single graph's vertex count. By
\eqref{eq:ub-conditional-cylinder-count}, its conditional expectation
is at most
\[
 C N^{-\xi+\alpha(1+\xi)}+N\varepsilon_N
 \le C N^{-3\xi/4}+N\varepsilon_N.
\]
If \(N^{-\xi}<\mu_x<N^\xi\), put
\(q_* =\prod_{j>|x|}\max_\ell p_{j\ell}\). Since
\(\lambda_xq_*\le\max_w\lambda_w\),
\[
 q_*\le N^{-\eta+\xi},\qquad
 \prod_{j>|x|}\sum_\ell p_{j\ell}^{2(1-\alpha)}
 \le q_*^{1-2\alpha}.
\]
Self-intersection is at most the number of collision pairs in the
single block. Equation~\eqref{eq:ub-conditional-tail-collisions}
bounds its expectation by
\[
 C N^{-\eta+3\xi+2\alpha(1+\eta-2\xi)}+N^2\varepsilon_N
 \le C N^{-3\eta/5}+N^2\varepsilon_N.
\]
Thus the near-unit regime is controlled by its remaining collisions,
not by total block size.

Finally suppose \(\mu_x\ge N^\xi\). On the part of \(Z\)
satisfying \eqref{eq:ub-exploration-prefix-concentration}, its actual
size \(m=C_x\) obeys
\[
 \tfrac12\mu_x\le m\le2\mu_x,\qquad
 \max_{j,\ell}\left|\frac{D_{j\ell}(x)}m-p_{j\ell}\right|
 \le2N^{-\xi/3}.
\]
In particular \(m\ge N^{\xi/2}\) for large \(N\). The three
anchor columns remain \(g/2\)-good. Lemma~\ref{lem:ub-conditional-preservation} controls the change
in the empirical costs by \(O(N^{-\xi/3})\). Also
\(\log_Nm=\log_N\mu_x+o(1)\), and
\(\mu_x\le N\prod_{j\le|x|}\max_\ell p_{j\ell}\).
Consequently the actual induced array and the actual exponent
\(\log_Nm\) satisfy \eqref{eq:ub-empirical-negative-margin} with margin
at least $\eta$ for all sufficiently large $N$: explicitly, the
loss is at most $\frac72(e_N+kTz_N)+12k^4Tz_N$, where
$e_N=\log2/\log N$ and $z_N=2N^{-\xi/3}$. This is a
consequence of the whole-deck margin and
Lemma~\ref{lem:ub-conditional-preservation}, not an additional
recursive hypothesis.
Here the time grid and every exponent still use the original
\(\log N\); no replacement by \(\log m\) is made.
Lemma~\ref{lem:ub-exact-conditional-block-law} and
Proposition~\ref{prop:ub-empirical-first-moment} therefore bound
its conditional self-intersection by \(C_FN^{-\delta_F}\).
The exceptional \(Z\)-set has conditional probability at most
\(\varepsilon_N\), because \(H\) implies all the required
concentration assertions and \(h_i\in F_i\). Its contribution
is at most \(N\varepsilon_N\).

The current-run term from
\eqref{eq:ub-conditional-extra-copies} is at most
\(C N^{-\eta+\alpha(1+\eta)}+N\varepsilon_N\), by
\eqref{eq:ub-exploration-maximal-mass}. Combining all three regimes
with \eqref{eq:ub-exploration-block-count} gives
\[
 \sum_e q_e(h_i)^2
 \le C N^{-\eta+\alpha(1+\eta)}
 +C(\log N)N^{\rho\log k}
   \bigl(N^{-3\xi/4}+N^{-3\eta/5}+N^{-\delta_F}
                            +N^2\varepsilon_N\bigr)
 \le C'N^{-a/2}.
\]
The stretched-exponential error absorbs every displayed polynomial,
and \(\log N\le N^{a/4}\) for large \(N\).

The two matrices remain independent after their respective histories
are fixed. Hence the actual expected remaining common count is
\[
 \sum_e q_e(h_i)q_e(h'_i)
 \le\frac12\left(\sum_e q_e(h_i)^2+
                         \sum_e q_e(h'_i)^2\right)
 \le C'N^{-\zeta}.
\]
Both histories receive their own uniform bound; symmetry is not
used to select one. The backward argument reverses the spatial
order and uses the left reserve. For the empty histories apply
Proposition~\ref{prop:ub-empirical-first-moment} to the retained
whole deck, which has the same good anchor and buffered margin;
the count in any one direction is bounded by the total shared-edge
count. After the last vertex is revealed no edge remains.
\end{proof}

We finish by constructing a single-matrix typical event and the adapted
exploration processes.
Let \(\mathcal R\) be the event that \(H\) and all the forward
and backward history-good events hold, and that every initial pile
with \(n_p<N^\xi\) is edge-free. Then
\begin{equation}
\label{eq:ub-exploration-regularity-probability}
 \Pp(\mathcal R^c)
 \le r_N^{\rm bad}+2N\varepsilon_N+(k/2)N^{\xi-\eta}=o(1).
\end{equation}
For each large pile and each direction, reveal the two sorted lists
in that direction simultaneously. At time \(i\), \(0\le i\le n_p\),
let \(X_i\) count the common directional edges whose two endpoints
are among the first \(i\) revealed vertices. In particular
\(X_0=X_1=0\), and every increment is zero or one. Use the natural
pair history, take its good event to be \(F_i\cap F'_i\) for
\(1\le i<n_p\), and take the initial and terminal good events
to be the whole space. For a small or empty pile use the identically
zero process. There are at most \(2k\) processes.
Their filtrations are separate; they are not joined into one
global-deck filtration. The event \(\mathcal R\) intersects their
good events only for the final truncation and is never inserted
into any exploration history.
Proposition~\ref{lem:ub-conditional-first-moment} verifies every
conditional remaining-count inequality, including time zero.
On \(\mathcal R\times\mathcal R\),
Lemma~\ref{lem:ub-exploration-concentration} and edge-freeness of
the small piles give
\[
 |E(G)\cap E(G')|\le\sum_{p,\,\pm}X^{p,\pm}_{n_p}.
\]
Thus the conditional estimates, directional cover, and typical-set
probability required for the finite exploration criterion all follow
from the fixed-column model and the empirical first-moment theorem.

\begin{lemma}[Adapted exponential excess]
\label{lem:ub-portenko-exponential-moment}
Let $(\mathcal F_i)_{i=0}^m$ be a filtration on a finite probability
space. Let $X_i$ be nonnegative integer-valued, $\mathcal F_i$-measurable,
with $X_0=0$ and $X_j-X_{j-1}\in\{0,1\}$.
Let $B_i\in\mathcal F_i$ for $0\le i\le m$, and suppose that
$\gamma\ge0$ and
\begin{equation}
\label{eq:ub-adapted-tail-premise}
 \E[X_m-X_i\mid\mathcal F_i]\le\gamma
 \quad\hbox{on }B_i,\qquad 0\le i<m.
\end{equation}
For any event $C\subseteq\bigcap_{i=0}^mB_i$ and any $t\ge0$ such
that $(e^t-1)\gamma<1$,
\begin{equation}
\label{eq:ub-adapted-excess-bound}
 \E[(e^{tX_m}-1)\1_C]
 \le\frac{(e^t-1)\gamma}{1-(e^t-1)\gamma}.
\end{equation}
All conditional statements and event inclusions may be interpreted
almost surely. In particular the premise includes the initial history.
\end{lemma}

\begin{proof}
Set $D_i=\bigcap_{h=0}^iB_h$ and
\[
 Z_i=\sum_{j=1}^i(X_j-X_{j-1})\1_{D_j},\qquad Z_0=0.
\]
This is an adapted process with increments in $\{0,1\}$. Since
$D_j\subseteq D_i$ for $j>i$,
\[
 0\le Z_m-Z_i\le\1_{D_i}(X_m-X_i).
\]
Taking the conditional expectation and using $D_i\subseteq B_i$ gives
\begin{equation}
\label{eq:ub-stopped-tail-bound}
 \E[Z_m-Z_i\mid\mathcal F_i]\le\gamma
 \qquad(0\le i\le m).
\end{equation}
For every integer $r\ge0$, we claim
\begin{equation}
\label{eq:ub-stopped-binomial-moment}
 \E\left[\binom{Z_m-Z_i}{r}\mathrel{\Big|}\mathcal F_i\right]
 \le\gamma^r.
\end{equation}
The assertion is immediate for $r=0$. For $r\ge1$, the unit increments
give the pathwise identity
\[
 \binom{Z_m-Z_i}{r}
 =\sum_{j=i+1}^m(Z_j-Z_{j-1})\binom{Z_m-Z_j}{r-1}.
\]
Condition each summand first on $\mathcal F_j$. Its first factor is
$\mathcal F_j$-measurable, so induction bounds the conditional
expectation of its second factor by $\gamma^{r-1}$. Now condition on
$\mathcal F_i$ and use \eqref{eq:ub-stopped-tail-bound}. This proves
\eqref{eq:ub-stopped-binomial-moment}.

Put $q=e^t-1\ge0$. The binomial theorem and
\eqref{eq:ub-stopped-binomial-moment}, with $i=0$, imply
\[
 \E(e^{tZ_m}-1)
 =\sum_{r=1}^m q^r\E\binom{Z_m}{r}
 \le\sum_{r=1}^{\infty}(q\gamma)^r
 =\frac{q\gamma}{1-q\gamma}.
\]
On $C$, all $D_j$ hold and $Z_m=X_m$. Since the exponential excess is
nonnegative, $(e^{tX_m}-1)\1_C\le e^{tZ_m}-1$, proving
\eqref{eq:ub-adapted-excess-bound}. No history was enlarged by $C$.
\end{proof}

\begin{proposition}[Finite exploration criterion for total variation]
\label{lem:ub-shared-edge-exponential-moment}
Let $B\ge2$ and $d\ge1$ be integers, and let $S_1$ be a set of single
shuffle matrices whose graphs have components of size at most $B$.
On the product space of the independent matrices $S,S'$, suppose there
are $d$ explorations satisfying the hypotheses of
Lemma~\ref{lem:ub-portenko-exponential-moment}, with processes
$X^{(a)}_i$, terminal times $m_a$, adapted events $B^{(a)}_i$, and a
common bound $\gamma$, for $1\le a\le d$. Suppose further that
\begin{equation}
\label{eq:ub-direction-cover-premise}
 \begin{aligned}
 |E(G)\cap E(G')|&\le\sum_{a=1}^dX^{(a)}_{m_a}
 \quad\hbox{on }\{S\in S_1,S'\in S_1\},\\
 \{S\in S_1,S'\in S_1\}&\subseteq
                \bigcap_{a=1}^d\bigcap_{i=0}^{m_a}B^{(a)}_i.
 \end{aligned}
\end{equation}
Define $q=d_B^d-1$, with $d_B$ from
\eqref{eq:ub-component-constant}. If $q\gamma<1$, then
\begin{equation}
\label{eq:ub-exploration-excess}
 \E\left[(d_B^{|E(G)\cap E(G')|}-1)
       \1_{\{S\in S_1,S'\in S_1\}}\right]
 \le\frac{q\gamma}{1-q\gamma},
\end{equation}
and
\begin{equation}
\label{eq:ub-finite-exploration-tv}
 \dTV(P(n,K),\nuUnif)
 \le\frac12\sqrt{\frac{q\gamma}{1-q\gamma}}+\Pp(S\notin S_1).
\end{equation}
\end{proposition}

\begin{proof}
For nonnegative $x_1,\ldots,x_d$, convexity gives
\[
 d_B^{\sum_a x_a}-1
 \le\frac1d\sum_{a=1}^d\bigl(e^{d\log(d_B)x_a}-1\bigr).
\]
Use \eqref{eq:ub-direction-cover-premise}, multiply by the product-event
indicator, and apply \eqref{eq:ub-adapted-excess-bound} to every term,
with $t=d\log d_B$. This proves \eqref{eq:ub-exploration-excess}.
Lemma~\ref{lem:ub-component-factorial} and
\eqref{eq:ub-truncated-chi-square} prove
\eqref{eq:ub-finite-exploration-tv}.
\end{proof}

\begin{corollary}[Bounded-component truncation]
\label{lem:ub-regularity-sparsity}
Suppose $\mu_*\le N^{-\eta}$ for fixed $\eta>0$.
Let $\mathcal R$ be any set of single matrices, and put
\[
 S_1=\{S\in\mathcal R:\text{ every component of }G(S)
                            \text{ has at most }B\text{ vertices}\}.
\]
For every fixed $A>0$, choosing $B\ge2$ with $\eta B\ge A+1$ gives
\begin{equation}
\label{eq:ub-single-matrix-truncation}
 \Pp(S\notin S_1)\le\Pp(S\notin\mathcal R)
                         +\frac{N^{-A}}{(B+1)!}.
\end{equation}
If the explorations in
Proposition~\ref{lem:ub-shared-edge-exponential-moment} are constructed
for this $S_1$, with fixed $d$ and $\gamma=O(N^{-\delta})$, and
$\Pp(S\notin\mathcal R)=o(1)$, then
$\dTV(P(n,K),\nuUnif)\to0$.
\end{corollary}

\begin{proof}
The union bound and Lemma~\ref{lem:ub-largest-component-tail} give
\eqref{eq:ub-single-matrix-truncation}. The integers $B,d$ are fixed,
so $q=d_B^d-1$ is independent of $N$ and $q\gamma\to0$.
Equation \eqref{eq:ub-finite-exploration-tv} proves the last assertion.
\end{proof}

\begin{theorem}[Empirical upper bound]
\label{thm:ub-empirical-upper}
Fix $k\ge2$ and $T,\eta_0,g_0,\gamma_0>0$.
There is $\rho_0>0$ such that, for every fixed $0<\rho\le\rho_0$,
there are $c>0$ and $N_0$ with the following property.
For every $N\ge N_0$ and deterministic array of $K\le T\log N$
fixed-count columns on $N$ labels, suppose
\eqref{eq:ub-empirical-negative-margin} holds with $n=N$, $a=0$
and margin $\eta_0$. Suppose also that the local
$(g_0,\gamma_0,\rho)$-good-shuffle condition of
Definition~\ref{def:good-columns} holds on $[0,K/\log N]$.
Then
\begin{equation}
\label{eq:ub-empirical-tv-conclusion}
 \dTV(P(n,K),\nuUnif)\le N^{-c}.
\end{equation}
The constants are uniform over the arrays and $K$ in this range.
No macroscopic time profile is needed.
\end{theorem}

\begin{proof}
Put $g=\min\{g_0,1/(2k)\}$,
$\gamma=\min\{\gamma_0,1\}$ and
$\eta=\min\{\eta_0/4,1\}$. Shrinking the good threshold and
density preserves the assumed local condition.
Take $\xi=\eta/10$ and obtain the exponent $\delta_F>0$ in
Proposition~\ref{prop:ub-empirical-first-moment} for margin $\eta$
and good threshold $g/2$. Define
$a=\min\{3\xi/4,3\eta/5,\delta_F\}>0$ and choose $\rho_0$ so that
\begin{equation}
\label{eq:ub-completion-rho-budget}
 \rho_0\log k\le\min\{\eta/7,a/4,1/4\}.
\end{equation}
All these choices precede the cut array and $N$.

The terminal margin and $a_p\le\log k$ imply
$K\ge(1+\eta_0)\log N/\log k$. Thus the initial window of
$m_N=\lceil\rho\log N\rceil+1$ columns fits inside the horizon
for large $N$. Lemma~\ref{lem:ub-good-cluster} supplies the anchor
and two further good columns. Discard the columns preceding this
anchor. Lemma~\ref{lem:ub-conditional-preservation} and
\eqref{eq:ub-completion-rho-budget} show that the retained array
has margin at least $2\eta$ for large $N$. Its maximum full-word
mean is consequently at most $N^{-2\eta}$, and in particular it
satisfies \eqref{eq:ub-exploration-maximal-mass}.
Every retained integer window corresponds to a window of the
original array, so its local good-column condition is inherited.

The constructions leading to
Proposition~\ref{lem:ub-conditional-first-moment} apply, with the
parameters in \eqref{eq:ub-exploration-alpha-budget} and
\eqref{eq:ub-exploration-exponent-budget}. They give $2k$ separate
unit-increment explorations, a common conditional remainder bound
$\gamma_N=O(N^{-a/2})$, and the single-matrix event $\mathcal R$
with the failure estimate
\eqref{eq:ub-exploration-regularity-probability}.
Choose a fixed integer $B\ge2$ with $\eta B\ge2$, and restrict
$\mathcal R$ further to graphs with components of size at most $B$.
Lemma~\ref{lem:ub-largest-component-tail} bounds this additional
failure probability by $N^{-1}/(B+1)!$.

Apply Proposition~\ref{lem:ub-shared-edge-exponential-moment} with
$d=2k$. Its constant $q=d_B^{2k}-1$ is fixed, so $q\gamma_N<1$
for large $N$. Equation~\eqref{eq:ub-finite-exploration-tv} gives
\[
 \dTV(P_{\rm ret},\nuUnif)
 \le O(N^{-a/4})+O(N^{\xi-\eta})+O(N^{-1}),
\]
where $P_{\rm ret}$ denotes the inverse-shuffle law of the retained
array; the remaining errors are stretched exponential.
Convolution with the deleted independent shuffles cannot increase
distance to the uniform law. This proves
\eqref{eq:ub-empirical-tv-conclusion}, for instance with any
$0<c<\min\{a/4,\eta-\xi,1\}$ after increasing $N_0$.
\end{proof}

\begin{proposition}[Random time-profile upper bound]
\label{cor:random-time-profile-upper-bound}
Fix the profile, negative reduced margin, partition and ordered
parameters of Theorem~\ref{thm:deterministic-upper-bound}.  Let $X$
be randomized-almost-profile-like with those parameters.  Suppose
also that its realized cuts satisfy the local-good hypothesis with
probability $1-o(1)$.  Then
\[
 d_{\TV}(P(X,K_N),\nuUnif)\longrightarrow0.
\]
\end{proposition}
\begin{proof}
Let $\mathcal A_N$ be the intersection of the almost-profile and
local-good events.  Its complement has probability $o(1)$ by the
union bound; the two events need not be independent.  Conditional
on $X$, the law is exactly the deterministic law to which the
uniform theorem applies.  Convexity of total variation gives
\[
 d_{\TV}(P(X,K_N),\nuUnif)
 \le\Pp(\mathcal A_N^c)+
       \E\!\left[\mathbf1_{\mathcal A_N}
                    d_{\TV}(P(n(X),K_N),\nuUnif)\right].
\]
Here $n(X)$ denotes the realized deterministic cut array.  The first
term tends to zero by assumption, and the second is at most a
deterministic $b_N\to0$ supplied by the uniform deterministic theorem.
\end{proof}

The same argument applies directly to
Theorem~\ref{thm:ub-empirical-upper}: if its empirical-margin and
local-good hypotheses hold simultaneously with probability $1-o(1)$,
then the mixture distance is at most that failure probability plus
$N^{-c}$.  This version requires no macroscopic profile.


\section{Derivation from empirical upper margins}
\label{sec:buffered-upper-bound-derivation}

Only the two whole-deck profile margins are required.  The deterministic
transfer of bounded scalar costs is
Lemma~\ref{lem:lower-profile-transfer}, whose proof uses regular
partitions and the almost-profile condition and is independent of any
upper or lower probability estimate.

\begin{proof}[Proof of Theorem~\ref{thm:deterministic-upper-bound}]
Put $\eta_e=\eta/2$ and $e=\eta/100$.
Write $M=3\log k$.  Choose $\chi_0>0$ so small that
$2TC_k\chi_0\le e/4$.  Fix the regular partition in the theorem;
let $m$ be its number of cells.  The explicit quadrature proof of
Lemma~\ref{lem:lower-profile-transfer} applies to this partition,
with cell oscillation at most $C_k\chi$.
Choose $\rho_0$ below the empirical threshold
$\rho_0^{\rm emp}(k,T,\eta_e,g,\gamma)$ from
Theorem~\ref{thm:ub-empirical-upper}, and also so that
$4M\rho_0\le e/4$.  For each fixed $0<\rho\le\rho_0$, choose
$\varphi_0>0$ with $Mm\varphi_0T\le e/4$.
After $\varphi$ is fixed, the remaining explicit terms of order
$1/\log N$ are at most $e/4$ for all sufficiently large $N$.
Consequently, writing $S_N=K_N/\log N$, for every integer
$0\le i\le j\le K_N$ and every $2\le\theta\le4$,
\[
 |\widehat A(i,j)-A(i/\log N,j/\log N)|\le e,
 \qquad
 |\widehat B_\theta(i,j)-B_\theta(i/\log N,j/\log N)|\le e.
\]
Here the empirical costs in \eqref{eq:ub-empirical-costs} have
$n=N$ and $a=0$.
The omitted terminal interval has length less than $1/\log N$.
Thus, at every integer grid pair $0\le i\le j\le K_N$,
\begin{align*}
 \widehat\beta(K_N)
 &\le\Lambda_0(T)+e+\frac{\log k}{\log N},\\
 \inf_{3\le\theta\le4}
 \left\{\frac{\theta+3}{2}\widehat\beta(i)
     -\widehat B_\theta(i,j)-2\widehat B_2(j,K_N)\right\}
 &\le\Lambda_1(T)+\frac{13}{2}e
                         +\frac{2\log k}{\log N}.
\end{align*}
The second inequality compares the expressions pointwise in $\theta$
at $(s,u)=(i/\log N,j/\log N)$ before taking the infimum.
Both right sides are at most $-3\eta/4\le-\eta_e$ for large $N$.
The local condition on $[0,T]$ implies the same condition on $[0,S_N]$.
Theorem~\ref{thm:ub-empirical-upper} now applies.  The choices and
errors were uniform over the cut array, proving the stated uniformity.
\end{proof}

\begin{lemma}[Buffered pointwise upper bound]
\label{lem:buffered-pointwise-upper-bound}
Under the ordered approximation and local-good hypotheses of
Theorem~\ref{thm:deterministic-upper-bound}, if
$\Lambda_{\UB}(T)\le-\eta$ for some $\eta>0$, then
\[
 d_{\TV}(P(n,\lfloor T\log N\rfloor),\nuUnif)\longrightarrow0.
\]
\end{lemma}
\begin{proof}
The defining maximum gives $\Lambda_0(T),\Lambda_1(T)\le-\eta$.
Apply Theorem~\ref{thm:deterministic-upper-bound}.  The conditional estimates used by that theorem are already
proved for the induced finite arrays.
\end{proof}

\begin{lemma}[Monotonicity of the upper margins]
\label{lem:upper-margin-monotonicity}
Each function $\Lambda_i(T)$, $0\le i\le3$, is nonincreasing in $T$.
Consequently, the full margin $\Lambda_{\UB}(T)$ in
\eqref{eq:global-upper-margin} is nonincreasing.
\end{lemma}

\begin{proof}
We first recall the signs of the profile costs defined in
\eqref{eq:one-column-costs}--\eqref{eq:integrated-profile-costs}.  Since
$p_{\max}(p)\le1$, one has $a_p\ge0$.  Moreover, for every $\theta\ge1$,
\[
  \sum_{\ell=0}^{k-1}p_\ell^\theta
  \le\sum_{\ell=0}^{k-1}p_\ell=1,
\]
and hence $\psi_p(\theta)\ge0$.  Averaging as in
\eqref{eq:profile-costs} gives $a(t)\ge0$ and
$\psi_t(\theta)\ge0$.  It follows from
\eqref{eq:integrated-profile-costs} that $A$ and $B_\theta$ are
nonnegative and additive over adjacent time intervals.  In particular,
$A(0,T)$ is nondecreasing in $T$, while $\beta(T)$ in
\eqref{eq:global-beta} is nonincreasing; explicitly,
\begin{equation}
\label{eq:beta-monotonicity}
  0\le s_1\le s_2
  \quad\Longrightarrow\quad
  \beta(s_2)\le\beta(s_1).
\end{equation}
The identity
$\Lambda_0(T)=\beta(T)$ from \eqref{eq:global-lambda-zero} therefore proves
the claim for $\Lambda_0$.

We next treat $\Lambda_3$.  Fix $0<T_1<T_2$.  If $0\le s\le T_1$, the
additivity and nonnegativity of $B_2$ give
\begin{align}
  \beta(s)-B_2(s,T_2)
  &=\beta(s)-B_2(s,T_1)-B_2(T_1,T_2)\nonumber\\
  &\le\beta(s)-B_2(s,T_1)
  \le\Lambda_3(T_1),
  \label{eq:lambda-three-old-times}
\end{align}
where the last inequality is the defining supremum in
\eqref{eq:global-lambda-three}.  If $T_1<s\le T_2$, then
\begin{equation}
\label{eq:lambda-three-new-times}
  \beta(s)-B_2(s,T_2)
  \le\beta(s)
  \le\beta(T_1)
  =\beta(T_1)-B_2(T_1,T_1)
  \le\Lambda_3(T_1).
\end{equation}
Here the second inequality uses \eqref{eq:beta-monotonicity},
and the last again uses \eqref{eq:global-lambda-three}.  Taking the
supremum over $s\in[0,T_2]$ in
\eqref{eq:lambda-three-old-times}--\eqref{eq:lambda-three-new-times} yields
$\Lambda_3(T_2)\le\Lambda_3(T_1)$.

It remains to treat $\Lambda_1$ and $\Lambda_2$.  To recall their exact
coefficients from \eqref{eq:global-lambda-one} and
\eqref{eq:global-lambda-two}, set
\begin{equation}
\label{eq:upper-margin-coefficients}
  c_1(\theta)=\frac{\theta+3}{2},
  \qquad
  c_2(\theta)=\frac{\theta-1}{2}.
\end{equation}
Both coefficients are positive on $[3,4]$.  For $i\in\{1,2\}$ and
$0\le s\le u\le T$, define
\begin{equation}
\label{eq:upper-margin-integrand}
  F_i^T(s,u)
  :=\inf_{\theta\in[3,4]}
  \left\{c_i(\theta)\beta(s)-B_\theta(s,u)-2B_2(u,T)\right\}.
\end{equation}
Equations \eqref{eq:global-lambda-one} and
\eqref{eq:global-lambda-two} then read
\begin{equation}
\label{eq:upper-margin-as-supremum}
  \Lambda_i(T)=\sup_{0\le s\le u\le T}F_i^T(s,u),
  \qquad i\in\{1,2\}.
\end{equation}

Again fix $0<T_1<T_2$ and consider an arbitrary pair
$0\le s\le u\le T_2$.  There are three cases.  If $u\le T_1$, additivity
of $B_2$ in \eqref{eq:integrated-profile-costs} and
\eqref{eq:upper-margin-integrand} give the exact identity
\begin{equation}
\label{eq:upper-margin-case-one}
  F_i^{T_2}(s,u)
  =F_i^{T_1}(s,u)-2B_2(T_1,T_2)
  \le F_i^{T_1}(s,u)
  \le\Lambda_i(T_1),
\end{equation}
where the last inequality uses \eqref{eq:upper-margin-as-supremum}.

If $s\le T_1<u$, then, for every $\theta\in[3,4]$,
\[
  B_\theta(s,u)=B_\theta(s,T_1)+B_\theta(T_1,u),
\]
and both $B_\theta(T_1,u)$ and $B_2(u,T_2)$ are nonnegative.  Comparing
the defining integrands in \eqref{eq:upper-margin-integrand} pointwise in
$\theta$ and then taking their infima gives
\begin{equation}
\label{eq:upper-margin-case-two}
  F_i^{T_2}(s,u)
  \le F_i^{T_1}(s,T_1)
  \le\Lambda_i(T_1),
\end{equation}
where the last inequality again follows from
\eqref{eq:upper-margin-as-supremum}.

Finally, suppose $T_1<s\le u$.  Equation \eqref{eq:beta-monotonicity} gives
$\beta(s)\le\beta(T_1)$.  Since $c_i(\theta)>0$ by
\eqref{eq:upper-margin-coefficients} and the two $B$-terms are
nonnegative, a pointwise comparison in $\theta$ yields
\begin{equation}
\label{eq:upper-margin-case-three}
  F_i^{T_2}(s,u)
  \le F_i^{T_1}(T_1,T_1)
  \le\Lambda_i(T_1),
\end{equation}
with the last inequality supplied by
\eqref{eq:upper-margin-as-supremum}.  The three alternatives
\eqref{eq:upper-margin-case-one}--\eqref{eq:upper-margin-case-three} cover
every admissible pair $(s,u)$ at time $T_2$.  Taking the supremum in
\eqref{eq:upper-margin-as-supremum} proves
$\Lambda_i(T_2)\le\Lambda_i(T_1)$ for $i=1,2$.

We have now proved monotonicity for all four $\Lambda_i$.  Their maximum is
therefore nonincreasing by \eqref{eq:global-upper-margin}.
\end{proof}

\begin{proposition}[Four-margin critical-time upper bound]
\label{cor:full-margin-critical-upper-bound}
Under Definition~\ref{def:ub-arbitrarily-fine-approximation}, for every
finite $T>T_{\UB}$,
\[
 d_{\TV}(P(n,\lfloor T\log N\rfloor),\nuUnif)\longrightarrow0.
\]
\end{proposition}
\begin{proof}
Put $\mathcal G=\{S>0:\Lambda_{\UB}(S)<0\}$.
If $\mathcal G$ is empty, then $T_{\UB}=+\infty$ and there is no
finite $T>T_{\UB}$.  Otherwise, for any such $T$, the definition of
the infimum supplies $S\in\mathcal G$ with $S<T$.
Lemma~\ref{lem:upper-margin-monotonicity} implies
$\Lambda_{\UB}(T)\le\Lambda_{\UB}(S)<0$.
Fix $\eta_T=-\Lambda_{\UB}(T)>0$.
Choose a sufficiently fine regular partition from
Definition~\ref{def:ub-arbitrarily-fine-approximation} for the spatial
cutoff in Theorem~\ref{thm:deterministic-upper-bound}; that spatial
cutoff depends only on $k,T,\eta_T$, not on $g$ or $\gamma$.
Once the partition and its fixed $g,\gamma,\bar\rho$ are given,
choose $0<\rho\le\min\{\bar\rho,\rho_0\}$ and then
$0<\varphi\le\varphi_0$.  The two cut hypotheses hold eventually
by the definition.  Lemma~\ref{lem:buffered-pointwise-upper-bound} proves convergence.
No continuity of the margins at their critical time is needed.
\end{proof}

\section{Dominance of the auxiliary upper-bound margins}
\label{sec:dominance-upper-margins}

The four-margin critical-time upper bound is
Proposition~\ref{cor:full-margin-critical-upper-bound}.  The goal of this section is
to prove that the two auxiliary margins are pointwise dominated by
$\Lambda_0$, and then to identify the four-margin threshold with the
reduced threshold.  The pointwise reduction is
Theorem~\ref{thm:dominance-reduction}; its buffered and threshold consequences
are Proposition~\ref{cor:buffered-consequence} and
\eqref{eq:upper-threshold-identity}, respectively.

We record the notation for the dominance argument.
The simplex $D_k$ is defined in \eqref{eq:probability-simplex}.  The
one-column quantities $p_{\max}$, $a_p$, and $\psi_p$ are defined in
\eqref{eq:one-column-costs}; their time averages are defined in
\eqref{eq:profile-costs}; and the integrated costs $A$ and $B_\theta$ are
defined in \eqref{eq:integrated-profile-costs}.  The whole-deck exponent
$\beta$ and the margins $\Lambda_0,\Lambda_1,\Lambda_2,\Lambda_3$, and
$\Lambda_{\UB}$ are defined in
\eqref{eq:global-beta}--\eqref{eq:global-upper-margin}.  The full and
reduced thresholds are defined in \eqref{eq:full-upper-threshold} and
\eqref{eq:reduced-upper-threshold}.  We use these definitions below
without repeating them.

\begin{lemma}[The basic one-column inequality]
\label{lem:basic-one-column-inequality}
For every $p\in D_k$ and every $\theta\ge1$,
\begin{equation}
\label{eq:one-column-dominance}
  \psi_p(\theta)\ge(\theta-1)a_p.
\end{equation}
Consequently,
\begin{equation}
\label{eq:integrated-one-column-dominance}
  B_\theta(r,s)\ge(\theta-1)A(r,s).
\end{equation}
In particular,
\begin{equation}
\label{eq:basic-b-two-three-dominance}
  B_2(r,s)\ge A(r,s),
  \qquad
  B_3(r,s)\ge2A(r,s).
\end{equation}
\end{lemma}

\begin{proof}
Let $m=p_{\max}(p)$, where $p_{\max}$ is defined in
\eqref{eq:one-column-costs}. Since $p_\ell\le m$ for every $\ell$,
\[
  p_\ell^\theta=p_\ell p_\ell^{\theta-1}\le p_\ell m^{\theta-1}.
\]
Summing over $\ell$ gives
\[
  \sum_{\ell=0}^{k-1}p_\ell^\theta
  \le m^{\theta-1}\sum_{\ell=0}^{k-1}p_\ell
  =m^{\theta-1}.
\]
The last equality uses \eqref{eq:probability-simplex}. Taking $-\log$ and
using \eqref{eq:one-column-costs} yields
\[
  \psi_p(\theta)=-\log\sum_{\ell=0}^{k-1}p_\ell^\theta
  \ge-\log m^{\theta-1}
  =(\theta-1)\log\frac1m
  =(\theta-1)a_p.
\]
After averaging over $\mu_t$ as in \eqref{eq:profile-costs} and integrating
over $[r,s]$ as in \eqref{eq:integrated-profile-costs}, we get
\[
  B_\theta(r,s)\ge(\theta-1)A(r,s).
\]
\end{proof}

\begin{theorem}[Reduction to $\Lambda_0$ and $\Lambda_1$]
\label{thm:dominance-reduction}
For every $T>0$,
\begin{equation}
\label{eq:auxiliary-margin-dominance}
  \Lambda_2(T)\le\Lambda_0(T),
  \qquad
  \Lambda_3(T)\le\Lambda_0(T).
\end{equation}
Therefore
\begin{equation}
\label{eq:two-margin-reduction}
  \Lambda_{\UB}(T)=\max\{\Lambda_0(T),\Lambda_1(T)\}.
\end{equation}
\end{theorem}

\begin{proof}
First consider $\Lambda_3$. For any $0\le s\le T$,
\eqref{eq:basic-b-two-three-dominance} in
Lemma~\ref{lem:basic-one-column-inequality} gives
\[
  B_2(s,T)\ge A(s,T).
\]
Therefore, using \eqref{eq:global-beta},
\eqref{eq:integrated-profile-costs}, and
\eqref{eq:global-lambda-zero},
\[
  \beta(s)-B_2(s,T)
  \le\beta(s)-A(s,T)
  =1-A(0,s)-A(s,T)
  =1-A(0,T)
  =\Lambda_0(T).
\]
Taking the supremum over $s$ in \eqref{eq:global-lambda-three} gives
$\Lambda_3(T)\le\Lambda_0(T)$.

Now consider $\Lambda_2$. Fix $0\le s\le u\le T$. Since the infimum in
\eqref{eq:global-lambda-two} is over $[3,4]$, we may evaluate the
expression at $\theta=3$:
\[
  \inf_{\theta\in[3,4]}
  \left\{\frac{\theta-1}{2}\beta(s)-B_\theta(s,u)-2B_2(u,T)\right\}
  \le \beta(s)-B_3(s,u)-2B_2(u,T).
\]
Using \eqref{eq:basic-b-two-three-dominance} from
Lemma~\ref{lem:basic-one-column-inequality}, this is at most
\[
  \beta(s)-2A(s,u)-2A(u,T)=\beta(s)-2A(s,T),
\]
where the equality uses \eqref{eq:integrated-profile-costs}. Since
\[
  \beta(s)-2A(s,T)
  =1-A(0,s)-2A(s,T)
  =1-A(0,T)-A(s,T),
\]
by \eqref{eq:global-beta} and \eqref{eq:integrated-profile-costs}, and
$A(s,T)\ge0$ follows from
\eqref{eq:probability-simplex}--\eqref{eq:integrated-profile-costs}, we get
\[
  \beta(s)-2A(s,T)\le1-A(0,T)=\Lambda_0(T),
\]
where the last equality is \eqref{eq:global-lambda-zero}. Thus every
$(s,u)$ contributes at most $\Lambda_0(T)$ to $\Lambda_2(T)$, and taking
the supremum in \eqref{eq:global-lambda-two} gives
\[
  \Lambda_2(T)\le\Lambda_0(T).
\]
Together with the bound for $\Lambda_3$, this proves
\eqref{eq:auxiliary-margin-dominance}. The identity
\eqref{eq:two-margin-reduction} follows from
\eqref{eq:global-upper-margin} and
\eqref{eq:auxiliary-margin-dominance}.
\end{proof}

\begin{proposition}[Buffered consequence]
\label{cor:buffered-consequence}
If for some $T>0$ and $\eta>0$,
\[
  \Lambda_0(T)\le-\eta,
  \qquad
  \Lambda_1(T)\le-\eta,
\]
then automatically
\[
  \Lambda_2(T)\le-\eta,
  \qquad
  \Lambda_3(T)\le-\eta,
\]
and hence
\[
  \Lambda_{\UB}(T)\le-\eta.
\]
\end{proposition}

\begin{proof}
By \eqref{eq:auxiliary-margin-dominance} in
Theorem~\ref{thm:dominance-reduction},
\[
  \Lambda_2(T)\le\Lambda_0(T)\le-\eta,
  \qquad
  \Lambda_3(T)\le\Lambda_0(T)\le-\eta.
\]
Together with the assumptions on $\Lambda_0$ and $\Lambda_1$, all four
margins are at most $-\eta$.  The definition
\eqref{eq:global-upper-margin} then gives
$\Lambda_{\UB}(T)\le-\eta$.
\end{proof}

The two critical-time thresholds therefore agree.
Recall the reduced threshold from \eqref{eq:reduced-upper-threshold}.
The pointwise identity \eqref{eq:two-margin-reduction}, proved in
Theorem~\ref{thm:dominance-reduction}, shows that the negativity sets in the two
threshold definitions agree.  Comparing \eqref{eq:full-upper-threshold} with
\eqref{eq:reduced-upper-threshold} therefore gives the exact identity
\begin{equation}
\label{eq:upper-threshold-identity}
  T_{\UB}^{\mathrm{red}}=T_{\UB}.
\end{equation}
Thus Theorem~\ref{thm:upper-critical-reduced} may be stated using only
$\Lambda_0$ and $\Lambda_1$.  The quantities $\Lambda_2$ and $\Lambda_3$
remain useful proof diagnostics, but they do not affect its upper critical
time.

\begin{proof}[Proof of Theorem~\ref{thm:upper-critical-reduced}]
Theorem~\ref{thm:dominance-reduction} gives
$\Lambda_2(T),\Lambda_3(T)\le\Lambda_0(T)$, so the two negativity
sets coincide and $T_{\UB}=T_{\UB}^{\mathrm{red}}$, including the
case of an empty negativity set.  Proposition~\ref{cor:full-margin-critical-upper-bound} proves
the asserted convergence for every finite $T$ above this threshold.
For fixed $\varepsilon\in(0,1)$ and each such $T$, eventually
$t_{\mathrm{mix}}^{(N)}(\varepsilon)\le\lfloor T\log N\rfloor$.
Divide by $\log N$, take the limsup, and let $T$ decrease to the
finite threshold.  If the threshold is $+\infty$, the asserted
limsup inequality is automatic.
\end{proof}

\begin{remark}
\label{rem:dominance-reduction-assumptions}
The reduction in Theorem~\ref{thm:dominance-reduction} uses the admissible range
$[3,4]$ in the definition \eqref{eq:global-lambda-two}, specifically its
endpoint $\theta=3$.  It also uses the one-column definitions in
\eqref{eq:one-column-costs}.  Under different shuffle assumptions or a
different range of $\theta$, the proof would need to be rechecked.
\end{remark}

\begin{comment}

\section{Lower bound: two obstructions}

The lower bound has two mechanisms, matching the two main ways the upper-bound margins can fail.

\subsection{Largest prefix block}

\begin{definition}[Largest-prefix obstruction]
\label{def:lb-overview-largest-prefix-obstruction}
We say $\mathrm{LP}(T,\eta)$ holds if
\[
  A(0,T)\le1-\eta.
\]
For each $j$, choose
\[
  \ell_*(j)\in\arg\max_\ell p_\ell^{(j)},
\]
and let
\[
  x_*=(\ell_*(1),\ldots,\ell_*(K)).
\]
\end{definition}

\begin{proposition}[Largest-prefix lower bound]
\label{prop:lb-overview-largest-prefix-lower-bound}
If the cuts are almost-$\mu_\bullet$-like and $\mathrm{LP}(T,\eta)$ holds, then
\[
  \dTV(P(n,K),\nuUnif)\to1.
\]
\end{proposition}

\paragraph{Sketch.}
Let $Y_*$ count row strings equal to $x_*$. Then
\[
  \E Y_*
  =
  N\prod_{j=1}^K p_{\ell_*(j)}^{(j)}
  =
  N^{1-A(0,T)+o(1)}
  \ge N^{\eta/2}.
\]
Hypergeometric concentration gives $Y_*\ge N^{\eta/3}$ with high probability. These rows form one connected component of the shuffle graph, so the inverse shuffled permutation has a contiguous increasing subsequence of length $N^{\eta/3}$. A uniform permutation has such a subsequence with probability at most $N/(N^{\eta/3}!)=o(1)$.

\subsection{Cold spots}

Fix $\theta\in[3,4]$ and $u\in[0,T]$. Define
\[
  h_{\theta,u}
  =
  1-J_\theta(0,u)+\int_0^u H_t(\theta)\,dt,
\]
and
\[
  \gamma_{\theta,u}
  =
  2-2J_\theta(0,u)-2J_2(u,T)
  +\int_0^u H_t(\theta)\,dt+\int_u^T H_t(2)\,dt.
\]
Equivalently,
\[
  \gamma_{\theta,u}-\frac12h_{\theta,u}
  =
  \frac32+\frac{\theta-3}{2}J_\theta(0,u)
  -B_\theta(0,u)-B_2(u,T).
\]
Near the canonical point $J_\theta(0,u)=1/2$, write
\[
  \Xi(\theta,u)
  =
  \frac{3+\theta}{4}-B_\theta(0,u)-B_2(u,T).
\]

\begin{definition}[Cold-spot obstruction]
\label{def:lb-overview-cold-spot-obstruction}
We say $\mathrm{CS}(T,\eta)$ holds if there exist $\theta\in[3,4]$ and $u\in[0,T]$ such that
\[
  J_\theta(0,u)\le\frac12-\eta,
  \qquad
  h_{\theta,u}\ge\frac12+\eta,
  \qquad
  \gamma_{\theta,u}\ge\frac12h_{\theta,u}+\eta.
\]
A convenient near-critical sufficient condition is
\[
  |J_\theta(0,u)-1/2|\le\eta/10,
  \qquad
  \Xi(\theta,u)\ge\eta.
\]
\end{definition}

\begin{definition}[Collision-likely prefixes and cold set]
\label{def:lb-overview-collision-likely-prefixes-cold-set}
A prefix $x$ of length $\lfloor u\log N\rfloor$ is collision-likely if in each space-time cell $(q,i)$ its digit counts are close to the $\theta$-tilted law $(p^{(i)})^{[\theta]}$, with the same exceptional-digit truncation as above. Let $\PreCL(\theta,u)$ be the set of such prefixes, and write $J_x$ for the deterministic lexicographic cylinder interval of $x$ under the product measure determined by the column proportions. Define
\[
  H=H(n;\theta,u)
  =
  \Z\cap\bigcup_{x\in\PreCL(\theta,u)}NJ_x.
\]
\end{definition}

\begin{lemma}[Cold-set geometry]
\label{lem:lb-overview-cold-set-geometry}
Assume $\mathrm{CS}(T,\eta)$. Then
\[
  |H|=N^{h_{\theta,u}+o(1)},
  \qquad
  |\partial H|\le2|H|^{1/2},
\]
and each component of $H$ has size at least
\[
  N^{1-J_\theta(0,u)-o(1)}.
\]
\end{lemma}

\paragraph{Sketch.}
Each cold prefix has expected interval length $N^{1-J_\theta(0,u)+o(1)}$. The number of such prefixes is
\[
  N^{\int_0^u H_t(\theta)\,dt+o(1)}.
\]
The intervals are disjoint in lexicographic order, giving the formula for $|H|$. The boundary bound follows because the components are much larger than the number of components.

\begin{definition}[Collision-likely full strings]
\label{def:lb-overview-collision-likely-full-strings}
A full string is collision-likely if it starts with a collision-likely prefix, passes through a short deterministic buffer containing two good columns, and on the tail $(u,T]$ has digit counts close to the $2$-tilted law $p^{[2]}$ in each space-time cell.
Let $\CL(\theta,u)$ denote the set of these collision-likely full strings.
\end{definition}

\begin{lemma}[Localization]
\label{lem:lb-overview-localization}
With probability $1-o(1)$, every collision-likely full string occurring in the sorted row list has its index inside $H$.
\end{lemma}

\paragraph{Sketch.}
Uniform concentration of prefix interval endpoints puts every cold prefix close to its expected interval. The two good buffer columns create left and right buffers larger than the endpoint fluctuations, forcing the full string to lie in the corresponding cold interval.

\begin{lemma}[Cold collision count]
\label{lem:lb-overview-cold-collision-count}
Assume $\mathrm{CS}(T,\eta)$. With probability $1-o(1)$,
\[
  |E(G)\cap E(H)|\ge |H|^{1/2+\delta}
\]
for some $\delta>0$.
\end{lemma}

\paragraph{Sketch.}
For a collision-likely full string $v$, let $A_v$ be the event that at least two row strings equal $v$. Inclusion-exclusion gives the first moment
\[
  \E\sum_{v\in\CL(\theta,u)}\1_{A_v}=N^{\gamma_{\theta,u}+o(1)}.
\]
The same four-row calculation as in the homogeneous cold-spot proof gives a negligible relative variance. Thus there are $N^{\gamma_{\theta,u}-o(1)}$ collision-producing strings, all localized in $H$. Since $\gamma_{\theta,u}\ge h_{\theta,u}/2+\eta$, the claimed power bound follows.

\begin{lemma}[Lalley cold-spot test]
\label{lem:lb-overview-lalley-cold-spot-test}
Let $H$ be deterministic with $|H|\ge N^{1/2+\eta}$ and $|\partial H|\le2|H|^{1/2}$. If
\[
  \Pp(|E(G)\cap E(H)|\ge |H|^{1/2+\delta})\to1,
\]
then $P(n,K)$ is asymptotically singular from uniform.
\end{lemma}

\paragraph{Sketch.}
Edges of $G$ inside $H$ force ascents after sorting within graph components. Under a uniform permutation, ascents on odd and even edges are independent Bernoulli$(1/2)$ variables up to the boundary error. Hoeffding's inequality makes the same excess event exponentially unlikely under uniform.

\begin{proposition}[Cold-spot lower bound]
\label{prop:lb-overview-cold-spot-lower-bound}
If the cuts are almost-$\mu_\bullet$-like and $\mathrm{CS}(T,\eta)$ holds, then
\[
  \dTV(P(n,K),\nuUnif)\to1.
\]
\end{proposition}

\paragraph{Sketch.}
Construct $H(n;\theta,u)$ from the witnessing pair $(\theta,u)$. The geometry, localization, and collision-count lemmas verify the hypotheses of the Lalley test.

\begin{theorem}[Lower-bound blueprint]
\label{thm:lower-bound-blueprint}
If the deterministic cuts are almost-$\mu_\bullet$-like and either $\mathrm{LP}(T,\eta)$ or $\mathrm{CS}(T,\eta)$ holds, then
\[
  \dTV(P(n,K),\nuUnif)\to1.
\]
For randomized-$\mu_\bullet$-like cuts, the same conclusion follows after conditioning on typical deterministic realizations and using the many-to-one union bound from the homogeneous proof.
\end{theorem}

\paragraph{Sketch.}
The largest-prefix case is the largest-prefix proposition. The cold-spot case is the cold-spot proposition. In the random case, take the union of the deterministic tests over all typical cut sequences; the number of such sequences is $\exp(O((\log N)^2))$, while the false-positive probability under uniform is $\exp(-N^c)$.

\section{Consistency between upper and lower quantities}

The largest-prefix obstruction is exactly the failure of the upper singleton margin at the terminal time:
\[
  N_3(0,T;1)=1-A(0,T).
\]
Thus upper mixing requires this quantity to be negative, while the lower obstruction occurs when it is positive.

The cold-spot obstruction matches the corrected large-block first-moment margin. At the canonical point $J_\theta(0,u)=1/2$,
\[
  2\Xi(\theta,u)=N_1(0,0,u;1).
\]
Thus cold spots are the positive-margin counterpart of the large-block upper-bound estimate.

If $\mu_t\equiv\mu$, put
\(\psi_\mu(\theta)=\E_\mu\psi_p(\theta)\). Then
\[
  A(0,T)=T\,\E_\mu\left[\log\frac1{p_{\max}(p)}\right],
\]
so $\mathrm{LP}$ becomes
\[
  T<\frac1{\E_\mu[\log(1/p_{\max}(p))]}.
\]
For cold spots, choose $\theta_\mu$ satisfying
\[
  \psi_\mu(\theta_\mu)=2\psi_\mu(2),
\]
and choose $u$ so that $J_{\theta_\mu}(0,u)=1/2$. Then
\[
  \Xi(\theta_\mu,u)=\frac{3+\theta_\mu}{4}-T\psi_\mu(2),
\]
so the obstruction is
\[
  T<\frac{3+\theta_\mu}{4\psi_\mu(2)}.
\]
These are exactly the two lower-bound constants in the homogeneous theorem.
\end{comment}

\clearpage
\section*{Matching lower bound from prefix blocks and cold spots}
\addcontentsline{toc}{section}{Matching lower bound from prefix blocks and cold spots}

\section{Main lower-bound theorem}

Put $K_T=\lfloor T\log N\rfloor$. The quantities $\beta,\Lambda_0,\Lambda_1$
retain their definitions from the common setup; explicitly,
\[
 \beta(s)=1-A(0,s),\qquad \Lambda_0(T)=\beta(T).
\]
\[
 \Lambda_1(T)=\sup_{0\le s\le u\le T}\inf_{3\le\theta\le4}
 \left\{\frac{\theta+3}{2}\beta(s)-B_\theta(s,u)-2B_2(u,T)\right\}.
\]
The almost-profile hypothesis is used with its ordered parameter convention:
for each fixed positive margin, the finite regular partition and sufficiently
small positive mesh and frequency tolerances are fixed before $N$ tends to
infinity, and the required estimates hold for all sufficiently large $N$.
Precisely, the argument below supplies a finite partition, a number $\rho_0>0$,
and, for each $0<\rho\le\rho_0$, a number $\varphi_0>0$; it suffices that
Definition~\ref{def:almost-time-profile-like-cuts} holds with that partition
and $0<\varphi\le\varphi_0$. The local hypothesis is
Definition~\ref{def:good-columns} on $[0,T]$, with this $\rho$ and fixed
positive good-column threshold $\chi_g$ and density $\gamma$.
The constants may depend on these fixed quantities. No rate at which any
parameter vanishes with $N$ is assumed.

\begin{theorem}[Buffered reduced lower bound]
\label{thm:buffered-reduced-lower-bound}
Fix $T>0$ and $\eta>0$, under the standing almost-$\mu_\bullet$-like
and local good-shuffle hypotheses just specified. If
\[
 \max\{\Lambda_0(T),\Lambda_1(T)\}\ge\eta,
\]
then
\[
 d_{\TV}(P_N^{K_T},\nuUnif)\longrightarrow1.
\]
The convergence is uniform over the deterministic cuts satisfying the fixed
parameter bounds.
\end{theorem}

The proof is given after Proposition~\ref{prop:cold-lower}. We first transfer
only the bounded costs $a$ and $\psi$ to the actual finite columns. The
largest-word test is Proposition~\ref{prop:largest-prefix}; otherwise
Lemma~\ref{lem:cold-witness} constructs a finite cold witness. Its location,
collision count, and observable ascent test are proved independently of the
upper-bound profile and first-moment estimates.

\section{Empirical profile costs}

Write $L=\log N$ and $S_N=K_T/L$.  On
$((j-1)/L,j/L]$, $1\le j\le K_T$, use the actual profile $p^{(j)}$.
For $0\le r\le s\le S_N$, define
\[
 (A_N,B_{\alpha,N},J_{\alpha,N},\mathcal H_{\alpha,N})(r,s)
 =\int_r^s(a_{p_N(t)},\psi_{p_N(t)}(\alpha),
                 I_{p_N(t)}(\alpha),H_{p_N(t)}(\alpha))\,dt,
\]
where $p_N(t)=p^{(j)}$ on the displayed column interval.
These are integrals of bounded step functions.  In particular, at grid
endpoints they are the normalized finite sums of the actual column costs.
Set
\[
 \beta_N(s)=1-A_N(0,s),\qquad
 F_{\alpha,N}(s,u)=\frac{\alpha+3}{2}\beta_N(s)
                 -B_{\alpha,N}(s,u)-2B_{2,N}(u,S_N).
\]
Set $\Lambda_{1,N}=\sup_{0\le s\le u\le S_N}
\inf_{3\le\alpha\le4}F_{\alpha,N}(s,u)$.
No distributional assumption on the order of the columns is made here.

For $\alpha\in[2,4]$, the one-column functions satisfy
\begin{equation}\label{eq:repair-global-one-column-bounds}
 0\le a_p\le\log k,\quad
 0\le I_p(\alpha)\le\log k,\quad
 0\le H_p(\alpha)\le I_p(\alpha),\quad
 0\le\psi_p(\alpha)\le(\alpha-1)\log k.
\end{equation}
Indeed $I_p(1)=H(p)\le\log k$, and, restricting the finite sums to the
positive coordinates,
$\partial_\alpha I_p(\alpha)=-\operatorname{Var}_{p^{[\alpha]}}
(\log(1/p_\ell))\le0$.  For $q=p^{[\alpha]}$, the inequality $\log x\le x-1$ gives
$\sum_{q_\ell>0}q_\ell\log(p_\ell/q_\ell)\le
\sum_{q_\ell>0}p_\ell-1\le0$, which proves $H_p(\alpha)\le I_p(\alpha)$.  The power-sum bound follows from convexity.
The same formulas show continuity in $\alpha$; zero coordinates contribute
zero.  We shall also use
\begin{equation}\label{eq:repair-two-comparisons}
 \psi_p(2)\le2a_p\le2I_p(\alpha),\qquad
 \psi_p(4)\ge2\psi_p(2).
\end{equation}
Thus every empirical endpoint displacement by $r$ changes an integrated
$I$ or $H$ by at most $r\log k$, uniformly in the cuts and $N$.


\begin{lemma}[Uniform transfer of bounded profile costs]
\label{lem:lower-profile-transfer}
Given $e>0$, the partition and parameters in the almost-profile hypothesis
can be chosen so that, for all sufficiently large $N$,
\begin{equation}
\label{eq:lower-profile-transfer}
 \sup_{0\le r\le s\le S_N}|A_N(r,s)-A(r,s)|\le e,\qquad
 \sup_{\substack{0\le r\le s\le S_N\\2\le\alpha\le4}}
 |B_{\alpha,N}(r,s)-B_\alpha(r,s)|\le e.
\end{equation}
In particular,
\begin{equation}
\label{eq:lower-margin-transfer}
 |\beta_N(S_N)-\Lambda_0(T)|\le e+O_k(1/\log N),\qquad
 |\Lambda_{1,N}-\Lambda_1(T)|\le13e/2+O_k(1/\log N).
\end{equation}
The mesh $\rho$ may also be required to lie below any prescribed positive
constant.
\end{lemma}

\begin{proof}
Let $f$ be any one of the functions $a_p$ or $\psi_p(\alpha)$,
$2\le\alpha\le4$. They are jointly continuous in their displayed
parameters on a compact set, are nonnegative, and are bounded by
$M=3\log k$. Choose a finite regular partition into $m$ cells on each of
which the oscillation of every such $f$ is at most $\omega$.
For a union of complete $\rho$-grid intervals, replace $f$ in each cell by
its representative value and apply Definition~\ref{def:almost-time-profile-like-cuts}.
There are at most $T/\rho+2$ such intervals. For an arbitrary interval
$[r,s]\subset[0,T]$, discard at most two incomplete grid intervals;
each has length at most $\rho$ and contains at most $\rho\log N+2$
column times. It follows that, uniformly in $f,r,s$,
\begin{align*}
 &\left|\frac1{\log N}\sum_{\substack{r\log N<j\le s\log N\\j\le K_T}}
 f(p^{(j)})-\int_r^s\E_{\mu_t}f\,dt\right|\\
 &\quad\le 2T\omega+Mm\varphi T+4M\rho
          +\frac{4Mm(T/\rho+2)+4M}{\log N}.
\end{align*}
The rounding term also covers converting closed grid intervals in the
definition to half-open intervals: at most one column per grid endpoint
changes its membership. The empirical step integral differs from this
sum by at most $2M/\log N$. Choose the partition first, then $\rho$,
then $\varphi$, and finally $N$ to make these errors at most $e$.
This proves \eqref{eq:lower-profile-transfer}.

The interval $(S_N,T]$ has length less than $1/\log N$.
On the common domain $0\le s\le u\le S_N$, the error in $F$ is at most
$(7/2+1+2)e+O_k(1/\log N)=13e/2+O_k(1/\log N)$.
Projecting any $s,u\le T$ to $\min(s,S_N),\min(u,S_N)$ changes the
macroscopic expression by $O_k(1/\log N)$, uniformly in $\alpha$,
because all its scalar integrands are bounded by $3\log k$.
Taking infima and suprema preserves a uniform error bound, proving
\eqref{eq:lower-margin-transfer}.
\end{proof}

\section{Exact prefix counts}

For a word \(w=(w_1,\ldots,w_m)\) of length \(m\le K_T\), let
\[
 \lambda_w=\prod_{j=1}^m p^{(j)}_{w_j},\qquad
 t_w=\sum_{\substack{v\in[k]_0^m\\v<_{\rm lex}w}}\lambda_v,
 \qquad C(w)=\#\{i:s_i|_m=w\}.
\]
Here \(s_i|_m\) is the length-\(m\) prefix of the sorted row \(s_i\).
Define the zero-based cuts
\[
 L_w=\#\{i:s_i|_m<_{\rm lex}w\},\qquad
 U_w=\#\{i:s_i|_m\le_{\rm lex}w\}.
\]
Thus the actual ranks occupied by this prefix are
\(I(B_w)=(L_w,U_w]\cap\mathbb Z\), and \(U_w-L_w=C(w)\).
Its deterministic probability interval is
\(J_w=[t_w,t_w+\lambda_w)\). These definitions also apply to the empty
word, whose mass is one and whose count is \(N\).

\begin{lemma}[Successive thinning transforms and tails]
\label{lem:successive-thinning-tail}
Let \(M\ge1\) and let \(P_1,\ldots,P_m\) be independent uniform
subsets of a fixed set of \(M\) labels, with deterministic sizes
\(m_1,\ldots,m_m\). Put
\[
 X=|P_1\cap\cdots\cap P_m|,\qquad
 q_j=m_j/M,\qquad Q=\prod_jq_j,\qquad \mu=MQ.
\]
For an empty intersection use \(X=M\) and \(Q=1\). For every \(s>0\),
\begin{equation}
\label{eq:thinning-transform}
 \E s^X\le(1-Q+Qs)^M\le \exp(\mu(s-1)).
\end{equation}
Consequently, for \(t>0\),
\begin{equation}
\label{eq:thinning-bernstein}
 \Pp(|X-\mu|\ge t)\le
 2\exp\left(-\frac{t^2}{4(\mu+t)}\right),
\end{equation}
and, with \(c_{\rm Ch}=(1-\log2)/2>0\),
\begin{equation}
\label{eq:thinning-lower-tail}
 \Pp(X\le\mu/2)\le e^{-c_{\rm Ch}\mu},
 \qquad \Pp(X<2)\le2e^{-\mu/2}.
\end{equation}
These statements apply to the count of any prescribed word in the shuffle
matrix. They also hold conditionally after any admissible fixed prescriptions
on \(b<N\) labelled rows: delete those rows, replace \(M\) by \(N-b\),
and subtract the prescribed occurrences of each digit from its column count.
\end{lemma}

\begin{proof}
First consider a uniform \(m\)-element subset of \(M\) labels and a fixed
\(x\)-element set. Its intersection size \(V\) has the hypergeometric
distribution \(\operatorname{Hyp}(M,x,m)\). By symmetry one may instead
sample \(x\) positions without replacement from a list of \(m\) copies
of \(s\) and \(M-m\) copies of one. Hence \(\E s^V\) is the mean
of all products of \(x\) distinct entries of this positive list.
This elementary symmetric mean is at most the \(x\)-th power of the
arithmetic mean, so
\begin{equation}
\label{eq:one-step-thinning-transform}
 \E s^V\le(1-q+qs)^x,\qquad q=m/M,
 \qquad s>0.
\end{equation}
Here is an elementary verification of the symmetric-mean inequality.
If two positive entries \(a,b\) are replaced by their average, their sum
is unchanged and their product increases. In the elementary symmetric
polynomial, the terms involving this pair are a nonnegative multiple of
\(ab\), a nonnegative multiple of \(a+b\), and a constant. Thus the
polynomial cannot decrease. Repeated averaging of a largest and a smallest
entry converges to the constant list with the same mean. Indeed, if $V$
is the sum of squared deviations from that mean and $a,b$ are the largest
and smallest entries, one such averaging reduces $V$ by $(a-b)^2/2$.
Since $V\le M(a-b)^2$, the new squared deviation is at most
$(1-1/(2M))V$. Thus $V\to0$, and continuity proves the claim.
The cases \(x=0,1\) are equalities.

Let \(X_j=|P_1\cap\cdots\cap P_j|\), with \(X_0=M\).
Conditionally on the first \(j-1\) subsets,
\(X_j\sim\operatorname{Hyp}(M,X_{j-1},m_j)\).
Starting from \(z_m=s\), set
\(z_{j-1}=1-q_j+q_jz_j\). Every \(z_j\) is positive, and
\(z_0=1-Q+Qs\). Iterating \eqref{eq:one-step-thinning-transform}
gives \(\E s^X\le z_0^M\). The inequality \(1+v\le e^v\)
gives the second bound in \eqref{eq:thinning-transform}.

For \(0\le h\le1\), both
\(e^h-1-h\le h^2\) and \(e^{-h}-1+h\le h^2\).
Substitution of \(s=e^{\pm h}\) into
\eqref{eq:thinning-transform} therefore yields
\(\E e^{\pm h(X-\mu)}\le e^{\mu h^2}\).
Markov's inequality with \(h=t/(2(\mu+t))\) bounds each tail by
\(\exp[-t^2/(4(\mu+t))]\), proving
\eqref{eq:thinning-bernstein}.
Taking \(s=1/2\) in \eqref{eq:thinning-transform} and applying
Markov's inequality at the thresholds \(2^{-\mu/2}\) and \(1/2\)
proves the two bounds in \eqref{eq:thinning-lower-tail}.

In the shuffle matrix each target-digit pile is a uniform fixed-size
subset of the row labels, independently across columns. After prescribed
symbols have been fixed on selected rows, exchangeability makes each
residual column uniform with its residual counts. The conditioning event
factors over columns, so these residual columns remain independent.
The same argument therefore applies with exactly the residual parameters
specified in the statement.
\end{proof}

\begin{lemma}[Uniform prefix-count concentration]
\label{cor:iterated-hypergeom}
For every fixed \(\varepsilon>0\), with probability \(1-o(1)\),
simultaneously for every word \(w\) of length at most \(K_T\),
\begin{equation}
\label{eq:all-prefix-counts}
 |C(w)-N\lambda_w|\le
 N^\varepsilon(\sqrt{N\lambda_w}+1).
\end{equation}
The convergence is uniform over all deterministic column counts.
\end{lemma}

\begin{proof}
Set \(A=N^\varepsilon\), \(\mu=N\lambda_w\), and
\(t=A(\sqrt\mu+1)\). For \(A\ge1\),
\(t^2/(\mu+t)\ge A/2\).
Lemma~\ref{lem:successive-thinning-tail} bounds the failure probability
for this word by \(2e^{-N^\varepsilon/8}\).
There are at most \((K_T+1)k^{K_T}\le (T\log N+1)N^{T\log k}\)
words of the indicated lengths. A union bound proves
\eqref{eq:all-prefix-counts}, including words of zero mass, whose counts
are identically zero.
\end{proof}

\section{Largest-prefix obstruction}

\begin{proposition}[Finite largest-word bound]
\label{prop:largest-prefix}
Let
\(\mu_*=N\prod_{j=1}^{K_T}\max_\ell p^{(j)}_\ell\).
For every integer \(2\le m\le N\) with \(m\le\mu_*/2\),
\begin{equation}
\label{eq:finite-largest-word-tv}
 d_{\TV}(P_N^{K_T},\nuUnif)
 \ge1-e^{-c_{\rm Ch}\mu_*}-\frac{N-m+1}{m!}.
\end{equation}
In particular, if \(\mu_*\ge N^b\) for a fixed \(b>0\), the
total-variation distance tends to one, uniformly over these column counts.
\end{proposition}

\begin{proof}
Choose a maximizing digit in every column, breaking ties by the smallest
digit, and let \(w_*\) be the resulting full word. Its count has mean
\(\mu_*\). By \eqref{eq:thinning-lower-tail} it is at least \(m\)
with probability at least \(1-e^{-c_{\rm Ch}\mu_*}\).
Equal full words form one component of the shuffle graph. By
Lemma~\ref{lem:ub-graph-representation}, the inverse shuffled permutation
therefore has an increasing run of length \(m\) with at least this
probability. Under a uniform permutation each interval of \(m\) consecutive
positions is increasing with probability \(1/m!\). A union bound over
its \(N-m+1\) possible starting positions proves
\eqref{eq:finite-largest-word-tv}. For the last assertion take
\(m=\lfloor N^{b/2}\rfloor\); the hypotheses force \(b\le1\), and
for all sufficiently large \(N\) this choice satisfies the required
inequalities and \((N-m+1)/m!\to0\).
\end{proof}
\section{The cold witness}

\begin{lemma}[Finite cold witness]\label{lem:cold-witness}
Fix $\eta_f>0$, and put $\kappa=\eta_f/1000$ and
$\kappa_0=\min\{\kappa/100,1/1000\}$.
Suppose $\Lambda_{1,N}\ge\eta_f$ and
$\beta_N(S_N)<\kappa_0$.
Choose a bridge length $r_b=2\rho$ with
$r_b\log k\le\kappa/100$, where the local good-column hypothesis holds
on every interval of length $\rho$.  For all sufficiently large $N$ there
are grid endpoints $s=m_s/L$, $u=m_u/L$, a tilt $\theta\in[3,4]$, and an
integer bridge length $b_N=\lceil r_bL\rceil$ such that
$0\le m_s\le m_u\le K_T-b_N$, the bridge columns
$m_u<j\le m_u+b_N$ contain two good columns, and, writing
$b=\beta_N(s)$, $J=J_{\theta,N}(s,u)$,
$H=\mathcal H_{\theta,N}(s,u)$, $J_t=J_{2,N}(u,S_N)$,
$H_t=\mathcal H_{2,N}(u,S_N)$,
\begin{gather}
 F_{\theta,N}(s,u)\ge\eta_f/2,\qquad
 J\le b/2-3\kappa,\label{eq:repair-witness-upper}\\
 \theta=3\quad\hbox{or}\quad J\ge b/2-5\kappa,
 \label{eq:repair-witness-lower}\\
 h:=b-J+H\ge\kappa,\qquad H\le b-J-\kappa,
 \label{eq:repair-witness-geometry}\\
 \Gamma:=2b-2J-2J_t+H+H_t,\qquad
 \Gamma-h/2\ge\kappa,\label{eq:repair-witness-collision}\\
 b-J-J_t\le\beta_N(S_N)<\kappa_0.
 \label{eq:repair-witness-mean}
\end{gather}
The constants are uniform in the finite cut sequence.
\end{lemma}

\begin{proof}
Choose $s_0\le u_0$ with
$F_{\alpha,N}(s_0,u_0)\ge7\eta_f/8$ for every $\alpha\in[3,4]$.
Write $b_0=\beta_N(s_0)$.  Nonnegativity of the $B$ terms gives
$b_0\ge\eta_f/4$, so $b_0/2-4\kappa>0$.
If $J_{3,N}(s_0,u_0)\le b_0/2$, choose $\theta=3$.
Otherwise, if $J_{4,N}(s_0,u_0)\ge b_0/2$, choose $\theta=4$.
In the remaining case the continuity and monotonicity in $\alpha$ give
$\theta\in(3,4)$ with $J_{\theta,N}(s_0,u_0)=b_0/2$.
This selection does not interchange a supremum and an infimum.

If $J_{\theta,N}(s_0,u_0)\le b_0/2-4\kappa$, leave $u_1=u_0$;
this case can occur only for $\theta=3$.  Otherwise choose
$u_1\in[s_0,u_0]$ with
$J_{\theta,N}(s_0,u_1)=b_0/2-4\kappa$.
Existence follows from continuity and the value zero at $s_0$.
In the $\theta=3$ or interior cases the decrease of $J$ is at most
$4\kappa$.  By \eqref{eq:repair-two-comparisons},
\[
 F_{\theta,N}(s_0,u_1)-F_{\theta,N}(s_0,u_0)
 =B_{\theta,N}(u_1,u_0)-2B_{2,N}(u_1,u_0)
 \ge-4J_{\theta,N}(u_1,u_0)\ge-16\kappa.
\]
For $\theta=4$ the same difference is nonnegative, even if the
backward displacement is large.  Hence in every case
$F_{\theta,N}(s_0,u_1)\ge7\eta_f/8-16\kappa$.

Set
\[
 m_s=\lfloor s_0L\rfloor,\qquad
 m_u=\min\{\lfloor u_1L\rfloor,K_T-b_N\}.
\]
These integers are ordered for large $N$.  To check the only nontrivial
case, observe that
\[
 (S_N-s_0)\log k\ge A_N(s_0,S_N)
 =b_0-\beta_N(S_N)>\eta_f/4-\kappa_0.
\]
Consequently $s_0<S_N-r_b-2/L$ for large $N$, and clipping $m_u$ to
$K_T-b_N$ cannot move it below $m_s$.
The endpoint errors satisfy $0\le s_0-s<1/L$ and
$0\le u_1-u\le r_b+2/L$.
Equation~\eqref{eq:repair-global-one-column-bounds} bounds the resulting
change of $F$ by $5r_b\log k+20\log k/L$ and that of
$J-b/2$ by $r_b\log k+4\log k/L$.
These errors are smaller than $\kappa/10$ for large $N$.
The retained slack proves \eqref{eq:repair-witness-upper} and
\eqref{eq:repair-witness-lower}.
The bridge, whose length is at least $2\rho$, contains a complete
source grid interval of length $\rho$ lying inside $[0,S_N]$; hence it
contains at least $\gamma\rho L-1$ good columns, which exceeds two.

The exact identities $H=\theta J-B_{\theta,N}(s,u)$ and
$H_t=2J_t-B_{2,N}(u,S_N)$ imply
\[
 2\Gamma-h=F_{\theta,N}(s,u)+(\theta-3)(J-b/2)
 \ge\eta_f/2-5\kappa.
\]
This is greater than $2\kappa$.
Also $H\le J$ and $J\le b/2-3\kappa$ give
$H\le b-J-6\kappa$ and $h\ge b-J\ge b/2+3\kappa$.
Finally $J\ge A_N(s,u)$ and $J_t\ge A_N(u,S_N)$ give
\eqref{eq:repair-witness-mean}.
\end{proof}

\section{Finite type classes and cold sets}

The construction below uses only the actual columns.  In particular it
needs no temporal approximation for the tilted observables $I$ and $H$.
Fix a sufficiently fine $\xi$-regular simplex partition, with a regularity constant $C=C_k$ independent of $\xi$ and representatives
$p_R$, where $0<\xi<1/(4k)$.  There are $m=m(\xi)$ cells, independent of $N$.
On a chosen set $I$ of column positions, let $I_R$ be the positions in cell
$R$ and $n_R=|I_R|$.  The positions retain their chronological order.
For $\alpha\in[2,4]$ put
\[
 G_R=\{\ell:p_R(\ell)>2\xi\},\qquad
 q_R(\ell)=\frac{p_R(\ell)^\alpha\mathbf1_{G_R}(\ell)}
                 {\sum_{a\in G_R}p_R(a)^\alpha}.
\]
Choose integers $m_{R,\ell}\ge0$ supported on $G_R$, with
$\sum_\ell m_{R,\ell}=n_R$ and
$|m_{R,\ell}-n_Rq_R(\ell)|\le1$; floor the counts and distribute the
remaining units deterministically.  Let $\mathcal C^0_I$ be the words on
$I$ with exactly these counts in each cell.  This class is nonempty.
For every word write $\lambda_x=\prod_{j\in I}p^{(j)}(x_j)$.
Partition $\mathcal C^0_I$ into the integer mass bins
\[
 \{x:j\le-\log\lambda_x<j+1\},\qquad j\in\mathbb Z_{\ge0}.
\]
Select a bin of maximal cardinality, with deterministic tie breaking, and
call it $\mathcal C_I$.  If its index is $j_I$, define
\[
 J_I^*=j_I/L,\qquad H_I^*=L^{-1}\log|\mathcal C_I|.
\]
These quantities may depend on $N$; that dependence is intentional.

\begin{lemma}[Finite type estimates]\label{lem:type-estimates}
For $|I|\le TL$ and $\alpha\in[2,4]$, uniformly over the cuts, the chosen
position set, and the selected word,
\begin{gather}
 e^{-1}N^{-J_I^*}<\lambda_x\le N^{-J_I^*},\qquad
 |\mathcal C_I|=N^{H_I^*},\label{eq:repair-bin-exact}\\
 \left|J_I^*-L^{-1}\sum_{j\in I}I_{p^{(j)}}(\alpha)\right|
 +\left|H_I^*-L^{-1}\sum_{j\in I}H_{p^{(j)}}(\alpha)\right|
 \le C_{k,C}T\xi\log(1/\xi)
       +O_{k,m,\xi,T}\!\left(\frac{\log\log N+1}{\log N}\right).
 \label{eq:repair-types-error}
\end{gather}
Every digit occurring in a selected word has actual column probability
at least $\xi$.  Here $C$ is the constant in regularity, and is fixed before
$N$.
\end{lemma}

\begin{proof}
At a representative $r=p_R$, the omitted $\alpha$-power mass is at most
$k(2\xi)^\alpha$ and the full denominator is at least $k^{1-\alpha}$.
Thus the omitted tilted probability is $O_k(\xi^2)$, uniformly in
$\alpha\in[2,4]$.  Its cost is
$O_k(\xi^2\log(1/\xi))$, since
$x^\alpha\log(1/x)\le x^2\log(1/x)$ and this last function is increasing
on $(0,2\xi]$ when $\xi$ is sufficiently small.
Renormalization has the same or smaller order, because
$I_r(\alpha)\le\log k$.
For an actual $p$ in the same cell, retained coordinates satisfy
$|\log p_\ell-\log r_\ell|\le C\xi$, while omitted coordinates are at
most $3\xi$.  On the retained coordinates the normalized power weights
therefore have ratios between $e^{-8C\xi}$ and $e^{8C\xi}$.
Their costs are bounded by $\log(1/\xi)$.
Combining this ratio comparison with the omitted-cost estimate gives
\[
 \left|\sum_\ell q_R(\ell)\log(1/r_\ell)-I_p(\alpha)\right|
 \le C_{k,C}\xi\log(1/\xi).
\]
The identity $H(q_R)=\alpha\sum q_R\log(1/r)-\psi_R^G(\alpha)$,
where $\psi_R^G=-\log\sum_{G_R}r_\ell^\alpha$, and the same denominator
comparison give the analogous entropy bound.

Put $\bar J=L^{-1}\sum_R n_R\sum_\ell q_R(\ell)\log(1/p_R(\ell))$
and $\bar H=L^{-1}\sum_R n_RH(q_R)$.
The rounding of counts changes the logarithmic product at representatives
by at most $km\log(1/(2\xi))$.
Replacing representative coordinates by actual coordinates changes it by
at most $C\xi|I|$.  Consequently every $x\in\mathcal C^0_I$ satisfies
\[
 |{-L^{-1}\log\lambda_x}-\bar J|
 \le CT\xi+km\log(1/(2\xi))/L.
\]
The multinomial formula gives
$|\mathcal C^0_I|=\prod_Rn_R!/\prod_\ell m_{R,\ell}!$.
The elementary factorial estimate
$\log(n!)=n\log n-n+O(\log(n+1)+1)$ and the entropy continuity bound
\[
 |n_RH((m_{R,\ell}/n_R)_\ell)-n_RH(q_R)|
 \le C_k(\log(n_R+1)+1)
\]
show the asserted cardinality estimate: if $f(t)=t\log t$, with $f(0)=0$,
then $|f(x)-f(y)|\le\log(n+1)+2$ for $x,y\in[0,n]$ with
$|x-y|\le1$. This follows from $f'(t)=1+\log t$ away from zero and
the bound $|t\log t|\le1/e$ on $[0,1]$. Apply it to
$x=m_{R,\ell}$ and $y=n_Rq_R(\ell)$ and sum over the coordinates;
the empty-cell contribution is zero. Consequently,
$|L^{-1}\log|\mathcal C^0_I|-\bar H|
 =O_{k,m,T}((\log\log N+1)/L)$.
Since every allowed actual digit is at least $\xi$, there are at most
$2+TL\log(1/\xi)$ nonempty integer mass bins.
Keeping a largest bin therefore changes the logarithmic cardinality by
$O_{\xi,T}(\log\log N+1)$ only.  Its index divided by $L$ lies within
$1/L$ of the logarithmic mass of each selected word.
Combining these estimates proves both assertions.
\end{proof}

\subsection{Cold sets and bridges from the finite witness}

For a grid endpoint $s=m_s/\log N$, let $a_s$ be the length-$m_s$ word
which uses a maximizing digit in every column, with ties broken by the
smallest digit. For any word on a specified consecutive set of columns,
$\lambda_x$ denotes its product probability on those columns.
For $x$ on $(m_s,m_u]$, put
\[
 t_x^{\rm rel}=\sum_{\substack{x'\in[k]_0^{m_u-m_s}\\x'<_{\rm lex}x}}
                \lambda_{x'},\qquad
 J_x^{\rm rel}=[t_x^{\rm rel},t_x^{\rm rel}+\lambda_x).
\]

\begin{lemma}[Finite cold geometry and mass bounds]
\label{lem:cold-geometry}
Use the parameters of Lemma~\ref{lem:cold-witness}, with the additional
bridge budget $r_b\log k\le\kappa_0/100$. There are deterministic prefix
and tail families $\mathcal P_N,\mathcal T_N$, a bridge word $v_b$ with
reserve words $v_b^-,v_b^+$, and a deterministic set $H_{{\rm rel},N}$,
as constructed in the proof. The classes $\mathcal P_N,\mathcal T_N$
are obtained from Lemma~\ref{lem:type-estimates} on the column sets
$(m_s,m_u]$ with tilt $\theta$ and $(m_u+b_N,K_T]$ with tilt $2$,
respectively; denote their mass-bin exponents by $J_p^*,H_p^*$ and
$J_t^*,H_t^*$. Set
\[
 H_{{\rm rel},N}=\mathbb Z\cap
    \bigcup_{x\in\mathcal P_N}N\lambda_{a_s}J_x^{\rm rel},\qquad
 M_N=|H_{{\rm rel},N}|,\qquad H_q=(q+H_{{\rm rel},N})\cap[N].
\]
Define the real exponents
\[
 \ell_p=b-J_p^*,\quad h_N=\ell_p+H_p^*,\quad
 \ell_b=-\log_N\lambda_{v_b},\quad
 z_N=b-J_p^*-J_t^*-\ell_b,\quad
 g_N=2z_N+H_p^*+H_t^*.
\]
For all sufficiently large $N$,
\begin{gather}
 \ell_p\ge b/2+2\kappa,\qquad H_p^*\le\ell_p-\kappa,\qquad
 h_N\ge\kappa/2,\label{eq:cold-effective-geometry}\\
 M_N=N^{h_N+o(1)},\qquad
 N^{\kappa/3}\le M_N\le N,\qquad
 |\partial H_q|=O(\sqrt{M_N})\quad(0\le q\le N),
 \label{eq:cold-geometric-conclusions}\\
 g_N-h_N/2\ge\kappa/2,\qquad z_N\le2\kappa_0,
 \label{eq:cold-effective-collision}
\end{gather}
The full-word family $\mathcal Y_N=\{a_sxv_bz:x\in\mathcal P_N,z\in\mathcal T_N\}$
satisfies, uniformly in its words,
\begin{equation}
\label{eq:cold-full-word-masses}
 |\mathcal Y_N|=N^{H_p^*+H_t^*},\qquad
 e^{-2}N^{z_N}<N\lambda_y\le N^{z_N}.
\end{equation}
The reserves satisfy $v_b^-<_{\rm lex}v_b<_{\rm lex}v_b^+$ and
\begin{equation}
\label{eq:cold-reserve-means}
 \min_{x\in\mathcal P_N,\,\diamond\in\{-,+\}}
       N\lambda_{a_sxv_b^\diamond}\ge N^{b/2+\kappa}.
\end{equation}
Every digit selected by a full word or a reserve has column probability
at least a fixed $c>0$. All constants and vanishing errors are uniform
over the finite cut sequences under consideration.
\end{lemma}

\begin{proof}

Use the finite witness, and apply Lemma~\ref{lem:type-estimates} to
$I=(m_s,m_u]$ with $\alpha=\theta$.
Call the resulting class $\mathcal P_N$, with exponents $J_p^*,H_p^*$.
The ambient maximal prefix $a_s$ has exactly
$N\lambda_{a_s}=N^b$ because $s$ is a grid endpoint.
Set
\[
 \ell_p=b-J_p^*,\qquad h_N=\ell_p+H_p^*.
\]
Choose $\xi$ and then $N$ so that the error in
\eqref{eq:repair-types-error} is below a prescribed
$\varepsilon<\min(\kappa,\kappa_0)/100$.
Then
\begin{equation}\label{eq:repair-finite-geometric-buffers}
 \ell_p\ge b/2+2\kappa,\qquad
 H_p^*\le\ell_p-\kappa,\qquad h_N\ge\kappa/2.
\end{equation}
For each $x\in\mathcal P_N$ use its exact relative probability interval
$J_x^{\rm rel}$, and put
\[
 H_{{\rm rel},N}=\mathbb Z\cap\bigcup_{x\in\mathcal P_N}
                         N\lambda_{a_s}J_x^{\rm rel}.
\]
The intervals are disjoint and each has length between
$e^{-1}N^{\ell_p}$ and $N^{\ell_p}$.  Their integer counts differ from their
lengths by at most two.  Therefore, uniformly in $N$,
\[
 |H_{{\rm rel},N}|=N^{h_N+o(1)}\ge N^{\kappa/3},\qquad
 |\partial H_{{\rm rel},N}|\le2N^{H_p^*}
                         =o(|H_{{\rm rel},N}|^{1/2}).
\]
This does not assert that every probability interval is a distinct connected
component: adjacent selected intervals may merge, which only reduces the
boundary.  Translations clipped to $[N]$ have no larger cardinality and
have at most two additional boundary edges.

Choose two good bridge columns $g_L<g_R$ and define the word $v_b$ as follows.
At $g_L$ use the larger of two good digits; at $g_R$ use the smaller.
Use maximal digits at all other bridge positions.
For the reserve words $v_b^-,v_b^+$, change respectively the digit at
$g_L$ to the smaller and the digit at $g_R$ to the larger, keeping all
other digits equal to the central bridge word.
Then $v_b^-<_{\rm lex}v_b<v_b^+$, and each bridge word has probability at least
$\chi_g^2k^{-b_N}$, where $\chi_g$ is the fixed good-column threshold.
Writing $\ell_b=-L^{-1}\log\lambda_{v_b}$, we obtain
\[
 0\le\ell_b\le r_b\log k+O_{\chi_g,k}(1/L).
\]
The same estimate applies to both reserve words.  In particular every
reserve prefix $a_sxv_b^\pm$ has mean at least
$N^{b/2+\kappa}$ for large $N$, after the bridge budget is imposed.
This is the deterministic input for the later concentration/localization
argument; it makes no assertion of independence between reserve counts.

Let $\mathcal R=(m_u+b_N,K_T]$ and apply
Lemma~\ref{lem:type-estimates} there with $\alpha=2$.
Write $\mathcal T_N,J_t^*,H_t^*$ for the selected tail class and exponents.
The bridge deletion changes the empirical $J_2,H_2$ by at most
$(r_b+1/L)\log k$ each.
The full selected family is
\[
 \mathcal Y_N=\{a_sxv_bz:x\in\mathcal P_N, z\in\mathcal T_N\}.
\]
For $y\in\mathcal Y_N$ put $\mu_y=N\lambda_y$ and define
\[
 z_N=b-J_p^*-J_t^*-\ell_b,\qquad
 g_N=2z_N+H_p^*+H_t^*.
\]
The mass bins give the exact uniform bounds
\[
 e^{-2}N^{z_N}<\mu_y\le N^{z_N},\qquad
 |\mathcal Y_N|=N^{H_p^*+H_t^*},\qquad
 \sum_{y\in\mathcal Y_N}\mu_y^2=N^{g_N+o(1)}.
\]
Write
\[
 \delta_{p,J}=J_p^*-J,\quad \delta_{p,H}=H_p^*-H,\quad
 \delta_{t,J}=J_t^*-(J_t-d_J),\quad
 \delta_{t,H}=H_t^*-(H_t-d_H),
\]
where $d_J=J_{2,N}(u,u+b_N/L)$ and
$d_H=\mathcal H_{2,N}(u,u+b_N/L)$ are the deleted bridge integrals.
Each $\delta$ has absolute value at most $\varepsilon$ and
$0\le d_J,d_H\le d:=(r_b+1/L)\log k$.
For all sufficiently large $N$, our budgets give
$d,\ell_b\le\kappa_0/50$. Direct subtraction now yields
\[
 |h_N-h|\le2\varepsilon,\qquad
 g_N-h_N/2\ge\Gamma-h/2-7\varepsilon-d-2\ell_b\ge\kappa/2,
\]
and
\[
 z_N\le b-J-J_t+2\varepsilon+d
       <\kappa_0+2\varepsilon+d\le2\kappa_0.
\]
These prove \eqref{eq:cold-effective-collision}. The reserve mass lower
bound proved above gives \eqref{eq:cold-reserve-means}.
All selected digits have probability at least
$c=\min\{\xi,\chi_g,1/k\}>0$. Finally
$H_{{\rm rel},N}\subset\mathbb Z\cap[0,N\lambda_{a_s})\subset[0,N)\cap\mathbb Z$,
so $M_N\le N$. The bounds already established give
\eqref{eq:cold-geometric-conclusions}, including its uniform clipped-boundary
estimate, and complete the proof.
\end{proof}

\section{Uniform localization}

\begin{lemma}[Finite relative endpoint estimate]
\label{lem:finite-relative-endpoints}
Let \(a\) be a deterministic prefix of length \(l\), let
\(l\le m\le K_T\), and put \(\mu_a=N\lambda_a\).
For a suffix \(x\) on columns \(l+1,\ldots,m\), let
\(\lambda_x\) be its product mass and \(t_x\) the total product
mass of the length-\((m-l)\) suffixes strictly smaller than \(x\).
Define \(L_x^{\mathrm{rel}},U_x^{\mathrm{rel}}\) by counting rows
with prefix \(a\) whose length-\((m-l)\) suffixes are respectively
smaller than, or at most, \(x\).
If \(A>0\) and
\[
 |C(v)-N\lambda_v|\le A(\sqrt{N\lambda_v}+1)
 \quad\hbox{for every prefix \(v\) of length at most \(K_T\)},
\]
then simultaneously for all such \(a,x\),
\begin{equation}
\label{eq:relative-endpoint-finite-error}
\begin{split}
 &|L_x^{\mathrm{rel}}-\mu_at_x|
  +|U_x^{\mathrm{rel}}-\mu_a(t_x+\lambda_x)|\\
 &\hspace{6mm}\le
 A\bigl[(2\sqrt{kK_T}+1)\sqrt{\mu_a}+2kK_T+1\bigr].
\end{split}
\end{equation}
\end{lemma}

\begin{proof}
Partition the preceding and following suffixes by the first column
where they differ from \(x\).  The resulting exact identities are
\begin{align}
\label{eq:relative-left-side-exact}
 L_x^{\mathrm{rel}}-\mu_at_x
 &=\sum_{r=l+1}^m\sum_{d<x_r}
   \bigl[C(ax_{l+1}\cdots x_{r-1}d)
          -N\lambda_{ax_{l+1}\cdots x_{r-1}d}\bigr],\\
\label{eq:relative-right-side-exact}
 U_x^{\mathrm{rel}}-\mu_a(t_x+\lambda_x)
 &=C(a)-\mu_a
  -\sum_{r=l+1}^m\sum_{d>x_r}
   \bigl[C(ax_{l+1}\cdots x_{r-1}d)
          -N\lambda_{ax_{l+1}\cdots x_{r-1}d}\bigr].
\end{align}
Each side contains at most \(kK_T\) disjoint subcylinders of the
ambient cylinder.  Their means \(\mu_j\) therefore have sum at
most \(\mu_a\).  Cauchy--Schwarz bounds either side sum by
\[
 A\sum_j(\sqrt{\mu_j}+1)
 \le A\bigl[\sqrt{kK_T\mu_a}+kK_T\bigr].
\]
The additional ambient fluctuation in
\eqref{eq:relative-right-side-exact} is at most
\(A(\sqrt{\mu_a}+1)\).  Adding these bounds proves
\eqref{eq:relative-endpoint-finite-error}, including empty suffixes
and zero or small means.
\end{proof}

\begin{lemma}[Uniform cold-word localization]
\label{lem:localization}
Use the finite witness, the selected prefix family \(\mathcal P_N\),
the full-word family \(\mathcal Y_N\), and the cold set
\(H_{{\rm rel},N}\).  Put
\[
 H_q=(q+H_{{\rm rel},N})\cap[N],\qquad q_*=L_{a_s},\qquad
 b=\beta_N(s),\qquad N\lambda_{a_s}=N^b.
\]
With probability \(1-o(1)\), simultaneously for every
\(y\in\mathcal Y_N\), all rows carrying \(y\) have sorted
positions in \(H_{q_*}\).  On this event each distinct full word
that occurs at least twice supplies a distinct edge of
\(E(G)\cap E(H_{q_*})\).  The probability bound is uniform over
the permitted finite witnesses and cut arrays.
\end{lemma}

\begin{proof}
Take \(\varepsilon=\kappa/4\) in
Lemma~\ref{cor:iterated-hypergeom}.  On its simultaneous
prefix-count event, Lemma~\ref{lem:finite-relative-endpoints} bounds
both relative endpoint errors for every \(x\in\mathcal P_N\)
by
\[
 e_N=N^{\kappa/4}
       \bigl[(2\sqrt{kK_T}+1)N^{b/2}+2kK_T+1\bigr].
\]
Since \(J_p^*\ge0\), the inequality
\(\ell_p=b-J_p^*\ge b/2+2\kappa\) in
\eqref{eq:repair-finite-geometric-buffers} implies \(b\ge4\kappa\).
As \(K_T\le T\log N\), uniformly for all these witnesses,
\begin{equation}
\label{eq:repair-endpoint-error-concrete}
 e_N\le N^{b/2+\kappa/2}
 \quad\hbox{for all sufficiently large \(N\)}.
\end{equation}

The reserve words \(v_b^-\) and \(v_b^+\) differ from the central
bridge \(v_b\) at only the first and second distinguished good
column, respectively.  At the distinguished columns their digit
probabilities are at least \(\chi_g\), and all other bridge digits
are maximal.  Hence their bridge products are at least
\(\chi_g^2k^{-b_N}\), where
\(b_N=\lceil r_b\log N\rceil\) is the bridge length.  By
\eqref{eq:repair-bin-exact} and
\eqref{eq:repair-finite-geometric-buffers}, every reserve prefix has
mean at least
\[
 \mu_x^\pm:=\E C(a_sxv_b^\pm)
 \ge e^{-1}\chi_g^2 k^{-1}
       N^{\ell_p-r_b\log k}
 \ge N^{b/2+\kappa}
\]
for all sufficiently large \(N\), because
\(r_b\log k\le\kappa/100\).  The fixed factor involving
\(\chi_g\) is absorbed by the remaining positive exponent gap.

Apply Lemma~\ref{lem:successive-thinning-tail} to each reserve
prefix.  There are at most \(2k^{K_T}\le2N^{T\log k}\) such
prefixes, and each mean is at least \(N^{b/2+\kappa}\ge N^{3\kappa}\).
A union bound shows that, with probability \(1-o(1)\), all reserve
counts are at least half their respective means.  Together with
\eqref{eq:repair-endpoint-error-concrete}, this gives
\begin{equation}
\label{eq:repair-reserve-count-dominates}
 C(a_sxv_b^-),\ C(a_sxv_b^+)
 \ge \tfrac12N^{b/2+\kappa}\ge e_N+2
 \qquad(x\in\mathcal P_N)
\end{equation}
for all sufficiently large \(N\).  All preceding estimates use only
fixed constants \(k,T,\kappa,\chi_g\), so the asserted uniformity
is preserved when the empirical witness depends on \(N\).

Work on the intersection of these two high-probability events.
For a fixed \(x\), write
\[
 \ell=N\lambda_{a_s}t_x^{\rm rel},\qquad
 r=N\lambda_{a_s}(t_x^{\rm rel}+\lambda_x).
\]
The relative endpoints differ from \(\ell,r\) by at most
\(e_N\).  Every one-based sorted row rank \(i\) carrying a word
\(a_sxv_bz\in\mathcal Y_N\) satisfies
\[
 q_*+L_x^{\mathrm{rel}}+C(a_sxv_b^-)+1
 \le i\le
 q_*+U_x^{\mathrm{rel}}-C(a_sxv_b^+),
\]
because \(v_b^-<_{\rm lex}v_b<_{\rm lex}v_b^+\).
Equation~\eqref{eq:repair-reserve-count-dominates} therefore gives
\(q_*+\ell<i<q_*+r\).  Thus
\[
 i\in q_*+(\mathbb Z\cap[\ell,r))\subseteq H_{q_*}.
\]
The two spare units in the reserve bound resolve the one-based-rank
and half-open-interval conventions.  No union bound over tail words
\(z\) and no estimate on the random translation \(q_*\) are needed.

Copies of a complete word are consecutive after lexicographic sorting.
If there are at least two, choose the edge between its first two copies.
The localization places both ranks in \(H_{q_*}\), and this pair
is an edge of \(G\).  Distinct words occupy disjoint constant-word
blocks, so their chosen edges are distinct.
\end{proof}
\section{Collision counts}

For a full word $w$, write $N_w=C(w)$ for its occupancy and
$\mu_w=N\lambda_w$ for its mean.

\begin{lemma}[Fixed-column factorial moments]
\label{lem:fixed-column-factorial}
Suppose $K\le C\log N$ and every pile requested by the distinct full words
$w_1,\ldots,w_q$ has size at least $cN$, for fixed $c>0,C<\infty$.
For nonnegative integers $r_i$ with $R=\sum_i r_i\le N^{1/5}$,
\begin{equation}
\label{eq:exact-fixed-column-factorials}
 \E\prod_{i=1}^q(N_{w_i})_{r_i}
 =\prod_{i=1}^q\mu_{w_i}^{r_i}
    \left(1+O_{c,C}\left(\frac{(K+1)R^2}{N}\right)\right).
\end{equation}
The bound is uniform in the words and the column counts.
\end{lemma}

\begin{proof}
Expanding the falling factorials chooses ordered labelled rows for each
word. As the words are distinct, only disjoint choices contribute. There
are $(N)_R$ such choices. Let $m_{j,a}$ be the number of their requests for
digit $a$ in column $j$; thus $\sum_a m_{j,a}=R$.
Uniform sampling of each fixed-count column gives the exact formula
\[
 \E\prod_i(N_{w_i})_{r_i}
 =(N)_R\prod_{j=1}^K
                \frac{\prod_a(n_a^{(j)})_{m_{j,a}}}{(N)_R}.
\]
For $n\ge cN$ and $0\le v\le R\le N^{1/5}$,
\[
 \log\frac{(n)_v}{n^v}
 =\sum_{h=0}^{v-1}\log(1-h/n)=O_c(v^2/N),
\]
uniformly for sufficiently large $N$. Apply this also with $n=N$,
and compare the exact formula with
\[
 \prod_i\mu_{w_i}^{r_i}
   =N^R\prod_{j,a}(n_a^{(j)}/N)^{m_{j,a}}.
\]
Since $\sum_a m_{j,a}^2\le R^2$, the logarithm of their ratio is
$O_c((K+1)R^2/N)=o(1)$. Exponentiating proves
\eqref{eq:exact-fixed-column-factorials}. The case $R=0$ is equality.
\end{proof}

\begin{lemma}[Finite factorial comparison with a controlled remainder]
\label{lem:central-window}
Let $X,Y$ be nonnegative integer-valued random variables, let $\mu,\nu>0$,
and let $P_\mu,P_\nu$ be independent Poisson variables. Suppose that $m\ge2$
is even and, for $r,s\ge0$ with $r+s\le2m$,
\[
 \left|\E(X)_r(Y)_s-\mu^r\nu^s\right|
 \le\epsilon\mu^r\nu^s,\qquad 0\le\epsilon\le1.
\]
Set $f(t)=1-(1+t)e^{-t}$ and $d_m(t)=(m-1)t^m/m!$. Then
\begin{align*}
 \left|\Pp(X\ge2,Y\ge2)-f(\mu)f(\nu)\right|
 &\le \epsilon(1+\mu)(1+\nu)e^{\mu+\nu}\\
 &\quad+3\bigl(d_m(\mu)+d_m(\nu)+d_m(\mu)d_m(\nu)\bigr).
\end{align*}
If only the one-variable moments through order $m$ are assumed, then
\[
 \left|\Pp(X\ge2)-f(\mu)\right|
 \le\epsilon(1+\mu)e^\mu+3d_m(\mu).
\]
\end{lemma}

\begin{proof}
Define the even Bonferroni polynomial and its last term by
\[
 U_m(n)=\sum_{r=2}^{m}(-1)^r(r-1)\binom nr,
 \qquad D_m(n)=(m-1)\binom nm.
\]
The finite binomial identities give the pointwise inequalities
\begin{equation}
\label{eq:collision-bonferroni-sandwich}
 0\le U_m(n)-\1_{\{n\ge2\}}\le D_m(n)
 \qquad(n\ge0).
\end{equation}
For completeness, put $S_j(n)=\sum_{r=2}^{j}(-1)^r(r-1)\binom nr$.
For $n>j\ge2$,
\[
 S_j(n)-1=(-1)^j
 \left(n\binom{n-2}{j-1}-\binom{n-1}{j}\right),
\]
and the expression in parentheses equals
$((j-1)n+1)\binom{n-2}{j-1}/j\ge0$.
For $2\le n\le j$, $S_j(n)=1$; for $n=0,1$, it is zero.
Thus $S_{m-1}(n)\le\1_{\{n\ge2\}}\le S_m(n)$; the case $m=2$
uses $S_1=0$. Their difference is $D_m(n)$, proving
\eqref{eq:collision-bonferroni-sandwich}.

Writing $I_X=\1_{\{X\ge2\}}$ and similarly for $Y$, the same inequality gives
\[
 0\le U_m(X)U_m(Y)-I_XI_Y
 \le D_m(X)+D_m(Y)+D_m(X)D_m(Y).
\]
Consequently the expected error is at most
$(1+\epsilon)(d_m(\mu)+d_m(\nu)+d_m(\mu)d_m(\nu))$.
The independent Poisson pair has the same bound with $\epsilon=0$.
The difference between the expectations of the two polynomial products is
at most
\[
 \epsilon\left(\sum_{r=2}^{m}\frac{(r-1)\mu^r}{r!}\right)
 \left(\sum_{s=2}^{m}\frac{(s-1)\nu^s}{s!}\right)
 \le\epsilon(1+\mu)(1+\nu)e^{\mu+\nu}.
\]
Adding the three bounds proves the assertion. The one-variable assertion
follows from the same sandwich. All expectations here are of polynomials of
total degree at most $2m$: in particular their values on $\{X+Y>2m\}$ have
been controlled, rather than discarded using only the probability of that event.
\end{proof}

\begin{lemma}[Uniform fixed-column collision probabilities]
\label{lem:two-category-poisson}
Fix $c>0$ and $C<\infty$. For each $N$, let $K\le C\log N$, and take $K$
independent uniformly permuted fixed-count columns on the same $N$ labelled
rows. Let $\mathcal W$ be any deterministic family of distinct full words
such that every pile selected by any word in $\mathcal W$ has size at least
$cN$. Put
\[
 \lambda_w=\prod_{t=1}^{K}\frac{n_{w_t}^{(t)}}N,
 \quad \mu_w=N\lambda_w,
 \quad A_w=\{N_w\ge2\},\quad p_w=\Pp(A_w),
 \quad f(t)=1-(1+t)e^{-t}.
\]
Uniformly over all such data and words, as $N\to\infty$,
\begin{align}
 p_w&=f(\mu_w)\bigl(1+O_{c,C}(N^{-1/100})\bigr),
 \label{eq:uniform-collision-marginal}\\
 \Pp(A_w\cap A_{w'})
 &=p_wp_{w'}\bigl(1+O_{c,C}(N^{-1/100})\bigr),\qquad w\ne w'.
 \label{eq:uniform-collision-pair}
\end{align}
In particular $p_w\asymp\min\{\mu_w^2,1\}$. For
$Z=\sum_{w\in\mathcal W}\1_{A_w}$,
\begin{equation}
\label{eq:uniform-collision-variance}
 \Var Z\le\E Z+O_{c,C}(N^{-1/100})(\E Z)^2.
\end{equation}
No upper or lower restriction on $\mu_w$ is required.
\end{lemma}

\begin{proof}
All constants below depend only on $c,C$. The exact labelled-row calculation
in Lemma~\ref{lem:fixed-column-factorial} gives, for distinct words $w,w'$
and $r+s\le R\le N^{1/5}$,
\begin{equation}
\label{eq:collision-uniform-factorials}
 \E(N_w)_r(N_{w'})_s
 =\mu_w^r\mu_{w'}^s\left(1+O\left(\frac{(K+1)R^2}{N}\right)\right).
\end{equation}
The same one-variable estimate holds in every model obtained by forcing
$b\le3$ labelled rows to carry a word in $\mathcal W$ and deleting these
rows: the residual population is $N-b$, the columns remain independent
uniform fixed-count columns, and every selected residual pile has size at
least $(c/2)(N-b)$ for all sufficiently large $N$.

For the occupancy $V$ of a fixed word in any of these original or residual
models, let $v=\E V$. Equation~\eqref{eq:thinning-transform} gives
\begin{equation}
\label{eq:collision-empty-singleton-tail}
 \E s^V\le e^{-(1-s)v}\quad(0<s\le1),\qquad
 \Pp(V<2)\le2e^{-v/2}.
\end{equation}
The latter bound is also \eqref{eq:thinning-lower-tail}.

Choose once and for all
\[
 \tau=N^{-1/100},\qquad B=\frac1{10}\log N,
 \qquad R=\lfloor N^{1/25}\rfloor,
 \qquad m=2\lfloor R/4\rfloor.
\]
For large $N$, $m\ge2$ is even and $2m\le R\le N^{1/5}$.
The relative factorial error in \eqref{eq:collision-uniform-factorials}
is bounded by
\[
 \epsilon\le C_1(\log N)N^{-23/25}.
\]
If $0<\mu,\nu\le B$, Lemma~\ref{lem:central-window} therefore gives
\begin{equation}
\label{eq:collision-central-absolute}
 \left|\Pp(N_w\ge2,N_{w'}\ge2)-f(\mu)f(\nu)\right|
 \le C_2N^{-7/10}.
\end{equation}
Its one-variable version gives the same bound whenever $0<\mu\le2B$,
including the residual models on $N-b$ rows, $b\le3$.
To verify these errors explicitly, the moment-error term is at most
$C(\log N)^3N^{-18/25}=O(N^{-7/10})$ in the two-variable range,
and at most $C(\log N)^2N^{-18/25}$ in the one-variable range $\mu\le2B$.
Moreover,
\[
 d_m(t)\le m\left(\frac{et}{m}\right)^m
 \le\exp(-c_1N^{1/25}\log N),\qquad 0<t\le2B,
\]
for all sufficiently large $N$. The product of two such terms obeys the
same bound. Changing $N$ to $N-b$ changes only the constants.

We first prove \eqref{eq:uniform-collision-marginal}.
If $\mu\le\tau$, the order-two/order-three Bonferroni inequalities are
\begin{equation}
\label{eq:collision-rare-bonferroni}
 \frac12\E(V)_2-\frac13\E(V)_3
 \le\Pp(V\ge2)\le\frac12\E(V)_2.
\end{equation}
The coefficient of $(V)_3$ is $1/3$.
Equation~\eqref{eq:collision-uniform-factorials} thus gives
$p_w=\mu^2(1+O(\tau+(K+1)/N))/2$.
Since $f(\mu)=\mu^2(1+O(\mu))/2$, the desired relative estimate follows,
uniformly even when $\mu$ is arbitrarily small.
If $\tau<\mu\le B$, then $f(\mu)\ge c_2\min\{\mu^2,1\}\ge c_3\tau^2$.
The one-variable central bound has relative error
$O(N^{-7/10}\tau^{-2})=O(N^{-17/25})$, again smaller than $O(\tau)$.
If $\mu>B$, \eqref{eq:collision-empty-singleton-tail} gives
$1-p_w\le2N^{-1/20}$, whereas
$1-f(\mu)=(1+\mu)e^{-\mu}\le(1+B)e^{-B}=O((\log N)N^{-1/10})$.
This proves \eqref{eq:uniform-collision-marginal} in all three regimes.
The same conclusion holds uniformly in the residual models, by the verified
population and pile-size bounds above.

This marginal result also provides the conditional stability needed when
one of the two categories is rare. Fix $b=2$ or $3$ distinct labelled rows
and force them to carry $w$. For $w'\ne w$, none of the forced rows counts
as a $w'$ row. The conditional mean of the remaining $w'$ count is
\[
 \nu_b=(N-b)\prod_{t=1}^{K}
 \frac{n_{w'_t}^{(t)}-b\1_{\{w'_t=w_t\}}}{N-b},
 \qquad \frac{\nu_b}{\nu}=1+O((K+1)/N),\quad\nu=\mu_{w'}.
\]
All these conditioning events are positive for large $N$ because their
requested piles have size at least $cN$. The function $f$ satisfies
\[
 0\le\frac{tf'(t)}{f(t)}\le2\qquad(t>0),
\]
since $f(t)=\int_0^t u e^{-u}\,du\ge e^{-t}t^2/2$.
Integrating this bound with respect to $\log t$ shows
$f(\nu_b)/f(\nu)=1+O((K+1)/N)$, uniformly for every $\nu>0$.
Applying \eqref{eq:uniform-collision-marginal} to the original and residual
models proves
\begin{equation}
\label{eq:collision-conditional-stability}
 \Pp(A_{w'}\mid\text{the prescribed $b$ rows carry $w$})
 =p_{w'}(1+O(\tau)),\qquad b=2,3.
\end{equation}
The constants are independent of the row labels and both words.

To prove \eqref{eq:uniform-collision-pair}, first suppose
$\mu=\mu_w\le\tau$; no restriction is imposed on $\nu=\mu_{w'}$.
Summing \eqref{eq:collision-conditional-stability} over the ordered labelled
$b$-tuples gives
\[
 \E[(N_w)_b\1_{A_{w'}}]
 =\mu^b p_{w'}(1+O(\tau)),\qquad b=2,3.
\]
Multiplying the pointwise Bonferroni sandwich
\eqref{eq:collision-rare-bonferroni} by $\1_{A_{w'}}$ therefore yields
\[
 \Pp(A_w\cap A_{w'})=\frac{\mu^2}2p_{w'}(1+O(\tau))
 =p_wp_{w'}(1+O(\tau)).
\]
This includes the case where both means are arbitrarily small.

It remains to consider $\mu,\nu>\tau$. If, say, $\nu>B$, then
\eqref{eq:collision-empty-singleton-tail} implies
\[
 |\Pp(A_w\cap A_{w'})-p_wp_{w'}|
 \le2\Pp(A_{w'}^c)\le4N^{-1/20}.
\]
Since $p_w\ge c\tau^2$ and $p_{w'}\ge1/2$ for large $N$, its relative
error is $O(N^{-3/100})=O(\tau)$.
Finally, if $\tau<\mu,\nu\le B$, use
\eqref{eq:collision-central-absolute} and the already proved marginal
relative estimates. Here $f(\mu)f(\nu)\ge c\tau^4$, so the additional
central relative error is
$O(N^{-7/10}\tau^{-4})=O(N^{-33/50})$.
This proves the pair estimate uniformly in every case.

It follows that
\[
 \sum_{w\ne w'}|\operatorname{Cov}(\1_{A_w},\1_{A_{w'}})|
 \le C\tau\left(\sum_w p_w\right)^2.
\]
Adding $\sum_w\Var(\1_{A_w})\le\E Z$ proves
\eqref{eq:uniform-collision-variance}.
\end{proof}

\begin{lemma}[Many localized collision edges]
\label{lem:many-edges}
For the finite witness and cold set of Lemma~\ref{lem:cold-geometry},
there is a fixed $\delta>0$ such that
\[
 \Pp\{|E(G)\cap E(H_{L_{a_s}})|\ge M_N^{1/2+\delta}\}
 \longrightarrow1.
\]
The convergence is uniform over the permitted finite witnesses and cuts.
\end{lemma}

\begin{proof}
Let $A_y=\{N_y\ge2\}$ and $Z=\sum_{y\in\mathcal Y_N}\1_{A_y}$.
Every selected pile has size at least $cN$ for a fixed $c>0$, by
Lemma~\ref{lem:cold-geometry}. Hence
Lemma~\ref{lem:two-category-poisson} applies to this whole deterministic
family. Its marginal bound and \eqref{eq:cold-full-word-masses} give
\[
 \Pp(A_y)\ge c_1\min\{\mu_y^2,1\}
 \ge c_2N^{2z_N-4\kappa_0},\qquad
 \E Z\ge c_2N^{g_N-4\kappa_0},
\]
where the second inequality uses $z_N\le2\kappa_0$.
By \eqref{eq:cold-effective-collision} and $\kappa_0\le\kappa/100$,
\[
 g_N-4\kappa_0-h_N/2\ge\kappa/2-4\kappa_0>\kappa/3.
\]
Together with \eqref{eq:cold-geometric-conclusions} this yields, for all
sufficiently large $N$,
\[
 \E Z\ge M_N^{1/2}N^{\kappa/3},\qquad \E Z\longrightarrow\infty.
\]
Equation~\eqref{eq:uniform-collision-variance} and Chebyshev's inequality
now imply
\[
 \Pp(Z<\tfrac12\E Z)
 \le\frac4{\E Z}+O(N^{-1/100})=o(1).
\]
In particular $Z\ge M_N^{1/2}N^{\kappa/4}$ with probability $1-o(1)$.
Intersect this event with that of Lemma~\ref{lem:localization}. Its
injective assignment of edges gives
$|E(G)\cap E(H_{L_{a_s}})|\ge Z$. Since $M_N\le N$, any fixed
$0<\delta<\kappa/4$, for example $\delta=\kappa/5$, proves the claim.
\end{proof}

\section{Lalley ascent test}

For a permutation \(\sigma\in\mathfrak S_N\) and a set \(H\subseteq[N]\),
write
\[
  X_H(\sigma)=\sum_{(i,i+1)\in E(H)}
                 \1_{\{\sigma(i)<\sigma(i+1)\}}.
\]
If \(G\) is a subgraph of the path on \([N]\), list its nonempty
components from left to right as \(C_1,\ldots,C_q\).  The permutation
\(\pi_G\) is obtained by sorting the values of \(\pi\) increasingly
on each \(C_j\), and we put \(X_{G,H}(\pi)=X_H(\pi_G)\).
All probability and expectation statements with subscript \(\pi\) below
refer to a uniform permutation, with \(G\) and \(H\) fixed.

\begin{lemma}[Independent relative orders]
\label{lem:relative-order-independence}
Let \(A_1,\ldots,A_m\) be pairwise disjoint sets of positions in
\([N]\).  Under a uniform permutation \(\pi\), the relative orders of
the values on these sets are mutually independent and uniform.  In
particular, if \(A\) and \(B\) are disjoint nonempty sets of sizes
\(a\) and \(b\), then
\[
  \Pp_\pi\{\max\pi(A)<\min\pi(B)\}
  =\binom{a+b}{a}^{-1}.
\]
\end{lemma}

\begin{proof}
Put \(d_j=|A_j|\) and \(d_0=N-\sum_{j=1}^m d_j\).  Fix one
relative order on each \(A_j\), using increasing position order to
enumerate its entries.  There are
\(N!/(d_0!\prod_jd_j!)\) ways to assign disjoint value sets of the
prescribed sizes to these position sets and their complement.  Each
assignment admits exactly one ordering on each \(A_j\) with the
prescribed relative order and \(d_0!\) orderings on the complement.
Thus precisely \(N!/\prod_jd_j!\) permutations realize the prescribed
orders.  Dividing by \(N!\) proves the product uniform law.  On
\(A\cup B\), exactly \(a!b!\) of the \((a+b)!\) relative orders
put all values on \(A\) before those on \(B\), proving the last claim.
\end{proof}

\begin{lemma}[Component-sorting concentration]
\label{lem:permutation-bd}
Let \(G\) be a deterministic path subgraph, let \(H\subseteq[N]\),
and let \(M>0\) satisfy \(|H|\le M\).  For every \(t>0\),
\begin{equation}
\label{eq:component-ascent-one-sided}
  \Pp_\pi\{X_{G,H}-\E_\pi X_{G,H}\ge t\},\quad
  \Pp_\pi\{X_{G,H}-\E_\pi X_{G,H}\le-t\}
  \le \exp(-t^2/M).
\end{equation}
Consequently,
\[
  \Pp_\pi\{|X_{G,H}-\E_\pi X_{G,H}|\ge t\}
  \le 2\exp(-t^2/M).
\]
No condition on \(\partial H\) is needed for these estimates.
\end{lemma}

\begin{proof}
Every edge of \(E(G)\cap E(H)\) is an ascent after component
sorting.  Let \(J\subseteq\{1,\ldots,q-1\}\) consist of the
indices for which the edge joining \(C_j\) to \(C_{j+1}\) belongs
to \(E(H)\), and put \(B=|J|\).  Then
\[
  X_{G,H}=|E(G)\cap E(H)|+\sum_{j\in J}Y_j,
  \qquad
  Y_j=\1_{\{\max\pi(C_j)<\min\pi(C_{j+1})\}}.
\]
If \(B=0\), this variable is constant and both probabilities in
\eqref{eq:component-ascent-one-sided} vanish.  Assume \(B>0\).
Separate \(J\) into its even and odd component indices.  Within
either part, the sets \(C_j\cup C_{j+1}\) are disjoint, so
Lemma~\ref{lem:relative-order-independence} makes the corresponding
Bernoulli variables mutually independent.  Their means need not equal
\(1/2\).

For completeness, every Bernoulli variable \(Y\) satisfies
\[
  \E\exp\{\lambda(Y-\E Y)\}\le\exp(\lambda^2/8)
  \qquad(\lambda\in\mathbb R).
\]
Indeed the logarithm of the left side has value and first derivative
zero at the origin, and its second derivative is the variance of a
Bernoulli variable under exponential tilting, hence is at most \(1/4\).
Taylor's theorem gives the asserted bound, including the constant
Bernoulli cases.  Let \(S_0,S_1\) be the centered sums over the two
parity classes.  Cauchy--Schwarz and the independence within each class
give
\[
  \E_\pi e^{\lambda(S_0+S_1)}
  \le
  \bigl(\E_\pi e^{2\lambda S_0}
        \E_\pi e^{2\lambda S_1}\bigr)^{1/2}
  \le \exp(\lambda^2B/4).
\]
Markov's inequality, with \(\lambda=2t/B\), bounds either one-sided
tail by \(\exp(-t^2/B)\).  Finally,
\(B\le|E(H)|\le|H|\le M\), which proves
\eqref{eq:component-ascent-one-sided}.  Adding its two estimates proves
the two-sided bound.
\end{proof}

\begin{lemma}[Finite component expectation comparison]
\label{lem:component-ascent-expectation}
For every deterministic path subgraph \(G\) and set \(H\subseteq[N]\),
put \(e=|E(H)|\), \(r=|E(G)\cap E(H)|\), and
\(d=|\partial H|\), with the boundary taken in the integer path.
Then
\begin{equation}
\label{eq:component-ascent-expectation}
  \E_\pi X_{G,H}\ge \frac e2+\frac r6-d.
\end{equation}
\end{lemma}

\begin{proof}
Call a component \(C\) of \(G\) partial if it meets both \(H\)
and its complement in \([N]\).  Every partial component contains an
edge of \(\partial H\), and different components contain different
such edges.  Thus there are at most \(d\) partial components.
Call an edge of \(E(H)\setminus E(G)\) bad if at least one of its
two incident components is partial.  A path component is incident to
at most two intercomponent edges, so there are at most \(2d\) bad
edges.  Each has ascent probability at least zero and therefore loses
at most \(1/2\) against the baseline \(1/2\).

The two components incident to any remaining intercomponent edge are
wholly contained in \(H\).  If their sizes are \(a,b\ge1\), its
ascent probability is \(\binom{a+b}{a}^{-1}\), by
Lemma~\ref{lem:relative-order-independence}, and
\begin{equation}
\label{eq:component-boundary-deficit}
  \frac12-\binom{a+b}{a}^{-1}
  \le \frac{(a-1)+(b-1)}6.
\end{equation}
For \(a+b\ge5\), the right side is at least \(1/2\).
For \(a+b=2\) both sides are zero; for \(a+b=3\) both sides are
\(1/6\); and for \(a+b=4\) the left side is either \(1/4\)
or \(1/3\), while the right side is \(1/3\).
Each component wholly contained in \(H\) is charged at most twice
in \eqref{eq:component-boundary-deficit}.  Its \(|C|-1\) internal
edges belong to \(E(G)\cap E(H)\).  The total deficit from these
remaining intercomponent edges is therefore at most \(r/3\).
All \(r\) internal graph edges, including those in partial components,
are retained and contribute excess \(r/2\).  Combining the internal
excess, the deficit at most \(r/3\), and the bad-edge deficit at most
\(d\), gives \eqref{eq:component-ascent-expectation}.
\end{proof}

\begin{lemma}[Uniform translate-scan ascent bound]
\label{lem:uniform-ascent}
Let \(H_0\subset\mathbb Z\) be finite with \(M=|H_0|>0\),
and put \(H_q=(q+H_0)\cap[N]\) for \(0\le q\le N\).
For every \(s>0\), the event
\[
  \mathcal A(s)=
  \{\sigma:\exists q\in\{0,\ldots,N\},\
          X_{H_q}(\sigma)>|E(H_q)|/2+s\}
\]
satisfies
\begin{equation}
\label{eq:uniform-ascent-scan-finite}
  \nuUnif(\mathcal A(s))\le (N+1)\exp(-s^2/M).
\end{equation}
In particular, if \(M\ge N^a\) for a fixed \(a>0\), then for
every fixed \(\delta>0\),
\[
  \nuUnif\bigl(\mathcal A(M^{1/2+\delta/2})\bigr)\longrightarrow0.
\]
\end{lemma}

\begin{proof}
A comparison between two distinct positions of a uniform permutation
has mean \(1/2\), so \(\E_\pi X_{H_q}=|E(H_q)|/2\).
Apply Lemma~\ref{lem:permutation-bd} to the edgeless graph and
\(H_q\), for which \(|H_q|\le M\), and take a union bound over
the \(N+1\) choices of \(q\).  This proves
\eqref{eq:uniform-ascent-scan-finite}.  With
\(s=M^{1/2+\delta/2}\), its right side is at most
\((N+1)\exp(-N^{a\delta})\), which tends to zero.
\end{proof}

\begin{lemma}[Lalley component comparison]
\label{lem:lalley-comparison}
Let \(M>0\), \(|H|\le M\), and \(G\) be a deterministic
path subgraph.  Write \(e=|E(H)|\),
\(r=|E(G)\cap E(H)|\), and \(d=|\partial H|\).
If \(s>0\) and \(\Delta=r/6-d-s>0\), then
\begin{equation}
\label{eq:lalley-comparison-finite}
  \Pp_\pi\{X_{G,H}\le e/2+s\}
  \le\exp(-\Delta^2/M).
\end{equation}
Consequently, for fixed \(D\ge0\), \(\delta>0\), and any such
family with \(M\to\infty\), \(d\le D\sqrt M\), and
\(r\ge M^{1/2+\delta}\),
\[
  \Pp_\pi\{X_{G,H}>e/2+M^{1/2+\delta/2}\}\longrightarrow1
\]
uniformly over \(N,G,H\).
\end{lemma}

\begin{proof}
Lemma~\ref{lem:component-ascent-expectation} gives
\(\E_\pi X_{G,H}-(e/2+s)\ge\Delta\).
The event in \eqref{eq:lalley-comparison-finite} is therefore
contained in the lower-tail event from
\eqref{eq:component-ascent-one-sided}, with \(t=\Delta\).
For the last assertion take \(s=M^{1/2+\delta/2}\).
For all sufficiently large \(M\), depending only on \(D,\delta\),
\[
  D+M^{\delta/2}\le M^\delta/12,
  \qquad
  \Delta\ge M^{1/2+\delta}/12.
\]
Thus the failure probability is at most
\(\exp(-M^{2\delta}/144)\), uniformly as claimed.
\end{proof}

\begin{lemma}[Lalley cold-spot test]
\label{lem:lalley}
Let \(H_0\subset\mathbb Z\) be deterministic and finite, put
\(M=|H_0|\), and let \(H_q=(q+H_0)\cap[N]\).
Suppose that, for constants \(a,\delta>0\) and \(D<\infty\)
independent of \(N\),
\[
  M\ge N^a,\qquad |\partial H_q|\le D\sqrt M
  \quad(0\le q\le N),
\]
and that the shuffle graph satisfies
\[
  \Pp\{\exists q\in\{0,\ldots,N\}:\
         |E(G)\cap E(H_q)|\ge M^{1/2+\delta}\}\longrightarrow1.
\]
Then \(d_{\TV}(P_N^{K_T},\nuUnif)\longrightarrow1\).
\end{lemma}

\begin{proof}
Use the observable event
\(\mathcal A=\mathcal A(M^{1/2+\delta/2})\) from
Lemma~\ref{lem:uniform-ascent}.  That lemma gives
\(\nuUnif(\mathcal A)\to0\).
Use the representation in Lemma~\ref{lem:ub-graph-representation}
to sample the shuffle graph and an independent uniform permutation
\(\pi\), and then form \(\pi_G\).
On the event that a translate has at least
\(M^{1/2+\delta}\) internal graph edges, choose its smallest index
\(q\).  This choice depends only on \(G\).
Conditional on \(G\), Lemma~\ref{lem:lalley-comparison} bounds the
probability that this translate fails the ascent test by
\(\exp(-M^{2\delta}/144)\), for all sufficiently large \(N\).
The bound is uniform in \(G\) and the chosen translate.  Hence
\[
  P_N^{K_T}(\mathcal A)
  \ge
  \Pp\{\exists q:\ |E(G)\cap E(H_q)|\ge M^{1/2+\delta}\}
     -\exp(-M^{2\delta}/144)
  \longrightarrow1.
\]
Finally,
\(d_{\TV}(P_N^{K_T},\nuUnif)
\ge P_N^{K_T}(\mathcal A)-\nuUnif(\mathcal A)\to1\).
\end{proof}
\section{Completion of the lower bound}
\label{sec:completion-matching-lower-bound}

\begin{proposition}[Finite cold-spot lower bound]
\label{prop:cold-lower}
Fix $\eta_f>0$, set $\kappa=\eta_f/1000$ and
$\kappa_0=\min\{\kappa/100,1/1000\}$, and choose the good-column mesh
so that $2\rho\log k\le\kappa_0/100$. If
\[
 \Lambda_{1,N}\ge\eta_f,\qquad \beta_N(S_N)<\kappa_0,
\]
then $d_{\TV}(P_N^{K_T},\nuUnif)\to1$, uniformly over these finite
cut sequences satisfying the local good-column condition.
\end{proposition}

\begin{proof}
Choose the witness from Lemma~\ref{lem:cold-witness} with $r_b=2\rho$,
and the deterministic cold set from Lemma~\ref{lem:cold-geometry}.
Its size satisfies $M_N\ge N^{\kappa/3}$ and its clipped translates
have boundary $O(\sqrt{M_N})$, by \eqref{eq:cold-geometric-conclusions}.
Lemma~\ref{lem:many-edges} supplies the required excess of graph edges
on the translate indexed by $L_{a_s}\in\{0,\ldots,N\}$.
All hypotheses of Lemma~\ref{lem:lalley} therefore hold, proving the claim.
\end{proof}

\begin{proof}[Proof of Theorem~\ref{thm:buffered-reduced-lower-bound}]
Put $\eta_f=\eta/2$, $\kappa=\eta_f/1000$, and
$\kappa_0=\min\{\kappa/100,1/1000\}$. Choose
$0<e<\min\{\kappa_0/10,\eta/100\}$ in
Lemma~\ref{lem:lower-profile-transfer}, requiring in addition
$2\rho\log k\le\kappa_0/100$.
Fix the resulting partition and mesh, the frequency tolerance, and the
positive local good-column constants before letting $N$ tend to infinity.

If $\beta_N(S_N)\ge\kappa_0$, the maximal full word has the exact mean
\[
 \mu_*=N\prod_{j=1}^{K_T}\max_\ell p^{(j)}_\ell
       =N^{\beta_N(S_N)}\ge N^{\kappa_0}.
\]
Proposition~\ref{prop:largest-prefix} gives the asserted lower bound,
uniformly on this case. Otherwise $\beta_N(S_N)<\kappa_0$, and
\eqref{eq:lower-margin-transfer} implies
$\Lambda_0(T)<\kappa_0+e+O_k(1/\log N)<\eta$ for large $N$.
The assumed maximum therefore forces $\Lambda_1(T)\ge\eta$.
The same transfer estimate gives
$\Lambda_{1,N}\ge\eta/2=\eta_f$, so
Proposition~\ref{prop:cold-lower} applies.
Both cases have uniform errors tending to zero. This proves the assertion
even when the finite cut sequence alternates between the two cases as $N$
varies.
\end{proof}

\begin{thebibliography}{1}

\bibitem{Hoeffding1963}
Wassily Hoeffding.
\newblock Probability inequalities for sums of bounded random variables.
\newblock \emph{Journal of the American Statistical Association},
  58(301):13--30, 1963.
\newblock
  \href{https://doi.org/10.1080/01621459.1963.10500830}
  {doi:10.1080/01621459.1963.10500830}.

\end{thebibliography}

\end{document}
